% La vie culturelle 2 (la musique et la danse) — Allemand Tle A — Cours signé OnBuch+
% Programme officiel MINESEC. Compiler avec : tectonic AL23-la-vie-culturelle-2-la-musique-et-la-danse.tex
\newcommand{\DOCMATIERE}{Allemand}
\newcommand{\DOCNIVEAU}{Tle A}
\newcommand{\DOCMODULE}{Chapitre 5 : Citoyenneté}
\newcommand{\DOCLECON}{AL23}
\newcommand{\DOCTITRE}{La vie culturelle 2 (la musique et la danse)}
% OnBuch+ — préambule « cours signés » ENRICHI (compilé avec tectonic / XeLaTeX)
% Système visuel « Pop » d'OnBuch+ : fond crème (jamais blanc), bordures encre
% épaisses, ombres « dures » décalées, titres Archivo Black en capitales.
% Version MATHÉMATIQUES : graphes (pgfplots/tikz), boîtes pédagogiques
% (définition, propriété, méthode, attention, le savais-tu, exercice résolu,
% exercice, TP, bilan), page de garde et signature OnBuch+ ancrée en bas.
%
% Macros attendues dans main.tex AVANT \input{preamble} :
%   \DOCMATIERE  \DOCNIVEAU  \DOCMODULE  \DOCLECON (numéro)  \DOCTITRE
\documentclass[11pt]{article}

\usepackage{fontspec}
\usepackage[ngerman,french]{babel} % français = langue principale ; l'allemand via \all{..} / textebox
\usepackage[a4paper,margin=2cm,top=3.4cm,bottom=2.8cm,headheight=1.4cm,footskip=1.3cm]{geometry}
\usepackage{amsmath,amssymb,amsfonts}
\usepackage{mathtools} % \xmapsto, \coloneqq... souvent utilisés par les modèles
\usepackage{cancel} % simplifications barrées (bilans de forces)
% \abs{z} (module, valeur absolue) et \norm{u} (norme), utilisés par les modèles
\providecommand{\abs}{}\let\abs\relax\DeclarePairedDelimiter{\abs}{\lvert}{\rvert}
\providecommand{\norm}{}\let\norm\relax\DeclarePairedDelimiter{\norm}{\lVert}{\rVert}
\usepackage[table]{xcolor} % [table] : fournit aussi \rowcolors (alternance de lignes)
\usepackage{pagecolor}
\usepackage{tikz}
\usetikzlibrary{arrows.meta,positioning,calc,shapes.geometric,shapes.misc,shapes.arrows,decorations.pathmorphing,decorations.pathreplacing,decorations.markings,patterns,shadows,mindmap,trees,fit,backgrounds,matrix,intersections,angles,quotes}
\usepackage{pgfplots}
\usepackage[version=4]{mhchem} % équations nucléaires \ce{^{235}_{92}U}
\usepackage[european,siunitx]{circuitikz} % schémas électriques
\pgfplotsset{compat=1.18}
\usepgfplotslibrary{fillbetween,groupplots} % aires entre courbes (calcul intégral)
\usepackage{enumitem}
\usepackage{fancyhdr}
% [most] charge toutes les bibliothèques tcolorbox (listings, minted, raster,
% documentation...) et gonfle inutilement la mémoire TeX sur un document long
% (~300 boîtes) : on ne charge que ce qui est réellement utilisé.
\usepackage[skins,breakable]{tcolorbox}
\usepackage{graphicx}
\usepackage{colortbl}
\usepackage{booktabs}
\usepackage{tabularx}
\usepackage{diagbox} % case d'en-tête diagonale des tableaux à double entrée
% Colonnes extensibles centrée / gauche / droite (C, L, R), souvent utilisées par les modèles
\newcolumntype{C}{>{\centering\arraybackslash}X}
\newcolumntype{L}{>{\raggedright\arraybackslash}X}
\newcolumntype{R}{>{\raggedleft\arraybackslash}X}
\usepackage{array}
\usepackage{multicol}
\usepackage{multirow}
\usepackage{siunitx}
% parse-numbers=false : accepte aussi \SI{2,5 \times 10^{-3}}{...}
\sisetup{locale=FR,output-decimal-marker={,},group-minimum-digits=4,parse-numbers=false}
\DeclareSIUnit\ohm{\ensuremath{\symup{\Omega}}} % U+2126 absent de la police
\DeclareSIUnit\division{div} % réglages d'oscilloscope (ms/div, V/div)
\DeclareSIUnit\tr{tr} % tours (vitesse de rotation, ex. \SI{1500}{\tr\per\minute})
\DeclareSIUnit\molar{M} % concentration molaire (ex. \SI{3}{\molar}, \SI{1}{\micro\molar})
\DeclareSIUnit\an{an} % année (le modèle confond parfois avec \year, un compteur TeX) — ex. \SI{5}{\kilo\meter\per\an}
\DeclareSIUnit\FCFA{FCFA} % franc CFA (ex. \SI{500}{\FCFA\per\kilo\gram})
\DeclareSIUnit\degreeBrix{\textdegree Brix} % teneur en sucre d'un sirop/jus (réfractométrie)
\DeclareSIUnit\month{mois} % le modèle utilise parfois \month, un compteur TeX, comme unité de durée
\usepackage{qrcode}
% \degree : commande fréquemment utilisée par les modèles pour un angle ou une
% température en dehors de \SI{}{} (non fournie par les paquets chargés).
\providecommand{\degree}{^{\circ}}
% \qed (fin de démonstration), utilisé par les modèles sans amsthm
\providecommand{\qed}{\hfill$\square$}
% \barre : double barre des valeurs interdites dans un tableau de variations (tkz-tab)
\providecommand{\barre}{\ensuremath{\|}}
% environnement proof (amsthm non chargé) utilisé par les modèles
\ifdefined\proof\else\newenvironment{proof}[1][Démonstration]{\par\noindent\textit{#1.}\ }{\qed\par}\fi
\usepackage{lastpage}
\usepackage{hyperref}
\hypersetup{hidelinks}

% ============================================================
% Palette « Pop » OnBuch+ — mêmes hex que la classe P (plus_ui.dart)
% ============================================================
\definecolor{poporange}{HTML}{F59321}
\definecolor{popdark}{HTML}{E07A0C}
\definecolor{popcream}{HTML}{FFF6E8}
\definecolor{popink}{HTML}{1C1714}
\definecolor{poppurple}{HTML}{7A5AE0}
\definecolor{popgreen}{HTML}{2E9E6A}
\definecolor{poppink}{HTML}{E0457B}
\definecolor{popblue}{HTML}{2F7DE1}
\definecolor{popgold}{HTML}{C98A12}
\definecolor{popmuted}{HTML}{8A7F76}
\definecolor{popline}{HTML}{EADFD2}
\definecolor{popgray}{HTML}{8A7F76}
\colorlet{popgrey}{popgray}
% Alias tolérants (le modèle nomme parfois les couleurs différemment)
\colorlet{popwhite}{white}
\colorlet{popblack}{popink}
\colorlet{popred}{poppink}
\colorlet{popyellow}{popgold}
% Teintes claires (fonds de tableaux/encadrés) — le modèle nomme parfois
% les couleurs avec un suffixe L (ex. popblueL) sans les avoir définies.
\colorlet{poporangeL}{poporange!20}
\colorlet{popdarkL}{popdark!20}
\colorlet{poppurpleL}{poppurple!20}
\colorlet{popgreenL}{popgreen!20}
\colorlet{poppinkL}{poppink!20}
\colorlet{popblueL}{popblue!20}
\colorlet{popgoldL}{popgold!20}
\colorlet{popredL}{popred!20}
\colorlet{popgrayL}{popgray!20}
% Teintes claires pour fonds de tableaux / zones
\colorlet{popblueL}{popblue!10!white}
\colorlet{poporangeL}{poporange!14!white}
\colorlet{popgreenL}{popgreen!12!white}
\colorlet{poppurpleL}{poppurple!12!white}
\colorlet{poppinkL}{poppink!12!white}
\colorlet{popgoldL}{popgold!14!white}

\pagecolor{popcream}

% ============================================================
% Typographie
% ============================================================
\defaultfontfeatures{Ligatures=TeX}
\setmainfont{PlusJakartaSans-Medium.ttf}[Path=../../../fonts/,BoldFont=PlusJakartaSans-Bold.ttf]
% Maths en sans-serif (Fira Math) avec chiffres et lettres droites de Plus
% Jakarta, pour que les formules se fondent dans le texte.
\usepackage[math-style=ISO,bold-style=ISO]{unicode-math}
\setmathfont{FiraMath-Regular.otf}[Path=../../../fonts/]
\setmathfont{PlusJakartaSans-Medium.ttf}[Path=../../../fonts/,range={up/num,up/{Latin,latin},\mathup}]
\usepackage{icomma} % virgule décimale sans espace en mode math
% Caractères mathématiques Unicode absents des polices de texte (Plus Jakarta,
% Archivo Black) : rendus par leur équivalent mathématique, en texte comme en
% formule — sinon ils s'affichent en carré vide (ex. « Arithmétique dans ℤ »).
\usepackage{newunicodechar}
\newunicodechar{ℝ}{\ensuremath{\mathbb{R}}}
\newunicodechar{ℤ}{\ensuremath{\mathbb{Z}}}
\newunicodechar{ℂ}{\ensuremath{\mathbb{C}}}
\newunicodechar{ℕ}{\ensuremath{\mathbb{N}}}
\newunicodechar{ℚ}{\ensuremath{\mathbb{Q}}}
\newunicodechar{𝔻}{\ensuremath{\mathbb{D}}}
\newunicodechar{α}{\ensuremath{\alpha}}
\newunicodechar{β}{\ensuremath{\beta}}
\newunicodechar{γ}{\ensuremath{\gamma}}
\newunicodechar{δ}{\ensuremath{\delta}}
\newunicodechar{ε}{\ensuremath{\varepsilon}}
\newunicodechar{θ}{\ensuremath{\theta}}
\newunicodechar{λ}{\ensuremath{\lambda}}
\newunicodechar{μ}{\ensuremath{\mu}}
\newunicodechar{σ}{\ensuremath{\sigma}}
\newunicodechar{τ}{\ensuremath{\tau}}
\newunicodechar{φ}{\ensuremath{\varphi}}
\newunicodechar{ψ}{\ensuremath{\psi}}
\newunicodechar{ω}{\ensuremath{\omega}}
\newunicodechar{Σ}{\ensuremath{\Sigma}}
% majuscules produites par \MakeUppercase dans les titres (α-aminés -> Α)
\newunicodechar{Α}{\ensuremath{\alpha}}
\newunicodechar{Β}{\ensuremath{\beta}}
\newunicodechar{Γ}{\ensuremath{\Gamma}}
\newunicodechar{Δ}{\ensuremath{\Delta}}
\newunicodechar{Ε}{\ensuremath{\varepsilon}}
\newunicodechar{Θ}{\ensuremath{\Theta}}
\newunicodechar{Λ}{\ensuremath{\Lambda}}
\newunicodechar{Μ}{\ensuremath{\mu}}
\newunicodechar{Τ}{\ensuremath{\tau}}
\newunicodechar{Φ}{\ensuremath{\Phi}}
\newunicodechar{Ψ}{\ensuremath{\Psi}}
\newunicodechar{Ω}{\ensuremath{\Omega}}
\newunicodechar{↦}{\ensuremath{\mapsto}}
\newunicodechar{∈}{\ensuremath{\in}}
\newunicodechar{∉}{\ensuremath{\notin}}
\newunicodechar{∧}{\ensuremath{\wedge}}
\newunicodechar{⇔}{\ensuremath{\Leftrightarrow}}
\newunicodechar{⟺}{\ensuremath{\Longleftrightarrow}}
\newunicodechar{⇒}{\ensuremath{\Rightarrow}}
\newunicodechar{⟹}{\ensuremath{\Longrightarrow}}
\newunicodechar{∘}{\ensuremath{\circ}}
\newunicodechar{∪}{\ensuremath{\cup}}
\newunicodechar{∩}{\ensuremath{\cap}}
\newunicodechar{≡}{\ensuremath{\equiv}}
\newunicodechar{⊂}{\ensuremath{\subset}}
\newunicodechar{⊥}{\ensuremath{\perp}}
\newunicodechar{∃}{\ensuremath{\exists}}
\newunicodechar{∀}{\ensuremath{\forall}}
\newunicodechar{∛}{\ensuremath{\sqrt[3]{\,}}}
\newunicodechar{∜}{\ensuremath{\sqrt[4]{\,}}}
\newunicodechar{⃗}{\ensuremath{{}^{\to}}}
\newunicodechar{ⁿ}{\ifmmode^{n}\else\textsuperscript{\ensuremath{n}}\fi}
\newunicodechar{ˣ}{\ifmmode^{x}\else\textsuperscript{\ensuremath{x}}\fi}
\newunicodechar{ᵇ}{\ifmmode^{b}\else\textsuperscript{\ensuremath{b}}\fi}
\newunicodechar{ʳ}{\ifmmode^{r}\else\textsuperscript{\ensuremath{r}}\fi}
\newunicodechar{ᵘ}{\ifmmode^{u}\else\textsuperscript{\ensuremath{u}}\fi}
\newunicodechar{ʸ}{\ifmmode^{y}\else\textsuperscript{\ensuremath{y}}\fi}
\newunicodechar{ᵃ}{\ifmmode^{a}\else\textsuperscript{\ensuremath{a}}\fi}
\newunicodechar{ᵉ}{\ifmmode^{e}\else\textsuperscript{\ensuremath{e}}\fi}
\newunicodechar{ᵏ}{\ifmmode^{k}\else\textsuperscript{\ensuremath{k}}\fi}
\newunicodechar{ᵖ}{\ifmmode^{p}\else\textsuperscript{\ensuremath{p}}\fi}
\newunicodechar{ᴺ}{\ifmmode^{N}\else\textsuperscript{\ensuremath{N}}\fi}
\newunicodechar{ᶜ}{\ifmmode^{c}\else\textsuperscript{\ensuremath{c}}\fi}
\newunicodechar{ʰ}{\ifmmode^{h}\else\textsuperscript{\ensuremath{h}}\fi}
\newunicodechar{ᵠ}{\ifmmode^{q}\else\textsuperscript{\ensuremath{q}}\fi}
\newunicodechar{ₙ}{\ifmmode_{n}\else\textsubscript{\ensuremath{n}}\fi}
\newunicodechar{ₐ}{\ifmmode_{a}\else\textsubscript{\ensuremath{a}}\fi}
\newunicodechar{ᵢ}{\ifmmode_{i}\else\textsubscript{\ensuremath{i}}\fi}
\newunicodechar{ᵣ}{\ifmmode_{r}\else\textsubscript{\ensuremath{r}}\fi}
\newunicodechar{ₖ}{\ifmmode_{k}\else\textsubscript{\ensuremath{k}}\fi}
\newunicodechar{ᵦ}{\ifmmode_{\beta}\else\textsubscript{\ensuremath{\beta}}\fi}
\newunicodechar{ₓ}{\ifmmode_{x}\else\textsubscript{\ensuremath{x}}\fi}
\newfontfamily\popdisplay{ArchivoBlack-Regular.ttf}[Path=../../../fonts/]
\newfontfamily\poplogo{SpaceGrotesk-Bold.ttf}[Path=../../../fonts/]
\newfontfamily\popsemi{PlusJakartaSans-SemiBold.ttf}[Path=../../../fonts/]
\newfontfamily\popxbold{PlusJakartaSans-ExtraBold.ttf}[Path=../../../fonts/]
\newfontfamily\popxboldls{PlusJakartaSans-ExtraBold.ttf}[Path=../../../fonts/,LetterSpace=3.0]

\setlist[enumerate]{leftmargin=1.7em,itemsep=5pt,topsep=4pt}
\setlist[itemize]{leftmargin=1.7em,itemsep=4pt,topsep=4pt,label={\color{poporange}\textbullet}}
\setlength{\parindent}{0pt}
\setlength{\parskip}{5pt}
\renewcommand{\arraystretch}{1.25}

% pgfplots : style Pop par défaut
\pgfplotsset{
  every axis/.append style={axis line style={popink,line width=1pt},tick style={popink},
    grid style={popline,line width=0.5pt},label style={font=\small},tick label style={font=\footnotesize}},
  popcurve/.style={poporange,line width=1.8pt,no markers,smooth},
}
% Même style disponible dans un \draw TikZ simple (hors pgfplots)
\tikzset{popcurve/.style={poporange,line width=1.8pt}}
\tikzset{popgrid/.style={popline,very thin}}
\colorlet{popcurve}{poporange} % le nom du style est parfois utilisé comme couleur

% ============================================================
% Pill / squiggle
% ============================================================
\newcommand{\poppill}[3]{%
  \tikz[baseline=-0.55ex]\node[draw=popink,line width=1.2pt,rounded corners=2.6pt,%
    fill=#2,inner xsep=6.5pt,inner ysep=3pt]{%
    \popxboldls\fontsize{7}{7}\selectfont\color{#3}\MakeUppercase{#1}};%
}
\newcommand{\popsquiggle}[1]{%
  \tikz[baseline=-0.4ex]{
    \draw[#1,line width=2.6pt,line cap=round]
      plot [smooth] coordinates {(0,0.09) (0.34,0.01) (0.68,0.21) (1.02,0.01) (1.36,0.10)};
  }%
}
% Mot-clé surligné (marqueur orange sous le texte)
\newcommand{\cle}[1]{{\popxbold\color{popdark}#1}}

% ============================================================
% En-tête / pied
% ============================================================
\pagestyle{fancy}
\fancyhf{}
\renewcommand{\headrulewidth}{0pt}
\renewcommand{\footrulewidth}{0pt}
\lhead{\raisebox{-0.15ex}{\poplogo\fontsize{13}{13}\selectfont\color{popink}OnBuch\color{poporange}+}}
\rhead{\poppill{\DOCMATIERE\ \textbullet\ \DOCNIVEAU\ \textbullet\ Leçon \DOCLECON}{white}{popink}}
\lfoot{\popxbold\fontsize{7.6}{7.6}\selectfont\color{popmuted}\MakeUppercase{Cours signé OnBuch+}\ \textbullet\ {\popsemi\DOCTITRE}}
\rfoot{\tikz[baseline=-0.6ex]\node[rounded corners=8.5pt,fill=popink,minimum height=17pt,minimum width=17pt,inner xsep=5pt,inner sep=0]{\popxbold\fontsize{8}{8}\selectfont\color{popcream}\thepage\,/\,\pageref*{LastPage}};}

% ============================================================
% Titres
% ============================================================
\newcommand{\chaptitre}[2]{%
  \begin{tcolorbox}[enhanced,colback=poporange,colframe=popink,boxrule=2.6pt,arc=11pt,
      drop shadow=popink,left=15pt,right=15pt,top=12pt,bottom=11pt,before skip=3pt,after skip=18pt]
    \poppill{#2}{popink}{popcream}\par\vspace{9pt}
    {\raggedright\hyphenpenalty=10000\exhyphenpenalty=10000\popdisplay\fontsize{22}{24}\selectfont\color{popcream}\MakeUppercase{#1}\par}
  \end{tcolorbox}
}
% Section numérotée : pastille orange avec le numéro + titre Archivo
\newcounter{popsec}
\newcommand{\coursec}[1]{%
  \stepcounter{popsec}\setcounter{popsub}{0}%
  \par\vspace{16pt}\noindent
  \begin{minipage}{\linewidth}
  \tikz[baseline=-0.7ex]\node[draw=popink,line width=1.6pt,fill=poporange,rounded corners=4pt,minimum size=22pt,inner sep=0]{\popdisplay\fontsize{12}{12}\selectfont\color{popcream}\thepopsec};%
  \hspace{8pt}{\popdisplay\fontsize{15}{17}\selectfont\color{popink}\MakeUppercase{#1}}\par
  \vspace{4pt}\noindent{\color{poporange}\rule{36pt}{3.6pt}}
  \end{minipage}\par\nopagebreak\vspace{8pt}
}
\newcounter{popsub}
\newcommand{\courssub}[1]{%
  \stepcounter{popsub}%
  \par\vspace{9pt}\noindent{\popxbold\fontsize{11.5}{13}\selectfont\color{popink}{\color{poporange}\thepopsec.\thepopsub}\hspace{6pt}#1}\par\nopagebreak\vspace{4pt}
}

% ============================================================
% Boîtes pédagogiques — même recette : fond blanc, bordure épaisse colorée,
% ombre dure encre, badge plein. Titre optionnel [#1].
% ============================================================
\tcbset{popbase/.style={breakable,enhanced,colback=white,boxrule=2.3pt,arc=9pt,
  shadow={3pt}{-3pt}{0pt}{popink},left=14pt,right=14pt,top=11pt,bottom=11pt,before skip=11pt,after skip=15pt,
  fontupper=\popsemi\fontsize{10.6}{14.5}\selectfont\color{popink},
  fontlower=\popsemi\fontsize{10.6}{14.5}\selectfont\color{popink}}}
% Coche dessinée (le glyphe ✓ n'existe pas dans les polices du document)
\newcommand{\popcheck}[1]{\tikz[baseline=-0.5ex]\draw[#1,line width=1.6pt,line cap=round,line join=round](-0.13,0.0)--(-0.04,-0.09)--(0.13,0.1);}
\providecommand{\coche}{\popcheck{popgreen}}
\providecommand{\resultat}[1]{\par\smallskip\noindent\fcolorbox{popgreen}{popgreenL}{\parbox{\dimexpr\linewidth-2\fboxsep-2\fboxrule\relax}{\textbf{Résultat.} #1}}\par\smallskip}
\newcommand{\popboxhead}[3]{\poppill{#1}{#2}{white}\ifx\relax#3\relax\else\hspace{6pt}{\popxbold\fontsize{10.6}{12}\selectfont\color{popink}#3}\fi\par\vspace{7pt}}

\newtcolorbox{aretenir}[1][]{popbase,colframe=popgreen,before upper={\popboxhead{À retenir}{popgreen}{#1}}}
% Texte en allemand (document de lecture, dialogue, extrait) : cadre
% d'étude. Un titre optionnel (source, auteur) s'affiche dans l'en-tête.
\newtcolorbox{textebox}[1][]{popbase,colframe=popblue,before upper={\popboxhead{Text}{popblue}{#1}\begin{otherlanguage}{ngerman}},after upper={\end{otherlanguage}}}
% Marques ✓ / ✗ (forme correcte / fautive) souvent utilisées par les modèles
\providecommand{\cross}{\textcolor{poppink}{\ensuremath{\times}}}
\providecommand{\xmark}{\cross}\providecommand{\crossmark}{\cross}\providecommand{\cmark}{\ensuremath{\checkmark}}
% Ligne de réponse / justification à compléter (inventée par les modèles)
\providecommand{\justification}[1]{\par\noindent\textit{Justification~:} \dotfill\par}
% \textipa (package tipa non chargé) : la police ne couvre pas l'API -> simple passage
\providecommand{\textipa}[1]{#1}
% Soulignement (ulem non chargé) : \uline, \ul
\providecommand{\uline}[1]{\underline{#1}}\providecommand{\ul}[1]{\underline{#1}}
% \glossaire{\item ...} inventée par les modèles : liste de glossaire
\providecommand{\glossaire}[1]{\begin{itemize}#1\end{itemize}}
% \alert (beamer) inventée par les modèles : texte mis en évidence
\providecommand{\alert}[1]{\textbf{\textcolor{poppink}{#1}}}
% variantes ASCII d'ordinaux écrites en macro
\providecommand{\iere}{\textsuperscript{re}}\providecommand{\ieme}{\textsuperscript{e}}\providecommand{\ere}{\textsuperscript{re}}\providecommand{\eme}{\textsuperscript{e}}\providecommand{\er}{\textsuperscript{er}}\providecommand{\ie}{\textsuperscript{e}}
% Ordinaux français écrits en macro par les modèles : 1\ière, 3\ième
\newcommand{\ière}{\textsuperscript{re}}\newcommand{\ième}{\textsuperscript{e}}
% \exemple (inventée par les modèles) : étiquette « Exemple : »
\providecommand{\exemple}{\textbf{Exemple~:}\ }
% \square utilisable aussi hors mode math (cases à cocher)
\AtBeginDocument{\let\squareMath\square\renewcommand{\square}{\ensuremath{\squareMath}}}
% Mot ou phrase en allemand à l'intérieur d'un paragraphe français
\newcommand{\all}[1]{\ifmmode\text{\textit{#1}}\else\begingroup\selectlanguage{ngerman}\itshape #1\endgroup\fi}
\let\esp\all % alias toléré
\newtcolorbox{exemplebox}[1][]{popbase,colframe=poporange,before upper={\popboxhead{Exemple}{poporange}{#1}}}
\newtcolorbox{definition}[1][]{popbase,colframe=popblue,before upper={\popboxhead{Définition}{popblue}{#1}}}
\newtcolorbox{propriete}[1][]{popbase,colframe=poppurple,before upper={\popboxhead{Propriété}{poppurple}{#1}}}
\newtcolorbox{methode}[1][]{popbase,colframe=popgold,before upper={\popboxhead{Méthode}{popgold}{#1}}}
\newtcolorbox{attention}[1][]{popbase,colframe=poppink,before upper={\popboxhead{Attention}{poppink}{#1}}}
\newtcolorbox{experience}[1][]{popbase,colframe=popblue,colback=popblueL,before upper={\popboxhead{Expérience}{popblue}{#1}}}
% « Le savais-tu ? » : carte violette pleine (contexte, culture, Cameroun)
\newtcolorbox{savaistu}[1][]{popbase,colback=poppurple,colframe=popink,
  fontupper=\popsemi\fontsize{10.4}{14.2}\selectfont\color{white},
  before upper={\poppill{Le savais-tu\,?}{white}{poppurple}\ifx\relax#1\relax\else\hspace{6pt}{\popxbold\color{white}#1}\fi\par\vspace{7pt}}}
% Exercice résolu : énoncé en haut, \tcblower puis la solution sur fond crème
\newcounter{popexo}\newcounter{popexr}
\newtcolorbox{exoresolu}[1][]{popbase,colframe=poporange,colbacklower=popcream,
  segmentation style={popink,line width=1.2pt,dashed},
  before upper={\stepcounter{popexr}\popboxhead{Exercice résolu \thepopexr}{poporange}{#1}},
  before lower={\poppill{Solution}{popink}{popcream}\par\vspace{6pt}}}
% Exercice (non résolu en place) — la correction est dans la partie Corrigés
\newtcolorbox{exercice}[1][]{popbase,colframe=popink,
  before upper={\stepcounter{popexo}\popboxhead{Exercice \thepopexo}{popink}{#1}}}
% Corrigé
\newtcolorbox{corrige}[1][]{popbase,colframe=popgreen,colback=popgreenL,
  before upper={\popboxhead{Corrigé}{popgreen}{#1}}}
% Objectifs / prérequis / bilan
\newtcolorbox{objectifs}{popbase,colframe=popink,colback=poporangeL,
  before upper={\popboxhead{Objectifs de la leçon}{popink}{}}}
\newtcolorbox{prerequis}{popbase,colframe=popink,colback=popblueL,
  before upper={\popboxhead{Prérequis}{popblue}{}}}
\newtcolorbox{bilan}{popbase,colframe=popink,colback=white,boxrule=2.8pt,
  before upper={\popboxhead{Fiche bilan}{poporange}{L'essentiel à connaître pour l'examen}}}
% Figure encadrée avec légende
\newtcolorbox{popfigure}[1][]{enhanced,breakable=false,colback=white,colframe=popline,boxrule=1.4pt,arc=8pt,
  left=10pt,right=10pt,top=10pt,bottom=8pt,before skip=10pt,after skip=12pt,halign=center,
  lower separated=false,halign lower=center,
  fontlower=\popsemi\fontsize{9}{11.5}\selectfont\color{popmuted},
  #1}
\newcommand{\legende}[1]{\par\vspace{6pt}{\popsemi\fontsize{9}{11.5}\selectfont\color{popmuted}#1}}

% ============================================================
% Signature OnBuch+
% ============================================================
\newcommand{\poplogoinline}[1]{{\poplogo\fontsize{#1}{#1}\selectfont\color{popink}OnBuch\color{poporange}+}}
\newcommand{\signatureonbuch}{%
  \begin{tcolorbox}[enhanced,colback=popink,colframe=popink,boxrule=2.2pt,arc=10pt,
      drop shadow=poporange,left=14pt,right=14pt,top=10pt,bottom=10pt,before skip=0pt,after skip=0pt]
    \tikz[baseline=-0.7ex]\node[circle,fill=popgreen,draw=popcream,line width=1.3pt,minimum size=22pt,inner sep=0]{\popcheck{white}};%
    \hspace{9pt}\begin{minipage}[c]{0.62\linewidth}
      {\popxbold\fontsize{12}{13}\selectfont\color{popcream}Signé\ \textbullet\ {\poplogo OnBuch\color{poporange}+}}\par\vspace{2pt}
      {\raggedright\popsemi\fontsize{8}{10.5}\selectfont\color{popline}\MakeUppercase{Équipe pédagogique OnBuch+}\\\MakeUppercase{Conforme au programme officiel MINESEC}\par}
    \end{minipage}\hfill
    {\poplogo\fontsize{22}{22}\selectfont\color{popcream}OnBuch\color{poporange}+}
  \end{tcolorbox}%
}

% ============================================================
% Page de garde
% #1 titre · #2 matière · #3 niveau · #4 module · #5 n° leçon · #6 durée ·
% #7 description / objectifs (texte ou itemize)
% ============================================================
\newcommand{\pagegarde}[7]{%
  \thispagestyle{empty}
  \newgeometry{a4paper,margin=1.7cm,top=1.6cm,bottom=1.4cm}%
  \begin{tikzpicture}[remember picture,overlay]
    \fill[poporange!18] ([xshift=-3.2cm,yshift=-3.6cm]current page.north east) circle (4.6cm);
    \fill[poppurple!14] ([xshift=2.4cm,yshift=7.6cm]current page.south west) circle (3.8cm);
  \end{tikzpicture}%
  {\poplogo\fontsize{30}{30}\selectfont\color{popink}OnBuch\color{poporange}+}\hfill
  \raisebox{0.6ex}{\poppill{Cours signé \textbullet\ Premium}{popink}{popcream}}\par
  \vspace{1pt}\popsquiggle{poppurple}\par
  \vspace{8pt}
  \begin{tcolorbox}[enhanced,colback=poporange,colframe=popink,boxrule=2.8pt,arc=13pt,
      drop shadow=popink,left=18pt,right=18pt,top=12pt,bottom=13pt,after skip=10pt]
    \poppill{#2 \textbullet\ #3}{popink}{popcream}\hspace{5pt}\poppill{#4}{white}{popink}\par\vspace{8pt}
    {\popdisplay\fontsize{14}{14}\selectfont\color{popink}LEÇON #5}\par\vspace{4pt}
    {\raggedright\hyphenpenalty=10000\exhyphenpenalty=10000\popdisplay\fontsize{25}{27}\selectfont\color{popcream}\MakeUppercase{#1}\par}
    \vspace{8pt}
    {\popxbold\fontsize{9.5}{9.5}\selectfont\color{popink}\MakeUppercase{Durée conseillée : #6}}
  \end{tcolorbox}
  \begin{tcolorbox}[enhanced,colback=white,colframe=popink,boxrule=2.2pt,arc=10pt,
      drop shadow=popink,left=16pt,right=16pt,top=10pt,bottom=10pt,after skip=8pt]
    \poppill{Dans cette leçon}{poppurple}{white}\par\vspace{5pt}
    \popsemi\fontsize{9.3}{12.6}\selectfont\color{popink}#7
  \end{tcolorbox}
  \noindent\begin{minipage}[t]{0.487\linewidth}
    \begin{tcolorbox}[enhanced,colback=popgreen,colframe=popink,boxrule=2.2pt,arc=10pt,
        drop shadow=popink,left=11pt,right=11pt,top=10pt,bottom=10pt,height=3.1cm,valign=center]
      \tcbox[colback=white,colframe=popink,boxrule=1.3pt,arc=5pt,boxsep=0pt,nobeforeafter,
        left=3pt,right=3pt,top=3pt,bottom=3pt,valign=center]{\qrcode[height=1.9cm]{https://play.google.com/store/apps/details?id=com.luvvix.onbuch&hl=fr}}%
      \hspace{8pt}\begin{minipage}[c]{0.5\linewidth}
        {\popdisplay\fontsize{11}{12.5}\selectfont\color{white}\MakeUppercase{Télécharge OnBuch}}\par\vspace{4pt}
        {\raggedright\popsemi\fontsize{8.4}{11}\selectfont\color{white}Gratuit \textbullet\ Google Play\\Tuteur IA, annales, concours, résultats\par}
      \end{minipage}
    \end{tcolorbox}
  \end{minipage}\hfill
  \begin{minipage}[t]{0.487\linewidth}
    \begin{tcolorbox}[enhanced,colback=poppurple,colframe=popink,boxrule=2.2pt,arc=10pt,
        drop shadow=popink,left=11pt,right=11pt,top=10pt,bottom=10pt,height=3.1cm,valign=center]
      \tikz[baseline=-0.7ex]\node[circle,fill=white,draw=popink,line width=1.3pt,minimum size=1.6cm,inner sep=0]{\popdisplay\fontsize{24}{24}\selectfont\color{poppurple}+};%
      \hspace{8pt}\begin{minipage}[c]{0.6\linewidth}
        {\popdisplay\fontsize{11}{12.5}\selectfont\color{white}\MakeUppercase{Passe à OnBuch+}}\par\vspace{4pt}
        {\raggedright\popsemi\fontsize{8.4}{11}\selectfont\color{white}Tous les cours signés, Tuteur IA illimité, fiches PDF — onglet OnBuch+ de l'app\par}
      \end{minipage}
    \end{tcolorbox}
  \end{minipage}
  \vspace{8pt}
  \popcontenu
  \vfill
  \signatureonbuch
  \restoregeometry
  \newpage
}

% Rangée « Ce cours contient » (page de garde)
\newcommand{\poptile}[3]{% #1 couleur #2 gros texte #3 légende
  \begin{tcolorbox}[enhanced,colback=#1,colframe=popink,boxrule=2pt,arc=9pt,shadow={2.5pt}{-2.5pt}{0pt}{popink},
      left=6pt,right=6pt,top=9pt,bottom=9pt,halign=center,valign=center,height=1.75cm,width=\linewidth,nobeforeafter]
    {\popdisplay\fontsize{12.5}{13}\selectfont\color{white}\MakeUppercase{#2}}\par\vspace{3pt}
    {\centering\popsemi\fontsize{7.8}{9.5}\selectfont\color{white}#3\par}
  \end{tcolorbox}}
\newcommand{\popcontenu}{%
  \par\noindent{\popxboldls\fontsize{7.5}{7.5}\selectfont\color{popmuted}CE COURS SIGNÉ CONTIENT}\par\vspace{7pt}
  \noindent\begin{minipage}[t]{0.235\linewidth}\poptile{popblue}{Cours}{complet, illustré, pas à pas}\end{minipage}\hfill
  \begin{minipage}[t]{0.235\linewidth}\poptile{poporange}{Méthodes}{et exercices résolus}\end{minipage}\hfill
  \begin{minipage}[t]{0.235\linewidth}\poptile{poppink}{Exercices}{type examen + corrigés}\end{minipage}\hfill
  \begin{minipage}[t]{0.235\linewidth}\poptile{popgreen}{Fiche bilan}{l'essentiel à retenir}\end{minipage}\par}

% Sommaire
\newtcolorbox{sommaire}{enhanced,colback=white,colframe=popink,boxrule=2.2pt,arc=10pt,
  drop shadow=popink,left=18pt,right=18pt,top=13pt,bottom=13pt,before skip=6pt,after skip=20pt,
  before upper={\poppill{Sommaire}{poppurple}{white}\par\vspace{9pt}\popsemi\fontsize{10.6}{14.5}\selectfont\color{popink}}}

% Signature de fin de cours — poussée tout en bas de la dernière page
\newcommand{\findecours}{%
  \par\vfill\nopagebreak
  \begin{center}\popsquiggle{poporange}\end{center}\vspace{4pt}
  \signatureonbuch
}
\AtBeginDocument{\def\llbracket{\mathopen{[\mkern-3.5mu[}}\def\rrbracket{\mathclose{]\mkern-3.5mu]}}} % intervalles d'entiers (glyphe ⟦ absent de la police)
% Glyphes absents de Fira Math : redéfinis à partir de glyphes disponibles
\AtBeginDocument{%
  \def\setminus{\mathbin{\backslash}}\let\smallsetminus\setminus
  \def\vdots{\mathord{\vbox{\baselineskip3.2pt\lineskiplimit0pt\kern4pt\hbox{.}\hbox{.}\hbox{.}}}}%
  \def\ddots{\mathinner{\mkern1mu\raise6.5pt\hbox{.}\mkern2mu\raise3.6pt\hbox{.}\mkern2mu\raise0.7pt\hbox{.}\mkern1mu}}%
  \def\triangle{\mathord{\scalebox{0.9}{$\symup{\Delta}$}}}%
  \def\nabla{\mathord{\raisebox{\depth}{\scalebox{1}[-1]{$\symup{\Delta}$}}}}%
  \def\checkmark{\popcheck{popdark}}%
  \def\star{\mathbin{\ast}}\def\bigstar{\mathord{\ast}}%
  \def\diamond{\mathbin{\scalebox{0.6}{$\Diamond$}}}%
  \def\bigcup{\mathop{\mathchoice{\raisebox{-0.3em}{\scalebox{1.7}{$\cup$}}}{\scalebox{1.25}{$\cup$}}{\cup}{\cup}}\displaylimits}%
  \def\bigcap{\mathop{\mathchoice{\raisebox{-0.3em}{\scalebox{1.7}{$\cap$}}}{\scalebox{1.25}{$\cap$}}{\cap}{\cap}}\displaylimits}%
  \def\lesseqqgtr{\mathrel{\lessgtr}}\def\bigoplus{\mathop{\oplus}\displaylimits}%
}
\providecommand{\ssi}{si et seulement si} % abréviation fréquente
\AtBeginDocument{\providecommand{\wideparen}{\overparen}} % arc de cercle

%% FIGFIX-BEGIN v5 — correctifs globaux des figures (docs/onbuch-generation/tools/figfix)
\usepackage{adjustbox}\usepackage{etoolbox}\tcbuselibrary{raster}
% toute figure TikZ/pgfplots plus large que la ligne est réduite à la largeur disponible
\BeforeBeginEnvironment{tikzpicture}{\begin{adjustbox}{max width=\linewidth}}
\AfterEndEnvironment{tikzpicture}{\end{adjustbox}}
% légendes pgfplots lisibles par-dessus les courbes
\pgfplotsset{every axis/.append style={legend style={fill=white,fill opacity=0.92,text opacity=1,draw=black!15,inner sep=3pt}}}
% étiquettes d'arêtes (syntaxe edge["..."]) avec fond blanc
\tikzset{every edge quotes/.append style={fill=white,inner sep=1.2pt}}
%% FIGFIX-END
\begin{document}

\pagegarde{La vie culturelle 2 (la musique et la danse)}{Allemand}{Tle A}{Chapitre 5 : Citoyenneté}{AL23}{2 h}{Cette leçon de type « texte » propose la lecture et le commentaire d'un texte original sur la musique et la danse au Cameroun et en Allemagne. Elle consolide le vocabulaire thématique, la comparaison avec « wie » et « als ob », la déclinaison de l'adjectif et le subjonctif II, avant une production guidée.
\vspace{4pt}
\begin{multicols}{2}\small
\begin{itemize}
\item À la fin de cette leçon, je sais comprendre globalement puis en détail un texte allemand sur la musique et la danse.
\item À la fin de cette leçon, je sais employer le vocabulaire thématique (die Musik, der Tanz, das Konzert, die Trommel, der Sänger, das Fest) avec les articles et les pluriels corrects.
\item À la fin de cette leçon, je sais construire des comparaisons avec « wie » et des comparaisons irréelles avec « als ob » + Konjunktiv II.
\item À la fin de cette leçon, je sais décliner correctement l'adjectif après article défini, indéfini et sans article.
\item À la fin de cette leçon, je sais former et employer le subjonctif II (présent et passé) pour exprimer l'irréel et le souhait.
\item À la fin de cette leçon, je sais répondre à des questions de compréhension en phrases complètes.
\end{itemize}\end{multicols}}

\begin{sommaire}
\begin{multicols}{2}
\begin{enumerate}
\item Texte support : Musik und Tanz in Kamerun und in Deutschland
\item Compréhension du texte
\item Lexique thématique : die Musik und der Tanz
\item Grammaire 1 : la comparaison avec « wie » et « als ob »
\item Grammaire 2 : la déclinaison de l'adjectif et le subjonctif II
\item Méthode du commentaire et production guidée
\item Activité d'intégration
\item Méthodes et exercices résolus
\item Exercices
\item Corrigés des exercices
\item Fiche bilan
\end{enumerate}
\end{multicols}
\end{sommaire}

\begin{objectifs}
\begin{itemize}
\item À la fin de cette leçon, je sais comprendre globalement puis en détail un texte allemand sur la musique et la danse.
\item À la fin de cette leçon, je sais employer le vocabulaire thématique (die Musik, der Tanz, das Konzert, die Trommel, der Sänger, das Fest) avec les articles et les pluriels corrects.
\item À la fin de cette leçon, je sais construire des comparaisons avec « wie » et des comparaisons irréelles avec « als ob » + Konjunktiv II.
\item À la fin de cette leçon, je sais décliner correctement l'adjectif après article défini, indéfini et sans article.
\item À la fin de cette leçon, je sais former et employer le subjonctif II (présent et passé) pour exprimer l'irréel et le souhait.
\item À la fin de cette leçon, je sais répondre à des questions de compréhension en phrases complètes.
\item À la fin de cette leçon, je sais rédiger un résumé et un court commentaire structuré (introduction, développement, conclusion).
\item À la fin de cette leçon, je sais produire un court paragraphe sur une fête musicale camerounaise en réutilisant les acquis de la leçon.
\end{itemize}
\end{objectifs}

\begin{prerequis}
\begin{itemize}
\item Les articles définis et indéfinis et les quatre cas (Nominatif, Accusatif, Datif, Génitif) vus en Première
\item Le présent et le prétérit des verbes faibles et forts
\item Les pronoms personnels et possessifs
\item Le vocabulaire de base de la vie culturelle 1 (das Theater, der Film, das Museum)
\item La phrase subordonnée et la place du verbe conjugué en fin de proposition
\end{itemize}
\end{prerequis}

\newpage

\coursec{Texte support : Musik und Tanz in Kamerun und in Deutschland}

\begin{experience}[Situation d'accroche : la fête de fin d'année]
C'est la fête de fin d'année au lycée de Yaoundé : les tambours résonnent, les élèves dansent le bikutsi et chantent en chœur. En Allemagne aussi, la musique et la danse rythment la vie : concerts classiques à Vienne, fêtes populaires en Bavière, carnaval à Cologne. Comment parler de ces moments forts en allemand ? Comment comparer une danse à une autre, décrire un chanteur avec les bons adjectifs et exprimer ce qu'on ferait si on était sur scène ? Cette leçon nous donne tous les outils pour raconter et commenter la vie musicale, ici et dans les pays germanophones.
\end{experience}

\textbf{Problématique.} Comment comprendre, raconter et commenter en allemand un événement musical ou une fête, au Cameroun comme dans les pays germanophones ?

\textbf{Objectifs de la leçon.} À la fin de cette leçon, tu seras capable de : comprendre un texte sur la musique et la danse ; employer le vocabulaire du thème avec le bon article et le bon pluriel ; construire des comparaisons avec \all{wie} et \all{als ob} ; accorder correctement l'adjectif épithète ; exprimer un souhait irréel au subjonctif II.

\courssub{Lecture globale du texte}

\begin{methode}[Comment lire un texte au Bac]
\begin{enumerate}
\item Lire le \cle{titre} et la \cle{première phrase} : ils annoncent le thème et le point de vue.
\item Repérer \cle{qui parle} (le narrateur), \cle{de quoi} il parle (le thème) et \cle{où} se situe la scène (le lieu).
\item Lire les \cle{questions} de compréhension avant la deuxième lecture : on sait ainsi ce qu'on cherche.
\item Souligner les \cle{mots-clés} : noms, verbes d'action, connecteurs, marques de comparaison (\all{wie}, \all{als ob}).
\end{enumerate}
\end{methode}

\begin{textebox}[Musik und Tanz in Kamerun und in Deutschland]
Letztes Wochenende war ein großes Fest in meinem Dorf. Schon am frühen Morgen hörte man die Trommeln, die so laut klangen, dass das ganze Viertel wach wurde. Die alten Trommler spielten mit einer Kraft, die mich immer wieder überrascht. Um zehn Uhr versammelten sich alle Leute auf dem großen Platz vor dem Haus des Dorfchefs. Die Frauen trugen bunte Kleider, und die Kinder liefen singend zwischen den Gruppen hin und her.

Dann begannen die traditionellen Tänze. Die Tänzer bewegten sich so schnell, als ob ihre Füße den Boden nicht berührten. Der bekannte Sänger aus der Nachbarstadt trat auf, und die Menge jubelte, als ob sie einen König sähe. Seine Stimme war stark wie ein Fluss, und alle sangen den Refrain mit. Wir amüsierten uns bis spät in die Nacht.

In Deutschland ist die Musik anders, erzählt mein Onkel, der in München lebt. Dort besucht man oft ein klassisches Konzert in einem eleganten Saal, oder man tanzt auf einem Volksfest zu moderner Musik. Die Leute klatschen höflich, und niemand tanzt auf den Tischen — fast niemand, lacht mein Onkel. Wenn ich in Deutschland wäre, würde ich gern das Oktoberfest besuchen und ein echtes Konzert hören. Ich würde auch den deutschen Freunden unsere Trommeln zeigen.

Ob in Kamerun oder in Deutschland: Musik verbindet die Menschen wie eine universelle Sprache. Man braucht keine Wörter, um einen Rhythmus zu fühlen. Wenn die Trommeln spielen, versteht jeder die Einladung zum Tanz.
\end{textebox}

\begin{definition}[Glossaire du texte]
\begin{itemize}
\item \all{die Trommel, -n} = le tambour, le tam-tam ; \all{der Trommler, -} = le joueur de tam-tam
\item \all{der Sänger, -} / \all{die Sängerin, -nen} = le chanteur / la chanteuse
\item \all{das Konzert, -e} = le concert ; \all{das Fest, -e} = la fête ; \all{das Volksfest, -e} = la fête populaire
\item \all{der Rhythmus, die Rhythmen} = le rythme ; \all{der Tanz, „-e} = la danse
\item \all{auftreten} (tritt auf, trat auf, ist aufgetreten) = se produire sur scène, monter sur scène
\item \all{sich amüsieren} = s'amuser ; \all{jubeln} = jubiler, acclamer
\item \all{sich versammeln} = se rassembler ; \all{der Dorfchef, -s} = le chef du village
\item \all{die Stimme, -n} = la voix ; \all{der Refrain, -s} = le refrain ; \all{klatschen} = applaudir
\item \all{verbinden} (verbindet, verband, hat verbunden) = unir, relier
\end{itemize}
\end{definition}

\begin{attention}[Ne traduis pas mot à mot !]
Au premier contact, cherche le \cle{sens global} du texte, pas la traduction de chaque mot. Le glossaire est là pour les mots nouveaux : utilise-le. Un mot inconnu peut souvent être deviné par le contexte : \all{die Menge jubelte} — la foule fait quoi quand un chanteur arrive ? Elle \emph{jubile}. Attention aussi aux faux amis : \all{das Konzert} = le concert (pas « le concours »), \all{der Saal} = la salle (pas « le sel »).
\end{attention}

\courssub{Repérage : thème, narrateur, lieu, ton}

Après une première lecture silencieuse, puis une lecture à voix haute par le professeur, on répond aux quatre questions d'orientation :

\begin{center}
\begin{tabularx}{\linewidth}{|X|X|}
\hline
\rowcolor{popblueL} \textbf{Question} & \textbf{Réponse (avec preuve du texte)} \\
\hline
Quel est le thème ? & La musique et la danse, au Cameroun et en Allemagne (\all{Musik und Tanz in Kamerun und in Deutschland}). \\
\hline
Qui parle ? & Un jeune Camerounais, narrateur à la première personne (\all{in meinem Dorf}, \all{wir amüsierten uns}). \\
\hline
Où et quand ? & Dans son village, le week-end dernier (\all{Letztes Wochenende}) ; puis évocation de l'Allemagne (Munich). \\
\hline
Quel est le ton ? & Joyeux, fier et admiratif : \all{großes Fest}, \all{jubelte}, \all{wir amüsierten uns}. \\
\hline
\end{tabularx}
\end{center}

\begin{aretenir}[Les quatre repères de la lecture globale]
Toute compréhension de texte commence par quatre repères : le \cle{thème} (de quoi ?), le \cle{narrateur} (qui parle ?), le \cle{lieu et le moment} (où et quand ?), le \cle{ton} (joyeux, triste, critique, admiratif ?). Ces quatre éléments doivent toujours être justifiés par des citations du texte.
\end{aretenir}

\courssub{Lecture détaillée : observer la langue}

Relisons le texte en soulignant trois phénomènes grammaticaux qui seront étudiés dans les sections suivantes.

\textbf{1. Les comparaisons avec \all{wie} et \all{als ob}.}

\all{Die Tänzer bewegten sich so schnell, als ob ihre Füße den Boden nicht berührten.} « Les danseurs se déplaçaient si vite, comme si leurs pieds ne touchaient pas le sol. »

\all{Die Menge jubelte, als ob sie einen König sähe.} « La foule jubilait, comme si elle voyait un roi. »

\all{Seine Stimme war stark wie ein Fluss.} « Sa voix était forte comme un fleuve. »

\all{Musik verbindet die Menschen wie eine universelle Sprache.} « La musique relie les hommes comme une langue universelle. »

On observe : \all{wie} + nom introduit une comparaison directe (« comme ») ; \all{als ob} introduit une subordonnée de comparaison irréelle, avec le verbe à la fin et au \cle{subjonctif II} (\all{sähe}, \all{berührten}).

\textbf{2. Les adjectifs épithètes.}

\all{ein großes Fest} (une grande fête), \all{der bekannte Sänger} (le chanteur célèbre), \all{bunte Kleider} (des habits colorés), \all{die traditionellen Tänze} (les danses traditionnelles), \all{in einem eleganten Saal} (dans une salle élégante). Chaque adjectif placé devant le nom prend une \cle{désinence} qui dépend de l'article, du genre, du nombre et du cas : c'est la déclinaison de l'adjectif. Compare : \all{das Fest ist groß} (adjectif attribut, invariable) et \all{ein großes Fest} (adjectif épithète, décliné).

\textbf{3. Le subjonctif II.}

\all{Wenn ich in Deutschland wäre, würde ich gern das Oktoberfest besuchen.} « Si j'étais en Allemagne, je visiterais volontiers l'Oktoberfest. » \all{Ich würde auch den deutschen Freunden unsere Trommeln zeigen.} « Je montrerais aussi nos tam-tams aux amis allemands. » Le narrateur exprime un souhait irréel : \all{wäre} (subjonctif II de \all{sein}) et \all{würde} + infinitif.

\courssub{Relevé des champs lexicaux}

Le texte organise son vocabulaire autour de quatre champs lexicaux. Complétons-les avec des mots du glossaire et du thème :

\begin{center}
\begin{tabularx}{\linewidth}{|X|X|}
\hline
\rowcolor{popblueL} \textbf{Champ lexical} & \textbf{Mots relevés dans le texte} \\
\hline
Les instruments & die Trommel, die Trommeln, die Stimme \\
\hline
Les personnes & der Sänger, die Sängerin, der Trommler, die Tänzer, die Menge, das Publikum \\
\hline
Les événements & das Fest, das Volksfest, das Konzert, der Auftritt, das Oktoberfest \\
\hline
Les sentiments et réactions & jubeln, sich amüsieren, überraschen, klatschen, die Freude \\
\hline
\end{tabularx}
\end{center}

\begin{popfigure}
\begin{center}\tcbox[colback=popblue,colframe=popblue,coltext=white,arc=9pt,boxsep=4pt,left=6pt,right=6pt]{\bfseries\small Musik und Tanz}\end{center}
\vspace{-2pt}
\begin{tcbraster}[raster columns=2,raster equal height=rows,raster column skip=6pt,raster row skip=6pt,enhanced,blankest]
\begin{tcolorbox}[enhanced,colback=white,colframe=poporange,boxrule=0.9pt,arc=3pt,title={Instrumente},fonttitle=\bfseries\scriptsize,coltitle=white,colbacktitle=poporange,left=4pt,right=4pt,top=2pt,bottom=2pt,boxsep=1pt,toptitle=1.5pt,bottomtitle=1.5pt,halign=left]
\begin{itemize}[nosep,leftmargin=*,itemsep=1pt,font=\scriptsize]
\item die Trommel
\item die Stimme
\item die Gitarre
\end{itemize}
\end{tcolorbox}
\begin{tcolorbox}[enhanced,colback=white,colframe=popgreen,boxrule=0.9pt,arc=3pt,title={Personen},fonttitle=\bfseries\scriptsize,coltitle=white,colbacktitle=popgreen,left=4pt,right=4pt,top=2pt,bottom=2pt,boxsep=1pt,toptitle=1.5pt,bottomtitle=1.5pt,halign=left]
\begin{itemize}[nosep,leftmargin=*,itemsep=1pt,font=\scriptsize]
\item der Sänger
\item die Tänzer
\item das Publikum
\end{itemize}
\end{tcolorbox}
\begin{tcolorbox}[enhanced,colback=white,colframe=poppurple,boxrule=0.9pt,arc=3pt,title={Ereignisse},fonttitle=\bfseries\scriptsize,coltitle=white,colbacktitle=poppurple,left=4pt,right=4pt,top=2pt,bottom=2pt,boxsep=1pt,toptitle=1.5pt,bottomtitle=1.5pt,halign=left]
\begin{itemize}[nosep,leftmargin=*,itemsep=1pt,font=\scriptsize]
\item das Fest
\item das Konzert
\item das Volksfest
\end{itemize}
\end{tcolorbox}
\begin{tcolorbox}[enhanced,colback=white,colframe=poppink,boxrule=0.9pt,arc=3pt,title={Gefühle},fonttitle=\bfseries\scriptsize,coltitle=white,colbacktitle=poppink,left=4pt,right=4pt,top=2pt,bottom=2pt,boxsep=1pt,toptitle=1.5pt,bottomtitle=1.5pt,halign=left]
\begin{itemize}[nosep,leftmargin=*,itemsep=1pt,font=\scriptsize]
\item jubeln
\item sich amüsieren
\item die Freude
\end{itemize}
\end{tcolorbox}
\end{tcbraster}

\legende{Carte mentale du vocabulaire : quatre champs lexicaux autour de « Musik und Tanz ».}
\end{popfigure}

\begin{savaistu}[Musique en Allemagne : une passion nationale]
L'Allemagne compte plus de 80 orchestres professionnels permanents, un record en Europe : presque chaque grande ville possède le sien. Et l'Oktoberfest de Munich, la plus grande fête populaire du monde, attire chaque année environ 6 millions de visiteurs qui dansent et chantent au son des orchestres bavarois. La musique y est, comme chez nous au Cameroun avec le makossa ou le bikutsi, un langage qui rassemble.
\end{savaistu}

\begin{exoresolu}[Questions de compréhension globale]
Énoncé. Réponds en allemand, avec une phrase complète : 1) Wer erzählt die Geschichte? 2) Wo findet das Fest statt? 3) Warum jubelt die Menge? 4) Was würde der Erzähler in Deutschland machen?
\tcblower
Solution détaillée. 1) \all{Ein junger Kameruner erzählt die Geschichte.} (Un jeune Camerounais raconte l'histoire — indice : \all{in meinem Dorf}.) 2) \all{Das Fest findet in seinem Dorf statt, auf dem großen Platz vor dem Haus des Dorfchefs.} 3) \all{Die Menge jubelt, weil der bekannte Sänger auftritt.} (La foule jubile parce que le chanteur célèbre monte sur scène — attention au verbe à particule : \all{auftritt}, la particule \all{auf} se détache et va en fin de proposition.) 4) \all{Er würde das Oktoberfest besuchen, ein echtes Konzert hören und den deutschen Freunden die Trommeln zeigen.} Remarque la forme \all{würde} + infinitif à la fin : c'est le subjonctif II, obligatoire pour un souhait irréel.
\end{exoresolu}

\begin{exercice}[À toi de jouer]
1) Souligne dans le texte toutes les comparaisons et recopie-les en deux colonnes : \all{wie} / \all{als ob}. 2) Relève cinq adjectifs épithètes avec leur nom et leur article. 3) Réponds : \all{Wie ist die Musik in Deutschland laut dem Onkel?} 4) Produis trois phrases : décris une fête de ton quartier avec \all{die Trommeln}, \all{der Sänger} et \all{sich amüsieren}.
\end{exercice}

\begin{corrige}[Exercice « À toi de jouer »]
1) Avec \all{wie} : \all{stark wie ein Fluss} ; \all{wie eine universelle Sprache}. Avec \all{als ob} : \all{als ob ihre Füße den Boden nicht berührten} ; \all{als ob sie einen König sähe}. 2) Exemples : \all{ein großes Fest}, \all{der bekannte Sänger}, \all{bunte Kleider}, \all{die traditionellen Tänze}, \all{einem eleganten Saal}. 3) \all{In Deutschland besucht man oft ein klassisches Konzert oder tanzt auf einem Volksfest; die Leute klatschen höflich.} 4) Modèle : \all{Die Trommeln spielen laut. Der Sänger tritt um acht Uhr auf. Wir amüsieren uns die ganze Nacht.}
\end{corrige}

\begin{aretenir}[Bilan de la section]
Le texte \all{Musik und Tanz in Kamerun und in Deutschland} est un récit à la première personne, au ton joyeux, qui compare une fête de village camerounaise et la vie musicale allemande. Nous avons repéré le thème, le narrateur, le lieu et le ton, enrichi notre vocabulaire (instruments, personnes, événements, sentiments) et observé trois phénomènes grammaticaux — la comparaison avec \all{wie} et \all{als ob}, la déclinaison de l'adjectif épithète et le subjonctif II — que nous étudierons en détail dans les sections suivantes.
\end{aretenir}

\coursec{Compréhension du texte}

Dans la section précédente, nous avons découvert et lu le texte support \emph{Musik und Tanz in Kamerun und in Deutschland}, dans lequel un narrateur camerounais raconte une fête de village avec des tambours et compare cette expérience à la vie musicale en Allemagne. Nous allons maintenant exploiter ce texte : comprendre d'abord le sens général, puis les détails, justifier nos réponses par des citations, et enfin résumer le texte avec nos propres mots. C'est exactement la démarche exigée au baccalauréat.

\courssub{Rappel du texte support}

Avant de répondre aux questions, relisons le texte. Le voici à nouveau, avec son glossaire.

\begin{textebox}[Musik und Tanz in Kamerun und in Deutschland]
Letztes Wochenende war ein großes Fest in meinem Dorf bei Yaoundé. Die ganze Familie war da: meine Eltern, meine Geschwister, meine Onkel und Tanten und viele Nachbarn. Am Abend begann die Musik. Die Trommeln klangen so laut, dass alle Leute tanzten. Die Menge klatschte und sang mit den Sängern. Mein Großvater erzählte mir, dass diese Musik sehr alt ist: „Unsere Väter und Großväter haben schon auf diesen Trommeln gespielt“, sagte er stolz.

Ich tanzte mit meinen Freunden bis Mitternacht. Die Tänzer bewegten sich wie Vögel im Wind. Es war, als ob die ganze Welt nur aus Musik bestünde. Ein berühmter Sänger aus Douala trat auch auf. Alle waren glücklich, und niemand wollte nach Hause gehen.

In Deutschland ist die Musik anders, aber auch sehr wichtig. Viele junge Deutsche hören Popmusik, Rock oder Hip-Hop. In den großen Städten wie Berlin, München oder Wien gibt es viele Konzerte. Wenn ich in Deutschland wäre, würde ich ein klassisches Konzert besuchen, denn Deutschland ist das Land von Bach, Mozart und Beethoven. Ich würde auch gern ein Volksfest sehen, zum Beispiel das Oktoberfest in München, wo die Leute singen und tanzen.

Musik und Tanz verbinden die Menschen überall auf der Welt. Ob in Kamerun oder in Deutschland: Wenn die Musik spielt, tanzen alle, als ob sie keine Sorgen hätten. Musik ist wie eine Sprache, die jeder versteht.
\end{textebox}

\textbf{Glossaire (Wortschatz)}

\begin{center}
\begin{tabularx}{\linewidth}{|X|X|}
\hline
\rowcolor{popblueL} \textbf{Deutsch} & \textbf{Français} \\
\hline
das Fest, die Feste & la fête \\
\hline
das Dorf, die Dörfer & le village \\
\hline
die Trommel, die Trommeln & le tambour, le tam-tam \\
\hline
klingen (klang, hat geklungen) & résonner, sonner \\
\hline
die Menge, die Mengen & la foule \\
\hline
klatschen & applaudir \\
\hline
der Sänger, die Sänger / die Sängerin, die Sängerinnen & le chanteur / la chanteuse \\
\hline
auftreten (trat auf, ist aufgetreten) & se produire (sur scène) \\
\hline
berühmt & célèbre \\
\hline
das Konzert, die Konzerte & le concert \\
\hline
das Volksfest, die Volksfeste & la fête populaire \\
\hline
die Sorge, die Sorgen & le souci \\
\hline
verbinden (verband, hat verbunden) & unir, relier \\
\hline
\end{tabularx}
\end{center}

\courssub{La méthode de compréhension : du global au détail}

Comprendre un texte allemand ne se fait pas en une seule lecture. On procède en étapes, comme le montre l'organigramme suivant.

\begin{popfigure}
\begin{tikzpicture}[node distance=0.55cm,
box/.style={draw=popblue, fill=popblueL, rounded corners, align=center, text width=4.6cm, minimum height=0.9cm, font=\small},
arr/.style={-{Stealth[length=2.5mm]}, thick, poporange}]
\node[box, align=center] (a){\textbf{1. Lecture globale}\\ lire vite, saisir le thème général};
\node[box, below=of a, align=center] (b){\textbf{2. Questions}\\ lire les questions avant de relire};
\node[box, below=of b, align=center] (c){\textbf{3. Lecture détaillée}\\ souligner les mots-clés, repérer les réponses};
\node[box, below=of c, align=center] (d){\textbf{4. Réponses rédigées}\\ phrases complètes, citations entre guillemets};
\node[box, below=of d, align=center] (e){\textbf{5. Résumé}\\ reformuler avec des connecteurs};
\draw[arr] (a) -- (b);
\draw[arr] (b) -- (c);
\draw[arr] (c) -- (d);
\draw[arr] (d) -- (e);
\end{tikzpicture}
\legende{Organigramme de la méthode de compréhension d'un texte}
\end{popfigure}

\begin{methode}[Comment répondre à une question de compréhension]
\begin{enumerate}
\item Lire la question et repérer le mot interrogatif : \all{wer} (qui), \all{was} (que), \all{wo} (où), \all{wann} (quand), \all{wie} (comment), \all{warum} (pourquoi).
\item Repérer dans le texte la phrase qui contient la réponse.
\item \textbf{Reformuler la question dans la réponse} : reprendre les mots de la question pour construire une phrase complète.
\item Rédiger \textbf{toujours une phrase complète} (sujet + verbe conjugué en 2\textsuperscript{e} position), jamais un seul mot.
\item Vérifier l'ordre des mots allemands : verbe en 2\textsuperscript{e} position, particule séparable en fin de phrase.
\end{enumerate}
\end{methode}

\begin{attention}[Jamais de réponse en un seul mot]
À la question \all{Wo fand das Fest statt?}, la réponse \all{In einem Dorf} seule est insuffisante au Bac. Il faut écrire : \all{Das Fest fand in einem Dorf bei Yaoundé statt.} Remarquez la particule séparable \all{statt-} du verbe \all{stattfinden} renvoyée en fin de phrase : c'est un piège classique pour les francophones.
\end{attention}

\courssub{Compréhension globale (Globales Verstehen)}

La lecture globale répond à trois questions simples, sans entrer dans les détails.

\begin{exemplebox}[Questions de compréhension globale et réponses modèles]
\textbf{1. Worum geht es in dem Text ?} (De quoi parle le texte ?)\par
\all{Der Text handelt von Musik und Tanz in Kamerun und in Deutschland.} « Le texte parle de la musique et de la danse au Cameroun et en Allemagne. »

\textbf{2. Wer erzählt die Geschichte ?} (Qui raconte l'histoire ?)\par
\all{Ein junger Kameruner erzählt die Geschichte; er berichtet von einem Fest in seinem Dorf.} « Un jeune Camerounais raconte l'histoire ; il rapporte une fête de son village. » (On le devine grâce à \all{mein Dorf bei Yaoundé} et à la première personne \all{ich}.)

\textbf{3. Welche zwei Länder werden verglichen ?} (Quels sont les deux pays comparés ?)\par
\all{Die zwei Länder sind Kamerun und Deutschland.} « Les deux pays sont le Cameroun et l'Allemagne. »
\end{exemplebox}

\courssub{Compréhension détaillée (Detailliertes Verstehen)}

Pour la lecture détaillée, on cherche des informations précises. Observons d'abord deux phrases clés du texte et leur traduction.

\begin{exemplebox}[Phrases clés du texte, traduites]
\all{Die Trommeln klangen so laut, dass alle Leute tanzten.} « Les tambours résonnaient si fort que tout le monde dansait. » (construction \all{so ... dass} = « si ... que », verbe de la subordonnée en fin de phrase)\par\medskip
\all{Wenn ich in Deutschland wäre, würde ich ein klassisches Konzert besuchen.} « Si j'étais en Allemagne, j'assisterais à un concert classique. » (subjonctif II : hypothèse irréelle ; \all{wäre} = imparfait du subjonctif de \all{sein}, \all{würde besuchen} = conditionnel)
\end{exemplebox}

\begin{exoresolu}[Questions de compréhension détaillée]
\textbf{Énoncé.} Réponds par des phrases complètes :\par
a) Wo fand das Fest statt ?\par
b) Was machten die Leute, als die Trommeln klangen ?\par
c) Wie reagierte die Menge ?\par
d) Was würde der Erzähler in Deutschland besuchen ?
\tcblower
\textbf{Solution détaillée.}\par
a) \all{Das Fest fand in einem Dorf bei Yaoundé statt.} — On reprend la question (\all{das Fest ... stattfinden}) et on insère le lieu ; la particule \all{statt} va en fin de phrase.\par
b) \all{Als die Trommeln klangen, tanzten alle Leute.} — Le texte dit : \all{Die Trommeln klangen so laut, dass alle Leute tanzten.} On reformule avec la subordonnée temporelle \all{als} (verbe conjugué en fin de subordonnée ; dans la principale qui suit, le verbe reste en 2\textsuperscript{e} position, d'où \all{tanzten alle Leute}).\par
c) \all{Die Menge klatschte und sang mit den Sängern.} — Réponse presque littérale ; attention, \all{die Menge} est singulier en allemand, donc \all{klatschte} et non \all{klatschten}.\par
d) \all{In Deutschland würde der Erzähler ein klassisches Konzert besuchen.} — On conserve le subjonctif II (\all{würde ... besuchen}) car il s'agit d'un souhait irréel ; le complément de lieu en tête impose l'inversion du sujet.
\end{exoresolu}

\begin{attention}[La question « Wann ? » et la place du complément de temps]
À la question \all{Wann ?}, on répond avec un complément de temps \textbf{en tête ou en fin de phrase}, et le verbe reste \textbf{toujours en 2\textsuperscript{e} position}.\par
\all{Letztes Wochenende war ein großes Fest.} « Le week-end dernier, il y a eu une grande fête. » (complément de temps en tête $\rightarrow$ inversion : verbe puis sujet)\par
\all{Ein großes Fest war letztes Wochenende.} (complément de temps en fin)\par
Erreur typique du francophone : \emph{Letztes Wochenende ein großes Fest war} — le verbe doit suivre immédiatement le premier élément.
\end{attention}

\courssub{Richtig oder Falsch ? Justifier par une citation}

L'exercice « vrai / faux » est un classique du baccalauréat. Il ne suffit pas de cocher : il faut prouver.

\begin{propriete}[Règle d'or du « Richtig / Falsch » au Bac]
Toute réponse \all{Richtig} ou \all{Falsch} doit être justifiée par une \textbf{citation exacte} du texte, recopiée mot pour mot \textbf{entre guillemets allemands} „…“. Une réponse sans citation perd la moitié des points, même si elle est correcte. Si la phrase est fausse, on cite le passage qui prouve le contraire.
\end{propriete}

\begin{exoresolu}[Richtig oder Falsch ? Begründe mit einem Zitat]
\textbf{Énoncé.}\par
a) Das Fest war in der Stadt Douala.\par
b) Der Großvater des Erzählers sprach über die alte Musik.\par
c) In Deutschland gibt es keine Konzerte.\par
d) Der Erzähler tanzte bis Mitternacht.
\tcblower
\textbf{Solution détaillée.}\par
a) \all{Falsch.} Beleg : \all{„Letztes Wochenende war ein großes Fest in meinem Dorf bei Yaoundé.“} — Le texte dit « village près de Yaoundé », pas Douala.\par
b) \all{Richtig.} Beleg : \all{„Mein Großvater erzählte mir, dass diese Musik sehr alt ist.“}\par
c) \all{Falsch.} Beleg : \all{„In den großen Städten wie Berlin, München oder Wien gibt es viele Konzerte.“}\par
d) \all{Richtig.} Beleg : \all{„Ich tanzte mit meinen Freunden bis Mitternacht.“}
\end{exoresolu}

\courssub{Recherche d'équivalents dans le texte}

Autre exercice fréquent : retrouver dans le texte le mot allemand qui correspond à un mot français donné. Stratégie : repérer d'abord la \emph{nature} du mot cherché (verbe ? nom ? adjectif ?), puis parcourir le texte en cherchant cette nature.

\begin{exemplebox}[Équivalents à trouver dans le texte]
\begin{center}
\begin{tabularx}{\linewidth}{|X|X|X|}
\hline
\rowcolor{popblueL} \textbf{Français} & \textbf{Mot du texte} & \textbf{Nature / remarque} \\
\hline
se produire (sur scène) & \all{auftreten} & verbe à particule séparable : \all{trat ... auf} \\
\hline
la foule & \all{die Menge} & nom féminin, singulier \\
\hline
célèbre & \all{berühmt} & adjectif : \all{ein berühmter Sänger} \\
\hline
\end{tabularx}
\end{center}
Dans le texte : \all{Ein berühmter Sänger aus Douala trat auch auf.} « Un chanteur célèbre de Douala s'est aussi produit. » Remarquez que \all{auftreten} est séparable : au Präteritum, \all{trat} se place en 2\textsuperscript{e} position et \all{auf} à la fin.
\end{exemplebox}

\begin{attention}[Faux amis et pièges de recherche]
\begin{itemize}
\item \all{die Menge} signifie « la foule, la quantité » ; ne pas la confondre avec le français « mélange » ni avec « la mangeoire ».
\item \all{das Konzert} = « le concert » (musique) ; pour « de concert avec » au sens d'accord, l'allemand dit \all{gemeinsam mit} ou \all{im Einvernehmen mit}.
\item \all{berühmt} = « célèbre » ; le français « fameux » au sens de « formidable » ne se traduit pas par \all{berühmt} : on dirait \all{toll} ou \all{großartig}.
\end{itemize}
\end{attention}

\courssub{Résumé guidé du texte (Zusammenfassung)}

Résumer, c'est reformuler l'essentiel en 3 à 4 phrases, avec ses propres mots, en reliant les idées par des connecteurs. Voici les connecteurs à mobiliser :

\begin{center}
\begin{tabularx}{\linewidth}{|X|X|X|}
\hline
\rowcolor{popblueL} \textbf{Connecteur} & \textbf{Français} & \textbf{Emploi} \\
\hline
\all{zuerst} & d'abord & première étape du récit \\
\hline
\all{dann} & puis, ensuite & enchaînement chronologique \\
\hline
\all{außerdem} & de plus, en outre & ajout d'une idée \\
\hline
\all{schließlich} & finalement, enfin & conclusion \\
\hline
\end{tabularx}
\end{center}

\begin{exemplebox}[Modèle de résumé en 4 phrases]
\all{Zuerst erzählt der Text von einem großen Fest in einem Dorf bei Yaoundé, wo die Leute zu Trommelmusik tanzten und sangen. Dann beschreibt der Erzähler die Freude der Menge und den Auftritt eines berühmten Sängers. Außerdem vergleicht er die Musik in Kamerun mit den Konzerten und Volksfesten in Deutschland. Schließlich sagt er, dass Musik die Menschen überall auf der Welt verbindet.}\par
« D'abord, le texte raconte une grande fête dans un village près de Yaoundé, où les gens dansaient et chantaient au son des tambours. Ensuite, le narrateur décrit la joie de la foule et la prestation d'un chanteur célèbre. De plus, il compare la musique au Cameroun avec les concerts et les fêtes populaires en Allemagne. Finalement, il dit que la musique unit les hommes partout dans le monde. »
\end{exemplebox}

\begin{attention}[Le résumé n'est pas une traduction]
Deux erreurs à éviter : recopier des phrases entières du texte (il faut \emph{reformuler}) et oublier l'inversion après un connecteur en tête de phrase : \all{Zuerst erzählt der Text...} et non \emph{Zuerst der Text erzählt...}. Le verbe reste en 2\textsuperscript{e} position. Autre piège : \all{außerdem} et \all{schließlich} occupent eux aussi la 1\textsuperscript{re} position, donc le verbe les suit immédiatement : \all{Außerdem vergleicht er...}, \all{Schließlich sagt er...}.
\end{attention}

\courssub{Expression orale guidée (Mündlicher Ausdruck)}

Pour terminer, on passe du texte à sa propre expérience. Question : \all{Welches Fest kennst du in deiner Region ?} « Quelle fête connais-tu dans ta région ? » Réponse attendue : 2 à 3 phrases complètes, au Präsens ou au Perfekt.

\begin{exemplebox}[Modèles de réponses orales]
\all{In meiner Region gibt es das Ngondo-Fest. Die Leute tanzen und singen am Wasser, und die Trommeln spielen den ganzen Tag.} « Dans ma région, il y a la fête du Ngondo. Les gens dansent et chantent au bord de l'eau, et les tambours jouent toute la journée. »\par\medskip
\all{Ich kenne ein großes Fest in meinem Dorf bei Bafoussam. Letztes Jahr habe ich dort mit meiner Familie getanzt, und ein berühmter Sänger ist aufgetreten.} « Je connais une grande fête dans mon village près de Bafoussam. L'année dernière, j'y ai dansé avec ma famille, et un chanteur célèbre s'est produit. » (Au Perfekt, \all{auftreten} se conjugue avec \all{sein} : \all{ist aufgetreten}.)
\end{exemplebox}

\begin{experience}[À toi de parler !]
Par groupes de deux, posez-vous la question \all{Welches Fest kennst du in deiner Region ?} Chaque élève répond en 2 ou 3 phrases en utilisant au moins trois mots du lexique de la leçon (\all{das Fest, die Trommel, die Menge, der Sänger, auftreten, tanzen, klatschen}) et un connecteur (\all{zuerst, dann, außerdem, schließlich}). Le partenaire pose ensuite une question de détail avec \all{wann}, \all{wo} ou \all{wie}.
\end{experience}

\begin{exercice}[Pour t'entraîner]
1. Réponds par des phrases complètes : a) Wann war das Fest ? b) Wer trat auf ? c) Was würde der Erzähler in Deutschland gern sehen ?\par
2. Richtig oder Falsch ? Justifie par une citation : a) Niemand wollte nach Hause gehen. b) Der Erzähler war schon in Deutschland.\par
3. Trouve dans le texte un mot qui signifie : « applaudir », « le souci », « unir ».\par
4. Résume le texte en 3 phrases avec \all{zuerst}, \all{außerdem} et \all{schließlich}.
\end{exercice}

\begin{corrige}[Exercice « Pour t'entraîner »]
1. a) \all{Das Fest war letztes Wochenende.} b) \all{Ein berühmter Sänger aus Douala trat auf.} c) \all{Der Erzähler würde gern ein Volksfest sehen, zum Beispiel das Oktoberfest in München.}\par
2. a) \all{Richtig.} Beleg : \all{„Alle waren glücklich, und niemand wollte nach Hause gehen.“} b) \all{Falsch.} Beleg : \all{„Wenn ich in Deutschland wäre, würde ich ein klassisches Konzert besuchen.“} — le subjonctif II montre qu'il n'y est jamais allé.\par
3. applaudir = \all{klatschen} ; le souci = \all{die Sorge} ; unir = \all{verbinden}.\par
4. Exemple : \all{Zuerst beschreibt der Erzähler ein Fest mit Trommeln und Tanz in seinem Dorf. Außerdem spricht er über die Musik und die Konzerte in Deutschland. Schließlich erklärt er, dass Musik alle Menschen verbindet.}
\end{corrige}

\begin{aretenir}[L'essentiel de la section]
\begin{itemize}
\item On lit deux fois : d'abord globalement (thème, narrateur, pays comparés), puis en détail (lieux, actions, réactions).
\item Toute réponse est une \textbf{phrase complète} qui reformule la question ; jamais un mot isolé.
\item \all{Richtig/Falsch} exige une \textbf{citation exacte} entre guillemets „…“.
\item Le verbe conjugué reste en 2\textsuperscript{e} position, même après un complément de temps ou un connecteur en tête de phrase.
\item Le résumé reformule avec des connecteurs : \all{zuerst, dann, außerdem, schließlich}.
\end{itemize}
\end{aretenir}

\coursec{Lexique thématique : die Musik und der Tanz}

Dans la section précédente, nous avons lu et compris le texte \all{Musik und Tanz in Kamerun und in Deutschland}. Nous y avons rencontré de nombreux mots nouveaux : \all{die Trommel}, \all{der Sänger}, \all{das Konzert}, \all{die Bühne}... Pour bien parler de la musique et de la danse, il ne suffit pas de comprendre ces mots une fois : il faut les \cle{mémoriser en réseau}, avec leur article, leur pluriel et leurs familles. C'est l'objet de cette section.

\courssub{Qu'est-ce qu'un champ lexical ?}

\begin{definition}[Le champ lexical]
Un \cle{champ lexical} est un ensemble de mots qui se rapportent à un même thème ou à une même réalité. Par exemple, \all{die Trommel}, \all{das Konzert}, \all{singen} et \all{der Sänger} appartiennent tous au champ lexical de la musique. Apprendre le vocabulaire \textbf{en réseau} (et non mot par mot, isolément) facilite la mémorisation : chaque mot nouveau s'accroche à des mots déjà connus, comme les branches d'un même arbre.
\end{definition}

\begin{methode}[Comment apprendre un mot allemand]
Pour chaque nom allemand, mémorise toujours \textbf{trois informations} :
\begin{enumerate}
\item le \cle{mot} lui-même : \all{Konzert} ;
\item son \cle{article} (le genre) : \all{das Konzert} ;
\item son \cle{pluriel} : \all{die Konzerte}.
\end{enumerate}
Un nom sans article est un mot à moitié appris : sans le genre, impossible de décliner correctement l'adjectif ou l'article (nous le verrons dans la section grammaire).
\end{methode}

\courssub{Un texte pour observer le vocabulaire en contexte}

Avant les tableaux, lisons ce court texte original. Relève tous les mots qui appartiennent au champ lexical de la musique et de la danse.

\begin{textebox}[Ein Fest in Yaoundé]
Letzten Samstag war in unserem Viertel in Yaoundé ein großes Fest. Schon am Nachmittag hatte sich eine große Menge vor der Bühne versammelt. Zuerst trat ein Chor aus Douala auf: zwölf Sänger und Sängerinnen sangen traditionelle Lieder, und zwei Musiker spielten dazu Trommeln. Der Gesang klang laut und schön, und das Publikum klatschte nach jedem Lied.

Dann kam ein bekannter Sänger aus Deutschland, der mit seiner Gitarre auf der Bühne stand. Er machte moderne Musik, aber auch rhythmische Stücke zum Tanzen. Die Tänzer und Tänzerinnen einer Gruppe aus Bafoussam zeigten traditionelle Tänze. Alle feierten, tanzten und amüsierten sich bis spät in die Nacht.

Meine Schwester spielt Flöte in einem Chor, und mein Bruder lernt Klavier. „Auf einem Volksfest in Deutschland“, sagt er, „gibt es oft ein Schlagzeug und eine Geige, aber keine Trommeln.“ Vielleicht treten wir nächstes Jahr beide bei einem Konzert auf. Musik und Tanz verbinden die Menschen — in Kamerun und in Deutschland.
\end{textebox}

\textbf{Glossaire} : das Viertel (die Viertel) : le quartier ; der Gesang (meist nur Singular) : le chant ; das Stück (-e) : le morceau ; verbinden : relier, unir ; vielleicht : peut-être.

\textbf{Questions de compréhension} :
\begin{enumerate}
\item \all{Wer trat zuerst auf?} (Qui est monté sur scène en premier ?)
\item \all{Was spielte der Sänger aus Deutschland?} (Que jouait le chanteur allemand ?)
\item \all{Welches Instrument lernt der Bruder?} (Quel instrument le frère apprend-il ?)
\end{enumerate}

\courssub{Champ lexical 1 : les instruments de musique}

\begin{center}
\begin{tabularx}{\linewidth}{|l|l|X|}
\hline
\rowcolor{popblueL} \textbf{Deutsch} & \textbf{Plural} & \textbf{Français} \\
\hline
die Trommel & die Trommeln & le tambour, le tam-tam \\
\hline
die Gitarre & die Gitarren & la guitare \\
\hline
das Klavier & die Klaviere & le piano \\
\hline
die Flöte & die Flöten & la flûte \\
\hline
die Geige & die Geigen & le violon \\
\hline
das Schlagzeug & die Schlagzeuge & la batterie \\
\hline
\end{tabularx}
\end{center}

\begin{exemplebox}[Les instruments en contexte]
\all{Meine Schwester spielt Flöte in einem Chor.} « Ma sœur joue de la flûte dans un chœur. »

\all{Zwei Musiker spielten Trommeln.} « Deux musiciens jouaient du tambour / des tambours. »

\all{Mein Bruder lernt Klavier.} « Mon frère apprend le piano. »
\end{exemplebox}

\begin{attention}[Jouer \emph{de} un instrument]
En allemand, on dit simplement \all{ein Instrument spielen}, sans article : \all{Ich spiele Gitarre.} « Je joue de la guitare. » Ne traduis pas le « de la » français : on ne dit pas \emph{*Ich spiele die Gitarre} pour parler de l'activité en général.
\end{attention}

\courssub{Champ lexical 2 : les personnes}

\begin{center}
\begin{tabularx}{\linewidth}{|l|l|X|}
\hline
\rowcolor{popgreenL} \textbf{Deutsch} & \textbf{Plural} & \textbf{Français} \\
\hline
der Sänger / die Sängerin & die Sänger / die Sängerinnen & le chanteur / la chanteuse \\
\hline
der Musiker & die Musiker & le musicien \\
\hline
der Tänzer / die Tänzerin & die Tänzer / die Tänzerinnen & le danseur / la danseuse \\
\hline
das Publikum & (die Publika, rare) & le public \\
\hline
die Menge & die Mengen & la foule \\
\hline
der Chor & die Chöre & le chœur, la chorale \\
\hline
\end{tabularx}
\end{center}

\begin{propriete}[Le féminin en \all{-in}]
En allemand, les noms de métiers et de fonctions forment leur féminin en ajoutant \cle{-in} au masculin ; le pluriel du féminin est en \cle{-innen} :
\begin{itemize}
\item \all{der Sänger} $\rightarrow$ \all{die Sängerin} $\rightarrow$ \all{die Sängerinnen}
\item \all{der Tänzer} $\rightarrow$ \all{die Tänzerin} $\rightarrow$ \all{die Tänzerinnen}
\end{itemize}
Attention au pluriel de \all{der Chor} : \all{die Chöre} (avec Umlaut).
\end{propriete}

\courssub{Champ lexical 3 : les événements et les lieux}

\begin{center}
\begin{tabularx}{\linewidth}{|l|l|X|}
\hline
\rowcolor{poporangeL} \textbf{Deutsch} & \textbf{Plural} & \textbf{Français} \\
\hline
das Konzert & die Konzerte & le concert \\
\hline
das Fest & die Feste & la fête \\
\hline
das Volksfest & die Volksfeste & la fête populaire \\
\hline
die Bühne & die Bühnen & la scène \\
\hline
die Aufführung & die Aufführungen & la représentation, le spectacle \\
\hline
der Ball & die Bälle & le bal \\
\hline
\end{tabularx}
\end{center}

\begin{attention}[Genres et pluriels piégeux]
\begin{itemize}
\item \all{das Konzert} est \textbf{neutre} : on ne dit jamais \emph{*die Konzert}, même si « concert » est masculin en français.
\item \all{das Fest} prend le pluriel \all{die Feste} : \all{Es gibt viele Feste in Kamerun.} « Il y a beaucoup de fêtes au Cameroun. »
\item \all{der Ball} (le bal) a pour pluriel \all{die Bälle} : ne pas confondre avec \all{der Ball} (le ballon), qui a le même pluriel mais un autre sens !
\end{itemize}
\end{attention}

\courssub{Les verbes utiles}

\begin{center}
\begin{tabularx}{\linewidth}{|X|l|X|}
\hline
\rowcolor{poppurpleL} \textbf{Infinitif} & \textbf{Präteritum / Perfekt} & \textbf{Français} \\
\hline
singen & sang / hat gesungen & chanter \\
\hline
tanzen & tanzte / hat getanzt & danser \\
\hline
spielen & spielte / hat gespielt & jouer \\
\hline
auftreten (tritt auf) & trat auf / ist aufgetreten & monter sur scène, se produire \\
\hline
klingen & klang / hat geklungen & sonner, retentir \\
\hline
feiern & feierte / hat gefeiert & fêter \\
\hline
sich amüsieren & amüsierte sich / hat sich amüsiert & s'amuser \\
\hline
applaudieren, klatschen & applaudierte, klatschte & applaudir \\
\hline
\end{tabularx}
\end{center}

\begin{propriete}[Points de conjugaison à retenir]
\begin{itemize}
\item \all{singen} est un \cle{verbe fort} : \all{ich singe}, \all{ich sang}, \all{ich habe gesungen}.
\item \all{auftreten} est un \cle{verbe à particule séparable} et fort : au Präsens, \all{der Sänger tritt auf} (la particule \all{auf} va en fin de phrase) ; au Perfekt, il prend l'auxiliaire \all{sein} : \all{Der Sänger ist aufgetreten.}
\item \all{sich amüsieren} est \cle{pronominal} : \all{ich amüsiere mich}, \all{wir haben uns amüsiert}.
\end{itemize}
\end{propriete}

\begin{exemplebox}[Exemples traduits]
\all{Der bekannte Sänger trat auf der Bühne auf.} « Le chanteur célèbre est monté sur scène. »

\all{Wir haben auf dem Fest getanzt und uns amüsiert.} « Nous avons dansé à la fête et nous nous sommes amusés. »

\all{Das Publikum klatschte nach jedem Lied.} « Le public applaudissait après chaque chanson. »

\all{Der Gesang klang laut und schön.} « Le chant sonnait fort et beau. »
\end{exemplebox}

\courssub{Adjectifs et expressions figées}

\begin{center}
\begin{tabularx}{\linewidth}{|X|X|}
\hline
\rowcolor{popgoldL} \textbf{Deutsch} & \textbf{Français} \\
\hline
laut / leise & fort (bruyant) / doux (silencieux) \\
\hline
bekannt / berühmt & connu / célèbre \\
\hline
rhythmisch & rythmé \\
\hline
traditionell / modern & traditionnel / moderne \\
\hline
Musik machen & faire de la musique \\
\hline
auf der Bühne stehen & être sur scène \\
\hline
\end{tabularx}
\end{center}

\begin{exemplebox}
\all{Die Tänzerinnen zeigten traditionelle Tänze.} « Les danseuses ont présenté des danses traditionnelles. »

\all{Er machte moderne, rhythmische Musik.} « Il faisait de la musique moderne et rythmée. »
\end{exemplebox}

\courssub{Familles de mots : apprendre en réseau}

\begin{propriete}[La famille de mots]
Une \cle{famille de mots} regroupe les mots construits sur une même racine. Connaître la racine permet de deviner le sens des mots dérivés :
\begin{itemize}
\item \all{der Tanz} (la danse) $\rightarrow$ \all{tanzen} (danser) $\rightarrow$ \all{der Tänzer} (le danseur) $\rightarrow$ \all{die Tänzerin} (la danseuse)
\item \all{die Musik} (la musique) $\rightarrow$ \all{der Musiker} (le musicien) $\rightarrow$ \all{musikalisch} (musical)
\item \all{der Gesang} (le chant) $\rightarrow$ \all{singen} (chanter) $\rightarrow$ \all{der Sänger} (le chanteur) $\rightarrow$ \all{die Sängerin} (la chanteuse)
\end{itemize}
\end{propriete}

\begin{popfigure}
\begin{center}\tcbox[colback=popblue,colframe=popblue,coltext=white,arc=9pt,boxsep=4pt,left=6pt,right=6pt]{\bfseries\small singen \footnotesize chanter}\end{center}
\vspace{-2pt}
\begin{tcbraster}[raster columns=2,raster equal height=rows,raster column skip=6pt,raster row skip=6pt,enhanced,blankest]
\begin{tcolorbox}[enhanced,colback=popblue,colframe=popblue,coltext=white,boxrule=0.9pt,arc=3pt,left=4pt,right=4pt,top=3pt,bottom=3pt,boxsep=1pt,halign=left]\bfseries\scriptsize der Sänger \footnotesize le chanteur\end{tcolorbox}
\begin{tcolorbox}[enhanced,colback=poporange,colframe=poporange,coltext=white,boxrule=0.9pt,arc=3pt,left=4pt,right=4pt,top=3pt,bottom=3pt,boxsep=1pt,halign=left]\bfseries\scriptsize die Sängerin \footnotesize la chanteuse\end{tcolorbox}
\begin{tcolorbox}[enhanced,colback=popgreen,colframe=popgreen,coltext=white,boxrule=0.9pt,arc=3pt,left=4pt,right=4pt,top=3pt,bottom=3pt,boxsep=1pt,halign=left]\bfseries\scriptsize der Gesang \footnotesize le chant\end{tcolorbox}
\begin{tcolorbox}[enhanced,colback=poppurple,colframe=poppurple,coltext=white,boxrule=0.9pt,arc=3pt,left=4pt,right=4pt,top=3pt,bottom=3pt,boxsep=1pt,halign=left]\bfseries\scriptsize das Lied \footnotesize la chanson\end{tcolorbox}
\end{tcbraster}

\legende{Arbre de la famille de mots autour de la racine \all{singen}}
\end{popfigure}

\begin{savaistu}[Das Lied et le « lied » français]
Le mot allemand \all{das Lied} (pluriel \all{die Lieder}), « la chanson », est entré dans le vocabulaire musical français sous la forme « lied » : il désigne un chant pour voix solo avec accompagnement de piano, genre cultivé par les compositeurs autrichiens et allemands du XIX\textsuperscript{e} siècle. Quand ton professeur de musique prononce « lied », il parle... allemand sans le savoir !
\end{savaistu}

\courssub{Exercice résolu et entraînement}

\begin{exoresolu}[Complète avec le bon article et la bonne forme]
Énoncé : complète les phrases avec le mot entre parenthèses, à la forme correcte (article défini + nom, éventuellement au pluriel).
\begin{enumerate}
\item \all{\ldots{} (Sänger) trat auf \ldots{} (Bühne) auf.}
\item \all{\ldots{} (Publikum) klatschte nach jedem Lied.}
\item \all{Wir haben auf \ldots{} (Fest) getanzt.}
\item \all{Meine Schwester spielt in \ldots{} (Chor).}
\end{enumerate}
\tcblower
Solution détaillée :
\begin{enumerate}
\item \all{\textbf{Der} Sänger trat auf \textbf{der} Bühne auf.} — \all{der Sänger} est masculin (sujet, nominatif) ; \all{auf der Bühne} : datif féminin après la préposition \all{auf} (position, pas de mouvement).
\item \all{\textbf{Das} Publikum klatschte nach jedem Lied.} — \all{das Publikum} est neutre ; sujet au nominatif.
\item \all{Wir haben auf \textbf{dem} Fest getanzt.} — \all{das Fest} est neutre ; après \all{auf} (position), datif : \all{dem}.
\item \all{Meine Schwester spielt in \textbf{einem} Chor.} — \all{der Chor} est masculin ; après \all{in} (position), datif : \all{einem}.
\end{enumerate}
\end{exoresolu}

\begin{attention}[Pièges fréquents des francophones]
\begin{itemize}
\item Ne dis pas \emph{*die Konzert} : c'est \all{das Konzert}. Le genre français ne prédit pas le genre allemand.
\item \all{auftreten} se sépare : \all{Der Chor tritt auf} (et non \emph{*Der Chor auftritt}).
\item « S'amuser » se dit \all{sich amüsieren} avec le pronom réfléchi : \all{wir haben \textbf{uns} amüsiert}.
\item Majuscule obligatoire à tous les noms : \all{das Fest}, \all{die Trommel}, \all{der Ball}.
\end{itemize}
\end{attention}

\begin{exercice}[À toi de jouer]
\begin{enumerate}
\item Traduis en allemand : « La chanteuse célèbre a chanté une chanson traditionnelle, et la foule a applaudi. »
\item Forme le féminin et les pluriels : \all{der Musiker} (féminin ?), \all{der Tänzer} (féminin ?), \all{das Konzert} (pluriel ?), \all{der Ball} (pluriel ?).
\item Conjugue au Perfekt : \all{wir (sich amüsieren)}, \all{der Sänger (auftreten)}, \all{ich (singen)}.
\end{enumerate}
\end{exercice}

\begin{corrige}[Exercice « À toi de jouer »]
\begin{enumerate}
\item \all{Die berühmte Sängerin hat ein traditionelles Lied gesungen, und die Menge hat geklatscht.}
\item \all{die Musikerin} ; \all{die Tänzerin} ; \all{die Konzerte} ; \all{die Bälle}.
\item \all{wir haben uns amüsiert} ; \all{der Sänger ist aufgetreten} ; \all{ich habe gesungen}.
\end{enumerate}
\end{corrige}

\begin{aretenir}[L'essentiel du lexique]
\begin{itemize}
\item Apprends chaque nom avec \textbf{article + pluriel} : \all{das Konzert, die Konzerte} ; \all{die Bühne, die Bühnen} ; \all{der Ball, die Bälle}.
\item Le féminin des personnes se forme en \all{-in} : \all{der Sänger} $\rightarrow$ \all{die Sängerin}.
\item Verbes clés : \all{singen} (sang, gesungen), \all{auftreten} (tritt auf, ist aufgetreten), \all{sich amüsieren}.
\item Expressions figées : \all{Musik machen}, \all{auf der Bühne stehen}.
\item Pense en \textbf{familles de mots} : \all{der Tanz} $\rightarrow$ \all{tanzen} $\rightarrow$ \all{der Tänzer} $\rightarrow$ \all{die Tänzerin}.
\end{itemize}
\end{aretenir}

Ce vocabulaire désormais acquis, nous pourrons l'employer dans des phrases plus riches : la prochaine section nous apprendra à \textbf{comparer} avec \all{wie} et \all{als ob} — par exemple : \all{Der Sänger tanzt wie ein Profi.} « Le chanteur danse comme un professionnel. »

\coursec{Grammaire 1 : la comparaison avec « wie » et « als ob »}

Dans la section précédente, nous avons enrichi notre vocabulaire sur le thème de la musique et de la danse : \all{die Trommel}, \all{der Sänger}, \all{das Konzert}, \all{das Fest}. Nous allons maintenant observer comment ces mots s'organisent dans des phrases qui \cle{comparent} : le musicien compare les sons, le journaliste compare les foules, le poète compare les émotions. En allemand, comparer obéit à des règles précises qui distinguent le \cle{réel} de l'\cle{imaginaire}. C'est toute la différence entre \all{wie} et \all{als ob}.

\courssub{Observation : deux phrases du texte support}

Reprenons deux phrases tirées du texte \emph{Musik und Tanz in Kamerun und in Deutschland} :

\begin{exemplebox}[Phrases à observer]
\all{Musik verbindet die Menschen wie eine universelle Sprache.}\\
« La musique relie les hommes comme une langue universelle. »

\medskip
\all{Die Menge jubelte, als ob sie einen König sähe.}\\
« La foule jubilait, comme si elle voyait un roi. »
\end{exemplebox}

Observons attentivement :

\begin{itemize}
\item Dans la première phrase, la comparaison est \cle{réelle} : la musique fonctionne effectivement \emph{comme} une langue universelle. On utilise simplement \all{wie} + nom, sans verbe répété.
\item Dans la deuxième phrase, la comparaison est \cle{irréelle} : la foule ne voyait pas vraiment un roi ; c'est une image, une exagération. On utilise \all{als ob} et le verbe \all{sähe} n'est pas à l'indicatif (\all{sah}) mais au \cle{Konjunktiv II}, placé \emph{à la fin} de la subordonnée.
\end{itemize}

Toute la leçon découle de cette observation : l'allemand marque grammaticalement la différence entre « comme » (vrai) et « comme si » (imaginaire), là où le français se contente souvent du même mot.

\courssub{La comparaison réelle : « wie », « so ... wie », comparatif + « als »}

\begin{propriete}[Les trois constructions de la comparaison réelle]
\begin{enumerate}
\item \cle{Égalité} : \all{so + adjectif + wie} = « aussi ... que ».\\
\all{Die Trommel ist so laut wie ein Motor.} « Le tambour est aussi bruyant qu'un moteur. »
\item \cle{Supériorité / infériorité} : \all{comparatif + als} = « plus ... que ».\\
\all{Das Konzert in Douala war größer als das Konzert in Bafoussam.} « Le concert de Douala était plus grand que celui de Bafoussam. »
\item \cle{Comparaison avec un nom} (sans répéter le verbe) : \all{wie} + nom.\\
\all{Er singt wie ein Vogel.} « Il chante comme un oiseau. »
\end{enumerate}
\end{propriete}

\begin{attention}[« als » ou « wie » ? Le piège capital]
Les francophones confondent souvent \all{als} et \all{wie} car le français utilise « que » dans les deux cas. Retenez :
\begin{itemize}
\item \all{wie} sert à l'\textbf{égalité} : \all{so groß wie} « aussi grand que » ;
\item \all{als} sert à la \textbf{différence} : \all{größer als} « plus grand que ».
\end{itemize}
On ne dit jamais *\all{so groß als} ni *\all{größer wie}. Astuce : \all{so} appelle \all{wie} ; le comparatif en \all{-er} appelle \all{als}.
\end{attention}

Rappelons aussi que \all{als} a un autre emploi à ne pas confondre : la subordonnée temporelle au passé ponctuel (\all{Als ich Kind war, tanzte ich gern.} « Quand j'étais enfant, j'aimais danser. »). Le contexte permet toujours de distinguer les deux emplois.

\begin{exemplebox}[Comparaisons réelles sur le thème de la musique]
\all{Der Sänger ist so berühmt wie ein Fußballstar.} « Le chanteur est aussi célèbre qu'une star du football. »

\medskip
\all{Die Bendskin-Fahrer hupen lauter als die Trommler.} « Les conducteurs de bendskins klaxonnent plus fort que les joueurs de tambour. »

\medskip
\all{Sie tanzt wie ihre Großmutter.} « Elle danse comme sa grand-mère. »

\medskip
\all{Das Fest war so schön wie letztes Jahr.} « La fête était aussi belle que l'année dernière. »
\end{exemplebox}

\courssub{La comparaison irréelle : « als ob » + Konjunktiv II}

\begin{propriete}[Règle fondamentale]
Après \all{als ob} (« comme si »), le verbe conjugué est \textbf{toujours au Konjunktiv II} et se place \textbf{en fin de proposition} (subordonnée). L'indicatif est interdit.\\
Structure : phrase principale + \all{als ob} + sujet + ... + verbe au KII en finale.\\
\all{Er tanzt, als ob er ein Profi wäre.} « Il danse comme s'il était un professionnel. »
\end{propriete}

Le Konjunktiv II marque ici l'\cle{irréel} : la comparaison est imaginaire, exagérée ou hypothétique. Le français fait le même choix avec l'imparfait ou le plus-que-parfait après « comme si » : « comme s'il \emph{était} », « comme si elle n'avait jamais \emph{eu} peur ».

\textbf{Le Konjunktiv II présent} sert quand l'irréel est simultané (même moment que la principale) :

\begin{center}
\begin{tabular}{|l|l|l|}
\hline
\rowcolor{popblueL} Infinitif & KII présent & Exemple \\
\hline
sein & wäre & als ob er müde wäre \\
\hline
haben & hätte & als ob sie Zeit hätte \\
\hline
werden & würde & als ob es regnen würde \\
\hline
können & könnte & als ob er fliegen könnte \\
\hline
tun & täte & als ob sie nichts täte \\
\hline
sehen (fort) & sähe & als ob sie einen König sähe \\
\hline
\end{tabular}
\end{center}

\textbf{Le Konjunktiv II passé} sert quand l'action imaginée est antérieure : \all{hätte/wäre + Partizip II} (auxiliaire au KII en finale, le participe juste avant) :

\all{Sie sang, als ob sie nie Angst gehabt hätte.} « Elle chantait comme si elle n'avait jamais eu peur. »

\all{Er kam herein, als ob er den ganzen Tag getanzt hätte.} « Il entra comme s'il avait dansé toute la journée. »

\all{Sie lächelte, als ob sie schon gekommen wäre.} « Elle souriait comme si elle était déjà venue. » (verbe de mouvement : auxiliaire \all{sein})

\begin{exemplebox}[Exemples traduits]
\all{Er tanzt, als ob er ein Profi wäre.} « Il danse comme s'il était un professionnel. »

\medskip
\all{Sie sang, als ob sie nie Angst gehabt hätte.} « Elle chantait comme si elle n'avait jamais eu peur. »

\medskip
\all{Der Trommler schlug die Trommel, als ob er sie bezwingen müsste.} « Le joueur de tambour frappait la trommel comme s'il devait la dompter. »

\medskip
\all{Das Publikum schrie, als ob das Stadion brennen würde.} « Le public criait comme si le stade brûlait. »
\end{exemplebox}

\courssub{La forme contractée : « als » + KII avec inversion}

L'allemand permet de supprimer \all{ob} ; dans ce cas, le verbe au Konjunktiv II se place \textbf{juste après} \all{als} (inversion), exactement comme dans une subordonnée conditionnelle sans \all{wenn} :

\begin{exemplebox}[Transformation]
\all{Die Menge jubelte, als ob sie einen König sähe.}\\
$\rightarrow$ \all{Die Menge jubelte, als sähe sie einen König.}\\
« La foule jubilait, comme si elle voyait un roi. »
\end{exemplebox}

\begin{propriete}[Deux constructions équivalentes]
\begin{enumerate}
\item \all{als ob + sujet + ... + verbe KII en finale} : \all{als ob sie einen König sähe}
\item \all{als + verbe KII + sujet + ...} (sans \all{ob}) : \all{als sähe sie einen König}
\end{enumerate}
Les deux formes ont exactement le même sens. La forme contractée est un peu plus littéraire ; la forme avec \all{ob} est plus fréquente à l'oral.
\end{propriete}

\courssub{Schéma récapitulatif : réel ou irréel ?}

\begin{popfigure}
\begin{tikzpicture}[grow=right, level distance=4.6cm,
level 1/.style={sibling distance=3.2cm},
level 2/.style={sibling distance=1.6cm, level distance=4.2cm},
every node/.style={align=center, font=\small},
edge from parent/.style={draw=popmuted, thick}]
\node[draw=popdark, fill=popblueL, rounded corners, font=\bfseries] {Comparaison}
child { node[draw=poporange, fill=poporangeL, rounded corners, align=center]{Irréelle\\« comme si »}
child { node[draw=popmuted, fill=popcream, rounded corners, align=center]{als + KII\\(inversion, sans ob)} }
child { node[draw=popmuted, fill=popcream, rounded corners, align=center]{als ob + KII\\verbe en finale} }
}
child { node[draw=popgreen, fill=popgreenL, rounded corners, align=center]{Réelle\\« comme / que »}
child { node[draw=popmuted, fill=popcream, rounded corners, align=center]{wie + nom\\Er singt wie ein Vogel} }
child { node[draw=popmuted, fill=popcream, rounded corners, align=center]{so ... wie\\(égalité)} }
child { node[draw=popmuted, fill=popcream, rounded corners, align=center]{comparatif + als\\größer als} }
};
\end{tikzpicture}
\legende{Les deux familles de la comparaison en allemand : réelle (indicatif) et irréelle (Konjunktiv II).}
\end{popfigure}

\begin{methode}[Comment choisir entre « wie » et « als ob »]
\begin{enumerate}
\item Je me demande : la comparaison est-elle \textbf{vraie ou plausible} ? (Il chante vraiment comme un oiseau, le concert est vraiment plus grand.) $\rightarrow$ j'utilise \all{wie}, \all{so ... wie} ou \all{comparatif + als}, à l'indicatif.
\item La comparaison est-elle \textbf{imaginaire, exagérée ou fausse} ? (La foule n'a pas vu de roi, il n'est pas vraiment un professionnel.) $\rightarrow$ j'utilise \all{als ob} + Konjunktiv II.
\item Je choisis le temps du KII : simultané $\rightarrow$ KII présent (\all{wäre, hätte, sähe}) ; antérieur $\rightarrow$ KII passé (\all{hätte gesehen, wäre gekommen}).
\item Je vérifie la place du verbe : en finale après \all{als ob}, ou juste après \all{als} si je supprime \all{ob}.
\end{enumerate}
\end{methode}

\courssub{Comparaison français / allemand}

\begin{center}
\begin{tabularx}{\linewidth}{|X|X|X|}
\hline
\rowcolor{popblueL} Français & Deutsch & Remarque \\
\hline
Il chante comme un oiseau. & Er singt wie ein Vogel. & réel : wie + nom \\
\hline
aussi beau que l'an dernier & so schön wie letztes Jahr & égalité : so ... wie \\
\hline
plus grand que Yaoundé & größer als Jaunde & différence : comparatif + als \\
\hline
comme s'il était roi & als ob er König wäre & irréel : KII présent \\
\hline
comme si elle avait gagné & als ob sie gewonnen hätte & irréel : KII passé \\
\hline
comme s'il pleuvait & als ob es regnen würde & würde + infinitif (verbes faibles) \\
\hline
\end{tabularx}
\end{center}

\begin{attention}[L'erreur n° 1 du francophone]
Ne jamais mettre l'indicatif après \all{als ob} ! La phrase *\all{Die Menge jubelte, als ob sie einen König sah.} est fautive : il faut le Konjunktiv II \all{sähe}. De même : *\all{als ob er ist} $\rightarrow$ \all{als ob er wäre}. Le français « comme si » + imparfait doit vous rappeler que l'allemand exige le KII. Autre piège : ne pas oublier la place du verbe en fin de subordonnée : *\all{als ob sie sähe einen König} est impossible.
\end{attention}

\courssub{Texte d'application}

\begin{textebox}[Das Schulfest in Yaoundé]
Am Samstag feierte unsere Schule ein großes Fest. Der Schulhof war so voll wie ein Stadion vor einem Fußballspiel. Die Trommler spielten lauter als die Lautsprecher, und die Menge bewegte sich wie ein einziger Körper. Ein Junge aus der Terminale tanzte, als ob er ein Profi wäre; alle klatschten, als ob sie einen Weltmeister sähen. Eine Lehrerin sang so schön wie eine Opernsängerin, und die Schüler hörten zu, als ob sie nie etwas Schöneres gehört hätten. Am Abend leuchteten die Laternen, als wäre der Himmel auf die Erde gefallen. Es war, als ob die ganze Stadt mit uns feierte. Musik verbindet die Menschen wirklich wie eine universelle Sprache.
\end{textebox}

\textbf{Glossaire} : \all{der Schulhof, -höfe} la cour de l'école ; \all{der Lautsprecher, -} le haut-parleur ; \all{der Körper, -} le corps ; \all{der Weltmeister, -} le champion du monde ; \all{die Opernsängerin, -nen} la chanteuse d'opéra ; \all{die Laterne, -n} la lanterne ; \all{leuchten} briller.

\textbf{Questions} : 1) Relevez toutes les comparaisons du texte et classez-les : réelles ou irréelles ? 2) Justifiez l'emploi du Konjunktiv II dans \all{als ob er ein Profi wäre}. 3) Transformez \all{als wäre der Himmel auf die Erde gefallen} en construction avec \all{als ob}.

\courssub{Exercice résolu et entraînement}

\begin{exoresolu}[Complétez avec la forme correcte]
1. Der Sänger ist so berühmt \_\_\_ ein Filmstar. (aussi ... que)\\
2. Das Fest war größer \_\_\_ letztes Jahr. (plus ... que)\\
3. Er spielt Trommel, als ob er es schon zwanzig Jahre \_\_\_ (tun, KII passé).\\
4. Sie lächelte, als ob sie das Konzert gewonnen \_\_\_ (haben, KII passé).\\
5. Das Mädchen tanzt, als ob es schweben \_\_\_ (können, KII présent).
\tcblower
\textbf{Solution détaillée :}\\
1. \all{wie} : égalité, construction \all{so ... wie}. « Le chanteur est aussi célèbre qu'une star de cinéma. »\\
2. \all{als} : comparatif de supériorité \all{größer als}. « La fête était plus grande que l'année dernière. »\\
3. \all{getan hätte} : comparaison irréelle au passé, auxiliaire \all{haben} au KII en finale. « ... comme s'il en jouait depuis vingt ans. »\\
4. \all{hätte} : \all{als ob sie das Konzert gewonnen hätte} « comme si elle avait gagné le concert » ; le participe \all{gewonnen} précède l'auxiliaire.\\
5. \all{könnte} : KII présent de \all{können} pour une action simultanée imaginaire. « ... comme si elle pouvait voler. »
\end{exoresolu}

\begin{exercice}[À vous de jouer]
Traduisez en allemand : 1) Le tambour est aussi vieux que le village. 2) Elle chante comme si elle était seule au monde. 3) Le public dansait comme s'il n'avait jamais entendu cette musique. 4) Transformez : \all{Er lief, als ob er den Bus verpassen würde.} en forme contractée sans \all{ob}.
\end{exercice}

\begin{corrige}[Exercice « À vous de jouer »]
1. \all{Die Trommel ist so alt wie das Dorf.} 2. \all{Sie singt, als ob sie allein auf der Welt wäre.} 3. \all{Das Publikum tanzte, als ob es diese Musik nie gehört hätte.} 4. \all{Er lief, als würde er den Bus verpassen.}
\end{corrige}

\begin{aretenir}[L'essentiel]
Comparaison \textbf{réelle} : \all{wie} + nom, \all{so ... wie} (égalité), \all{comparatif + als} (différence), toujours à l'indicatif. Comparaison \textbf{irréelle} : \all{als ob} + Konjunktiv II en fin de proposition (\all{wäre, hätte, sähe, könnte} ; passé : \all{hätte/wäre + Partizip II}), ou forme contractée \all{als} + KII avec inversion. Jamais d'indicatif après \all{als ob} !
\end{aretenir}

\coursec{Grammaire 2 : la déclinaison de l'adjectif et le subjonctif II}

Dans la section précédente, nous avons appris à comparer avec \all{wie} et \all{als ob} : \all{Er tanzt wie ein Profi} (« Il danse comme un professionnel »), \all{Sie singt, als ob sie berühmt wäre} (« Elle chante comme si elle était célèbre »). Deux problèmes grammaticaux se cachaient dans ces phrases : comment accorder l'adjectif allemand (\all{ein berühmter Sänger} ou \all{der berühmte Sänger} ?), et quelle est cette forme \all{wäre} qui traduit le conditionnel français ? C'est l'objet de cette section : la \cle{déclinaison de l'adjectif} et le \cle{subjonctif II} (Konjunktiv II), deux piliers de l'expression écrite au baccalauréat.

\courssub{La déclinaison de l'adjectif épithète}

\textbf{Observation.} Lisons d'abord ces phrases tirées de notre thème :

\begin{exemplebox}[Observer avant de généraliser]
\all{Der bekannte Sänger tritt heute in Yaoundé auf.} « Le chanteur connu se produit aujourd'hui à Yaoundé. »\par
\all{Ich habe ein bekanntes Lied gehört.} « J'ai entendu une chanson connue. »\par
\all{Das Konzert endete mit lautem Applaus.} « Le concert s'est terminé par des applaudissements nourris. »\par
\all{Gute Musik verbindet die Menschen.} « La bonne musique relie les hommes. »
\end{exemplebox}

On constate que l'adjectif épithète (placé devant le nom) prend une \cle{désinence} : \all{bekannte}, \all{bekannten}, \all{bekannter}, \all{bekanntes}, \all{lautem}... Le choix de cette désinence dépend de trois facteurs : le \cle{genre} du nom, le \cle{cas}, et la présence ou non d'un \cle{article} devant l'adjectif. Il existe donc \textbf{trois tableaux} de déclinaison.

\begin{propriete}[Règle fondamentale : qui porte l'information ?]
En allemand, la marque du genre et du cas doit apparaître \textbf{au moins une fois} dans le groupe nominal :
\begin{itemize}
\item si l'article est \textbf{défini} (\all{der, die, das}), il porte déjà l'information : l'adjectif se contente des désinences « pauvres » \cle{-e} ou \cle{-en} ;
\item si l'article est \textbf{indéfini} (\all{ein, eine}) ou \textbf{absent}, l'adjectif doit porter lui-même la marque du genre et du cas : désinences « riches » \cle{-er, -es, -e, -en, -em}.
\end{itemize}
\end{propriete}

\textbf{Tableau 1 : après l'article défini (der, die, das, die).}

\begin{center}
\begin{tabular}{|l|l|l|l|l|}
\hline
\rowcolor{popblueL} Cas & Masculin & Féminin & Neutre & Pluriel \\
\hline
Nominatif & der bekannte Sänger & die bekannte Sängerin & das bekannte Lied & die bekannten Lieder \\
\hline
Accusatif & den bekannten Sänger & die bekannte Sängerin & das bekannte Lied & die bekannten Lieder \\
\hline
Datif & dem bekannten Sänger & der bekannten Sängerin & dem bekannten Lied & den bekannten Liedern \\
\hline
Génitif & des bekannten Sängers & der bekannten Sängerin & des bekannten Liedes & der bekannten Lieder \\
\hline
\end{tabular}
\end{center}

\begin{propriete}[Après der/die/das : -e ou -en]
Après article défini, l'adjectif prend \cle{-e} au nominatif singulier (trois genres) et à l'accusatif féminin et neutre ; il prend \cle{-en} dans \textbf{tous les autres cas}, y compris tout le pluriel. Astuce : \textbf{5 fois -e, tout le reste -en}. La même règle vaut après \all{dieser, jener, jeder, welcher, alle}.
\end{propriete}

\begin{exemplebox}[Un groupe nominal décliné aux 4 cas]
\all{Der laute Applaus gefällt dem Sänger.} « Les applaudissements nourris plaisent au chanteur. » (nominatif)\par
\all{Wegen des lauten Applauses kam er zurück.} « À cause des applaudissements nourris, il est revenu. » (génitif)\par
\all{Er dankte dem Publikum für den lauten Applaus.} « Il a remercié le public pour les applaudissements nourris. » (accusatif)\par
\all{Mit dem lauten Applaus endete das Fest.} « La fête s'est terminée par des applaudissements nourris. » (datif)
\end{exemplebox}

\textbf{Tableau 2 : après l'article indéfini (ein, eine, ein) et les possessifs (mein, dein, kein).}

\begin{center}
\begin{tabular}{|l|l|l|l|}
\hline
\rowcolor{popgreenL} Cas & Masculin & Féminin & Neutre \\
\hline
Nominatif & ein bekannter Sänger & eine bekannte Sängerin & ein bekanntes Lied \\
\hline
Accusatif & einen bekannten Sänger & eine bekannte Sängerin & ein bekanntes Lied \\
\hline
Datif & einem bekannten Sänger & einer bekannten Sängerin & einem bekannten Lied \\
\hline
Génitif & eines bekannten Sängers & einer bekannten Sängerin & eines bekannten Liedes \\
\hline
\end{tabular}
\end{center}

Ici, l'article \all{ein} ne marque ni le genre au nominatif (\all{ein} peut être masculin ou neutre) ni le cas à l'accusatif neutre : l'adjectif prend alors la marque forte \cle{-er} ou \cle{-es} : \all{ein bekannter Sänger}, \all{ein bekanntes Lied}. Partout où \all{ein} est déjà décliné (\all{einen, einem, einer, eines}), l'adjectif prend \cle{-en}.

\textbf{Tableau 3 : sans article (déclinaison forte).}

\begin{center}
\begin{tabular}{|l|l|l|l|l|}
\hline
\rowcolor{poporangeL} Cas & Masculin & Féminin & Neutre & Pluriel \\
\hline
Nominatif & lauter Applaus & gute Musik & lautes Fest & gute Lieder \\
\hline
Accusatif & lauten Applaus & gute Musik & lautes Fest & gute Lieder \\
\hline
Datif & lautem Applaus & guter Musik & lautem Fest & guten Liedern \\
\hline
Génitif & lauten Applauses & guter Musik & lauten Festes & guter Lieder \\
\hline
\end{tabular}
\end{center}

Sans article, l'adjectif porte seul toute l'information : il prend les désinences de \all{der/die/das} (sauf au génitif masculin et neutre : \cle{-en}).

\begin{attention}[Sans article, l'adjectif porte la marque du cas]
Ne dites jamais *\all{mit lauter Applaus} : après la préposition \all{mit} (datif), il faut \all{mit lautem Applaus}. Comparez : \all{lauter Applaus} (nominatif masculin) mais \all{mit lautem Applaus} (datif). De même : \all{gute Musik} (nominatif) mais \all{wegen guter Musik} (génitif). En français, l'adjectif ne change pas ; en allemand, il change obligatoirement.
\end{attention}

\begin{popfigure}
\begin{tikzpicture}
\node[draw=popdark, fill=popblueL, rounded corners, font=\bfseries, align=center] (titre) at (0,0) {Désinences de l'adjectif épithète};
\node[draw=popblue, fill=popblue!15, rounded corners, align=center] (a) at (-5.2,-2.2) {\textbf{nach bestimmtem Artikel}\\ der, die, das\\[2pt] -e (5 formes)\\ -en (partout ailleurs)\\ \all{der bekannte Sänger}};
\node[draw=popgreen, fill=popgreenL, rounded corners, align=center] (b) at (0,-2.2) {\textbf{nach unbestimmtem Artikel}\\ ein, mein, kein\\[2pt] -er / -es si ein est\\ indécliné, sinon -en\\ \all{ein bekannter Sänger}};
\node[draw=poporange, fill=poporangeL, rounded corners, align=center] (c) at (5.2,-2.2) {\textbf{ohne Artikel}\\ (déclinaison forte)\\[2pt] -er, -e, -es, -en, -em\\ comme der/die/das\\ \all{lauter Applaus}};
\draw[-{Stealth}, thick, popblue] (titre) -- (a);
\draw[-{Stealth}, thick, popgreen] (titre) -- (b);
\draw[-{Stealth}, thick, poporange] (titre) -- (c);
\end{tikzpicture}
\legende{Les trois déclinaisons de l'adjectif allemand : tout dépend de l'article qui précède.}
\end{popfigure}

\begin{exoresolu}[Décliner l'adjectif]
Énoncé : complétez avec la bonne désinence. 1) \all{Der berühmte Dirigent leitet das berühmt\_\_ Orchester.} 2) \all{Sie kauft ein neu\_\_ Trommel.} 3) \all{Wir tanzen zu traditionell\_\_ Musik.}
\tcblower
Solution détaillée : 1) \all{das berühmte Orchester} : article défini + accusatif neutre $\rightarrow$ \cle{-e}. 2) \all{eine neue Trommel} : \all{die Trommel} est féminin ; accusatif féminin après \all{ein} $\rightarrow$ \cle{-e} (et \all{ein} devient \all{eine}). 3) \all{zu traditioneller Musik} : \all{zu} exige le datif ; sans article, l'adjectif porte la marque du datif féminin $\rightarrow$ \cle{-er}.
\end{exoresolu}

\courssub{Le subjonctif II (Konjunktiv II) : formation}

Revenons à notre phrase d'observation : \all{Sie singt, als ob sie berühmt wäre.} La forme \all{wäre} est le \cle{Konjunktiv II} (KII) de \all{sein}. Il correspond très souvent à notre \textbf{conditionnel français} : « elle serait ».

\begin{propriete}[Formation du KII présent]
Le KII présent se forme sur le \textbf{radical du prétérit} + les désinences \cle{-e, -est, -e, -en, -et, -en}. Pour les verbes forts dont le prétérit contient \textbf{a, o, u}, on ajoute un \cle{tréma} : \all{kam $\rightarrow$ käme}, \all{sah $\rightarrow$ sähe}, \all{gab $\rightarrow$ gäbe}, \all{fuhr $\rightarrow$ führe}. Quelques formes irrégulières à connaître par cœur : \all{sein $\rightarrow$ wäre}, \all{haben $\rightarrow$ hätte}, \all{geben $\rightarrow$ es gäbe}, \all{wissen $\rightarrow$ wüsste}, \all{werden $\rightarrow$ würde}.
\end{propriete}

\begin{center}
\begin{tabular}{|l|l|l|l|}
\hline
\rowcolor{poppurpleL} Infinitif & Präteritum & Konjunktiv II & Français \\
\hline
sein & war & ich wäre & je serais \\
\hline
haben & hatte & ich hätte & j'aurais \\
\hline
kommen & kam & ich käme & je viendrais \\
\hline
geben & gab & es gäbe & il y aurait \\
\hline
wissen & wusste & ich wüsste & je saurais \\
\hline
werden & wurde & ich würde & (auxiliaire du KII) \\
\hline
\end{tabular}
\end{center}

\begin{popfigure}
\begin{tikzpicture}
\node[draw=popblue, fill=popblueL, rounded corners, align=center] (i1) at (-4.5,0) {\textbf{Infinitiv}\\ \all{kommen}};
\node[draw=poporange, fill=poporangeL, rounded corners, align=center] (p1) at (0,0) {\textbf{Präteritum}\\ \all{kam}};
\node[draw=popgreen, fill=popgreenL, rounded corners, align=center] (k1) at (4.5,0) {\textbf{Konjunktiv II}\\ \all{käme} ou \all{würde kommen}};
\draw[-{Stealth}, very thick, popdark] (i1) -- (p1);
\draw[-{Stealth}, very thick, popdark] (p1) -- (k1) node[midway, above, font=\small, align=center]{tréma sur a/o/u\\ + désinences};
\node[draw=popmuted, fill=popcream, rounded corners, align=center, font=\small] at (0,-2) {En cas de doute : \all{würde} + infinitif $\rightarrow$ \all{ich würde kommen}};
\end{tikzpicture}
\legende{Frise de formation du KII : Infinitiv $\rightarrow$ Präteritum $\rightarrow$ Konjunktiv II.}
\end{popfigure}

\begin{methode}[Former le Konjunktiv II en trois étapes]
\begin{enumerate}
\item Trouver le \textbf{prétérit} du verbe : \all{kommen $\rightarrow$ kam}.
\item Ajouter le \textbf{tréma} sur a, o, u pour les verbes forts : \all{kam $\rightarrow$ käm-}.
\item Ajouter les \textbf{désinences} : \all{ich käme, du kämest, er käme, wir kämen}.
\end{enumerate}
En cas de doute (verbe faible, verbe rare), utiliser \all{würde} + infinitif : \all{ich würde tanzen} « je danserais ».
\end{methode}

Dans la pratique, on n'utilise les formes simples que pour \all{sein} (wäre), \all{haben} (hätte), les verbes modaux (\all{könnte, müsste, dürfte, sollte, wollte}) et quelques verbes fréquents (\all{gäbe, wüsste, käme, ginge}). Pour tous les autres verbes, on préfère \cle{würde + infinitif} : \all{Ich würde gern das Oktoberfest besuchen.} « J'aimerais visiter la fête de la bière. »

\courssub{Le subjonctif II passé et les emplois du KII}

\begin{propriete}[Le KII passé]
Le KII passé se forme avec \all{hätte} ou \all{wäre} + \textbf{participe II} : \all{Ich wäre gern zum Konzert gekommen.} « Je serais volontiers venu au concert. » ; \all{Ich hätte die Trommel gespielt.} « J'aurais joué du tambour. » On choisit \all{wäre} ou \all{hätte} selon l'auxiliaire du Perfekt (\all{sein} pour les verbes de mouvement et de changement d'état).
\end{propriete}

\textbf{Les quatre emplois essentiels du KII :}

\begin{itemize}
\item \textbf{Souhait irréel} : \all{Wenn ich doch mehr Zeit hätte!} « Si seulement j'avais plus de temps ! » ; \all{Ich wäre gern Sängerin.} « Je voudrais être chanteuse. »
\item \textbf{Hypothèse irréelle} : \all{Wenn ich Zeit hätte, würde ich das Konzert besuchen.} « Si j'avais le temps, j'irais au concert. »
\item \textbf{Conseil poli, demande atténuée} : \all{Du solltest mehr üben.} « Tu devrais t'entraîner davantage. » ; \all{Könnten Sie mir helfen?} « Pourriez-vous m'aider ? »
\item \textbf{Comparaison irréelle avec als ob} : \all{Er tanzt, als ob er ein Profi wäre.} « Il danse comme s'il était un professionnel. »
\end{itemize}

\begin{exemplebox}[Exemples traduits]
\all{Das berühmte Orchester spielt heute.} « L'orchestre célèbre joue aujourd'hui. »\par
\all{Wenn ich Zeit hätte, würde ich das Konzert besuchen.} « Si j'avais le temps, j'irais au concert. »\par
\all{Ich wäre gern Sängerin.} « Je voudrais être chanteuse. »\par
\all{Hätten wir Geld, würden wir nach Douala zum Festival fahren.} « Si nous avions de l'argent, nous irions au festival à Douala. »
\end{exemplebox}

\begin{attention}[« Je serais » n'est jamais *\all{ich würde sein}]
Les verbes \all{sein} et \all{haben} s'emploient presque toujours \textbf{sans} \all{würde} : « je serais » = \all{ich wäre}, « j'aurais » = \all{ich hätte}. La forme *\all{ich würde sein} est une erreur typique des francophones, qui calquent le conditionnel français « ser + ais ». Retenez : \all{wäre} et \all{hätte} sont déjà des formes de conditionnel.
\end{attention}

\begin{center}
\begin{tabularx}{\linewidth}{|X|X|X|}
\hline
\rowcolor{popblueL} Français & Deutsch & Remarque \\
\hline
je serais & ich wäre & jamais *würde sein \\
\hline
j'aurais & ich hätte & jamais *würde haben \\
\hline
j'irais au concert & ich würde zum Konzert gehen & würde + infinitif \\
\hline
si j'avais le temps & wenn ich Zeit hätte & imparfait français = KII allemand \\
\hline
comme si elle était célèbre & als ob sie berühmt wäre & als ob + KII \\
\hline
\end{tabularx}
\end{center}

\courssub{Texte d'application et entraînement}

\begin{textebox}[Ein Wunschkonzert in Yaoundé]
\all{Wenn ich reich wäre, würde ich ein großes Musikfestival in Yaoundé organisieren. Ich würde die besten Sänger Kameruns einladen, und auch ein berühmtes Orchester aus Deutschland käme zu diesem Fest. Die bekannten Musiker würden auf einer großen Bühne spielen, und das Publikum würde zu traditioneller und moderner Musik tanzen. Man hörte laute Trommeln und sähe bunte Tänze. Am Ende gäbe es lauten Applaus für alle Künstler.}\\[4pt]
\all{Meine Schwester hätte auch einen Wunsch: Sie wäre gern Sängerin. „Wenn ich eine schöne Stimme hätte“, sagt sie oft, „würde ich jeden Tag singen.“ Sie singt, als ob sie schon berühmt wäre! Ich sage ihr: „Du solltest mehr üben, dann könntest du wirklich Sängerin werden.“}\\[4pt]
\all{Leider haben wir wenig Geld. Wenn wir mehr Geld hätten, würden wir jedes Wochenende ein Konzert besuchen. Aber gute Musik kostet nicht immer viel: In unserem Viertel spielen junge Leute oft mit alten Trommeln, und die Nachbarn tanzen mit großer Freude. Das zeigt: Musik verbindet die Menschen, ob sie reich oder arm sind.}\\[4pt]
\textbf{Glossaire} : der Wunsch, die Wünsche (souhait) ; die Bühne, -n (scène) ; die Stimme, -n (voix) ; das Viertel, - (quartier) ; die Freude, -n (joie) ; verbinden (relier, unir).\par
\textbf{Questions} : 1) Was würde der Erzähler organisieren, wenn er reich wäre? 2) Was wäre seine Schwester gern? 3) Welchen Rat gibt er seiner Schwester? 4) Was zeigt das Beispiel des Viertels?
\end{textebox}

\begin{exoresolu}[Mettre au Konjunktiv II]
Énoncé : transformez au KII. 1) \all{Ich habe Zeit. Ich besuche das Konzert.} 2) \all{Sie ist Sängerin.} (souhait) 3) \all{Wir kommen zum Fest.} (passé irréel)
\tcblower
Solution détaillée : 1) \all{haben} $\rightarrow$ prétérit \all{hatte} $\rightarrow$ KII \all{hätte} ; pour \all{besuchen}, on utilise \all{würde} + infinitif : \all{Wenn ich Zeit hätte, würde ich das Konzert besuchen.} 2) \all{sein} $\rightarrow$ \all{wäre} (jamais *\all{würde sein}) : \all{Sie wäre gern Sängerin.} 3) KII passé avec \all{sein} car \all{kommen} est un verbe de mouvement : \all{Wir wären gern zum Fest gekommen.}
\end{exoresolu}

\begin{exercice}[À vous]
1) Complétez : \all{das modern\_\_ Lied / ein modern\_\_ Lied / modern\_\_ Musik}. 2) Traduisez : « Si j'avais une belle voix, je chanterais dans un chœur. » 3) Traduisez : « Nous serions volontiers venus au concert. » 4) Écrivez trois phrases sur vos souhaits avec \all{wäre}, \all{hätte} et \all{würde} + infinitif.
\end{exercice}

\begin{corrige}[Exercice : À vous]
1) \all{das moderne Lied} (article défini, nominatif/accusatif neutre : -e) ; \all{ein modernes Lied} (\all{ein} indécliné : l'adjectif porte -es) ; \all{moderne Musik} (sans article, nominatif/accusatif féminin : -e). 2) \all{Wenn ich eine schöne Stimme hätte, würde ich in einem Chor singen.} 3) \all{Wir wären gern zum Konzert gekommen.} 4) Exemple : \all{Ich wäre gern Musiker. Wenn ich ein Instrument hätte, würde ich jeden Tag spielen.}
\end{corrige}

\begin{aretenir}[L'essentiel de la section]
Après \all{der/die/das} : adjectif en \cle{-e} (5 formes) ou \cle{-en}. Après \all{ein} indécliné ou sans article : l'adjectif porte la marque du genre et du cas (\all{ein bekannter Sänger}, \all{mit lautem Applaus}). Le KII présent = radical du prétérit + tréma + désinences (\all{kam $\rightarrow$ käme}) ; formes clés : \all{wäre, hätte, gäbe, wüsste} ; sinon \all{würde} + infinitif. KII passé : \all{hätte/wäre} + participe II. « Je serais » = \all{ich wäre}, jamais *\all{ich würde sein}.
\end{aretenir}

\coursec{Méthode du commentaire et production guidée}

Dans les sections précédentes, nous avons lu et compris le texte \all{„Musik und Tanz in Kamerun und in Deutschland“}, enrichi notre lexique sur la musique et la danse, et travaillé la grammaire : la comparaison avec \all{wie} et \all{als ob}, la déclinaison de l'adjectif et le subjonctif II. Il s'agit maintenant de \cle{réinvestir} tout cela dans les deux épreuves écrites qui vous attendent au Baccalauréat : le \cle{commentaire de texte} et la \cle{production écrite guidée}. Cette dernière section vous donne une méthode complète, un modèle rédigé, les expressions utiles, un sujet de production avec sa grille d'auto-évaluation et une correction collective.

\courssub{La méthode du commentaire de texte au Bac}

Au Baccalauréat, le commentaire de texte (\all{der Kommentar}, \all{die Textanalyse}) n'est pas un résumé : c'est une \cle{analyse organisée} qui montre que vous avez compris le texte, que vous savez dégager ses idées principales et que vous pouvez donner votre avis en allemand correct. Il se construit toujours en trois parties.

\begin{methode}[Comment réussir un commentaire de texte en trois parties]
\begin{enumerate}
\item \textbf{L'introduction (\all{die Einleitung})} : présentez le texte en deux ou trois phrases. Indiquez le \cle{thème} (\all{das Thema}), éventuellement le type de texte (\all{ein Zeitungsartikel}, \all{eine Erzählung}), puis formulez la \cle{problématique} (\all{die Fragestellung}) : la question à laquelle le texte répond. Exemple : \all{Der Text handelt von der Musik in Kamerun und in Deutschland. Die Fragestellung lautet: Ist Musik eine universelle Sprache?}
\item \textbf{Le développement (\all{der Hauptteil})} : dégagez \cle{deux ou trois idées principales}. Pour chaque idée : annoncez-la en une phrase, puis \textbf{appuyez-la par une citation courte du texte entre guillemets}, puis commentez-la avec vos propres mots. Une idée sans citation n'est pas une analyse, c'est une opinion gratuite.
\item \textbf{La conclusion (\all{der Schluss})} : faites le \cle{bilan} (\all{die Bilanz}) en une phrase — qu'a montré le texte ? — puis donnez votre \cle{avis personnel} (\all{die eigene Meinung}) et, si possible, une \cle{ouverture} vers un problème plus large (la musique et la paix, la culture et la mondialisation).
\end{enumerate}
\end{methode}

\begin{popfigure}
\begin{tikzpicture}[
block/.style={rectangle, rounded corners=4pt, draw=popblue, fill=popblueL, text width=4.6cm, align=center, minimum height=2.2cm, font=\small},
fleche/.style={-{Stealth[length=3mm]}, thick, popdark}]
\node[block, align=center] (e) at (0,0){\textbf{Einleitung}\\ Thema + Fragestellung\\ \emph{„Der Text handelt von\dots“}};
\node[block, draw=poporange, fill=poporangeL, align=center] (h) at (6.2,0){\textbf{Hauptteil}\\ Idee 1 + Zitat\\ Idee 2 + Zitat\\ \emph{„Man kann sagen, dass\dots“}};
\node[block, draw=popgreen, fill=popgreenL, align=center] (s) at (12.4,0){\textbf{Schluss}\\ Bilanz + eigene Meinung\\ \emph{„Meiner Meinung nach\dots“}};
\draw[fleche] (e) -- (h);
\draw[fleche] (h) -- (s);
\end{tikzpicture}
\legende{La structure du commentaire de texte en trois blocs : introduction, développement, conclusion.}
\end{popfigure}

\courssub{Les expressions utiles pour commenter}

Un bon commentaire repose sur un petit stock de \cle{phrases toutes faites} (\all{die Redemittel}) qu'il faut mémoriser et réutiliser. Observez d'abord ces exemples, puis le tableau récapitulatif.

\begin{exemplebox}[Expressions de base pour le commentaire]
\all{Der Text handelt von einem Fest in einem kamerunischen Dorf.}\\
« Le texte traite d'une fête dans un village camerounais. »\\[4pt]
\all{Der Erzähler beschreibt, wie die Menschen tanzen und singen.}\\
« Le narrateur décrit comment les gens dansent et chantent. »\\[4pt]
\all{Meiner Meinung nach verbindet Musik die Menschen.}\\
« À mon avis, la musique unit les gens. »\\[4pt]
\all{Man kann sagen, dass Musik eine universelle Sprache ist.}\\
« On peut dire que la musique est une langue universelle. »
\end{exemplebox}

\begin{center}
\begin{tabularx}{\linewidth}{|X|X|}
\hline
\rowcolor{popblueL} \textbf{Deutsch} & \textbf{Français} \\
\hline
\all{Der Text handelt von\dots} & Le texte traite de\dots \\
\hline
\all{Das Thema des Textes ist\dots} & Le thème du texte est\dots \\
\hline
\all{Der Autor / Der Erzähler beschreibt\dots} & L'auteur / le narrateur décrit\dots \\
\hline
\all{Der Text zeigt, dass\dots} & Le texte montre que\dots \\
\hline
\all{Im Text steht: „\dots“} & Dans le texte, il est écrit : « \dots » \\
\hline
\all{Man kann sagen, dass\dots} & On peut dire que\dots \\
\hline
\all{Meiner Meinung nach\dots} & À mon avis\dots \\
\hline
\all{Ich finde es wichtig, dass\dots} & Je trouve important que\dots \\
\hline
\all{Zum Schluss kann man sagen\dots} & Pour conclure, on peut dire\dots \\
\hline
\end{tabularx}
\end{center}

\begin{attention}[La place du verbe dans les expressions de commentaire]
Après \all{Meiner Meinung nach}, le verbe conjugué reste en \textbf{deuxième position} : \all{Meiner Meinung nach \textbf{verbindet} Musik die Menschen.} (et non : \all{Meiner Meinung nach Musik verbindet\dots}). En revanche, après \all{dass}, le verbe conjugué va à la \textbf{fin} de la subordonnée : \all{Man kann sagen, dass Musik die Menschen \textbf{verbindet}.} C'est l'erreur la plus fréquente des candidats francophones, car en français l'ordre des mots ne change jamais.
\end{attention}

\courssub{Un commentaire modèle sur le texte support}

Lisons maintenant un commentaire complet, rédigé selon la méthode, sur le texte support de la leçon. Thème retenu : \cle{la musique comme langue universelle}.

\begin{textebox}[Kommentar — Musik als universelle Sprache]
\textbf{Einleitung.} Der Text handelt von der Musik und dem Tanz in Kamerun und in Deutschland. Er beschreibt ein traditionelles Fest in einem kamerunischen Dorf und das Musikleben in Deutschland. Die Fragestellung lautet: Ist Musik eine universelle Sprache, die alle Menschen verstehen?\\[6pt]
\textbf{Hauptteil.} Die erste Idee des Textes ist die Bedeutung der traditionellen Feste in Kamerun. Der Erzähler beschreibt ein Dorffest und schreibt: „Die Trommeln rufen die Menschen zusammen, und alle tanzen bis in die Nacht.“ Man kann sagen, dass die Musik hier das Herz der Gemeinschaft ist: Sie begleitet die Geburt, die Hochzeit und die Ernte. Die zweite Idee ist das Musikleben in Deutschland. Im Text steht: „In Deutschland gibt es viele Konzerte, Festivals und Chöre, in denen junge und alte Menschen zusammen singen.“ Das zeigt, dass die Musik auch dort die Menschen verbindet, aber auf eine andere Weise: weniger im Dorf, mehr in den Städten und in den Konzertsälen.\\[6pt]
\textbf{Schluss.} Zum Schluss kann man sagen, dass der Text die Musik als eine universelle Sprache darstellt. Meiner Meinung nach verbindet Musik die Menschen über alle Grenzen hinweg: Ein Kameruner und ein Deutscher können zusammen tanzen, ohne ein Wort zu sagen. Wenn alle Menschen so oft zusammen musizieren würden, gäbe es vielleicht mehr Frieden auf der Welt.
\end{textebox}

\textbf{Glossaire :} \all{die Bedeutung, -en} (la signification, l'importance) ; \all{die Gemeinschaft, -en} (la communauté) ; \all{die Hochzeit, -en} (le mariage) ; \all{die Ernte, -n} (la récolte) ; \all{der Chor, die Chöre} (le chœur) ; \all{der Konzertsaal, die Konzertsäle} (la salle de concert) ; \all{die Grenze, -n} (la frontière) ; \all{musizieren} (faire de la musique) ; \all{der Frieden} (la paix, sans pluriel).

Observez comment le modèle applique la méthode : chaque idée du développement est \textbf{annoncée}, \textbf{prouvée par une citation courte entre guillemets allemands „\dots“}, puis \textbf{commentée}. La conclusion contient le bilan (\all{Zum Schluss kann man sagen\dots}), l'avis personnel (\all{Meiner Meinung nach\dots}) et une ouverture au subjonctif II (\all{gäbe es vielleicht mehr Frieden}) — ce qui réinvestit la grammaire de la leçon. Notez aussi la forme \all{gäbe es} : c'est le subjonctif II de \all{es gibt} (\all{geben} $\rightarrow$ \all{gäbe}), très utile pour exprimer une hypothèse.

\courssub{La production écrite guidée : « Ein Fest in meiner Region »}

Au Bac, la production écrite est souvent un paragraphe de 8 à 10 lignes sur un sujet proche du texte support. Voici la démarche à suivre.

\begin{methode}[Rédiger un paragraphe guidé en cinq étapes]
\begin{enumerate}
\item \textbf{Lire la consigne deux fois} et souligner les exigences : nombre de lignes, mots du lexique à employer, structures grammaticales imposées.
\item \textbf{Faire un plan rapide au brouillon} : 2 ou 3 idées (déroulement de la fête, musique et danses, mon impression personnelle).
\item \textbf{Choisir les mots du lexique} de la leçon : \all{das Fest}, \all{die Trommel}, \all{der Sänger}, \all{die Musik}, \all{der Tanz}, \all{das Konzert}.
\item \textbf{Rédiger} en vérifiant à chaque phrase : verbe en 2e position dans la principale, verbe conjugué à la fin après \all{als ob}, \all{dass}, \all{wenn}, \all{weil} ; adjectif décliné devant le nom ; majuscule à tous les noms.
\item \textbf{Se relire} avec la grille d'auto-évaluation ci-dessous.
\end{enumerate}
\end{methode}

\begin{exoresolu}[Production guidée : Ein Fest in meiner Region]
\textbf{Énoncé.} Rédigez un paragraphe de 8 à 10 lignes intitulé \all{„Ein Fest in meiner Region“}. Utilisez au moins cinq mots du lexique de la leçon, une comparaison avec \all{als ob}, un adjectif décliné devant un nom et une phrase au subjonctif II.
\tcblower
\textbf{Solution détaillée (modèle de production).}\\[4pt]
\all{Ein Fest in meiner Region}\\[4pt]
\all{In meiner Region gibt es jedes Jahr ein großes Fest. Die ganze Familie kommt zusammen, und die Menschen tragen schöne traditionelle Kleider. Am Abend beginnt die Musik: Die Trommeln spielen laut, und ein berühmter Sänger singt alte Lieder. Die Kinder tanzen, als ob sie fliegen würden, und alle lachen und klatschen. Es ist, als ob das ganze Dorf eine einzige Familie wäre. Ich liebe dieses Fest, weil die Musik alle Menschen verbindet. Wenn ich mehr Zeit hätte, würde ich jedes Wochenende zu einem Konzert gehen. Solche Feste sind wichtig, denn sie zeigen den jungen Leuten die Kultur unserer Väter.}\\[6pt]
\textbf{Vérification point par point :}
\begin{itemize}
\item \textbf{Lexique (5 mots minimum) :} \all{das Fest}, \all{die Musik}, \all{die Trommeln}, \all{der Sänger}, \all{tanzen}, \all{das Konzert} — consigne remplie.
\item \textbf{Comparaison avec \all{als ob} :} \all{Die Kinder tanzen, als ob sie fliegen \textbf{würden}} — verbe conjugué à la fin, subjonctif II obligatoire après \all{als ob} (comparaison irréelle : les enfants ne volent pas vraiment).
\item \textbf{Adjectif décliné :} \all{ein groß\textbf{es} Fest}, \all{schöne traditionell\textbf{e} Kleider}, \all{ein berühmt\textbf{er} Sänger}, \all{alt\textbf{e} Lieder}.
\item \textbf{Subjonctif II :} \all{Wenn ich mehr Zeit \textbf{hätte}, \textbf{würde} ich jedes Wochenende zu einem Konzert \textbf{gehen}.} — hypothèse irréelle : \all{hätte} (subjonctif II de \all{haben}) dans la subordonnée, \all{würde} + infinitif dans la principale.
\item \textbf{Place du verbe :} vérifiée dans chaque phrase (2e position dans les principales, finale dans les subordonnées ; dans \all{Wenn ich mehr Zeit hätte, würde ich\dots}, la principale commence par le verbe car la subordonnée occupe la première place).
\end{itemize}
\end{exoresolu}

\begin{attention}[Les erreurs fréquentes dans la production écrite]
\begin{itemize}
\item \textbf{Place du verbe :} après \all{als ob}, \all{dass}, \all{wenn}, \all{weil}, le verbe conjugué va à la \textbf{fin}. Écrivez \all{weil die Musik alle Menschen \textbf{verbindet}}, et non \all{weil die Musik verbindet alle Menschen}.
\item \textbf{L'adjectif non décliné :} en français l'adjectif s'accorde avec le nom, en allemand il prend une \textbf{terminaison} : \all{ein groß\textbf{es} Fest}, \all{die schön\textbf{e} Musik}. Ne laissez jamais un adjectif « nu » devant un nom.
\item \textbf{La majuscule des noms :} \all{die musik} est fautif ; écrivez \all{die Musik}. Chaque nom, partout dans la phrase, prend une majuscule.
\item \textbf{Le subjonctif II oublié après \all{als ob} :} \all{als ob sie fliegen \textbf{würden}}, et non l'indicatif \all{als ob sie fliegen}.
\item \textbf{Le faux ami \all{bekommen} :} il signifie « recevoir », pas « devenir » ; « devenir » se dit \all{werden} (\all{ich werde}, \all{ich wurde}, \all{ich würde}).
\end{itemize}
\end{attention}

\courssub{Grille d'auto-évaluation et correction collective}

Avant de rendre votre copie, évaluez-vous vous-même avec cette grille. Elle reprend les grands critères utilisés par les correcteurs.

\begin{center}
\begin{tabularx}{\linewidth}{|X|c|X|}
\hline
\rowcolor{popblueL} \textbf{Critère} & \textbf{Points} & \textbf{Question à se poser} \\
\hline
Contenu et cohérence & /4 & Ai-je répondu au sujet avec 2-3 idées claires ? \\
\hline
Vocabulaire & /4 & Ai-je employé au moins 5 mots du lexique de la leçon ? \\
\hline
Grammaire : KII et \all{als ob} & /4 & Mes phrases au subjonctif II sont-elles correctes ? \\
\hline
Grammaire : déclinaison de l'adjectif & /4 & Mes adjectifs devant les noms ont-ils une terminaison ? \\
\hline
Orthographe, majuscules, ponctuation & /4 & Tous les noms ont-ils une majuscule ? \\
\hline
\textbf{Total} & \textbf{/20} & \\
\hline
\end{tabularx}
\end{center}

\begin{exoresolu}[Correction collective d'une production type]
\textbf{Énoncé.} Voici le début de la copie d'un élève. Repérez et corrigez les cinq erreurs.\\[4pt]
\all{„In meine Region gibt es ein groß Fest. Die musik spielt laut, und die Kinder tanzen, als ob sie fliegen. Meiner Meinung nach Musik verbindet die Menschen. Wenn ich Zeit hätte, ich würde jedes Wochenende tanzen gehen.“}
\tcblower
\textbf{Solution détaillée.}
\begin{enumerate}
\item \all{In meine Region} $\rightarrow$ \all{In mein\textbf{er} Region} : après \all{in} + datif (localisation), le possessif devant un nom féminin prend la terminaison \all{-er} (\all{meiner Region}).
\item \all{ein groß Fest} $\rightarrow$ \all{ein groß\textbf{es} Fest} : adjectif décliné devant un nom neutre après \all{ein} (terminaison \all{-es} au nominatif/accusatif neutre).
\item \all{Die musik} $\rightarrow$ \all{Die \textbf{M}usik} : majuscule obligatoire à tous les noms.
\item \all{als ob sie fliegen} $\rightarrow$ \all{als ob sie fliegen \textbf{würden}} : le subjonctif II est obligatoire après \all{als ob} (comparaison irréelle).
\item \all{Meiner Meinung nach Musik verbindet} $\rightarrow$ \all{Meiner Meinung nach \textbf{verbindet} Musik} : le verbe conjugué reste en deuxième position ; et \all{Wenn ich Zeit hätte, ich würde\dots} $\rightarrow$ \all{Wenn ich Zeit hätte, \textbf{würde} ich\dots} : quand la subordonnée précède, la principale commence par le verbe conjugué (inversion).
\end{enumerate}
\textbf{Version corrigée :} \all{„In meiner Region gibt es ein großes Fest. Die Musik spielt laut, und die Kinder tanzen, als ob sie fliegen würden. Meiner Meinung nach verbindet Musik die Menschen. Wenn ich Zeit hätte, würde ich jedes Wochenende tanzen gehen.“}
\end{exoresolu}

\begin{experience}[Atelier d'écriture en classe]
Par groupes de quatre : (1) chaque élève rédige son paragraphe \all{„Ein Fest in meiner Region“} en 15 minutes ; (2) les copies tournent dans le groupe et chacun corrige la copie d'un camarade avec la grille d'auto-évaluation ; (3) le groupe choisit la meilleure production et la lit à la classe ; (4) la classe identifie ensemble les erreurs restantes au tableau. L'enseignant note les erreurs récurrentes (place du verbe, déclinaison, majuscules) pour une séance de remédiation.
\end{experience}

\begin{savaistu}[La stratégie gagnante au Bac]
Au Baccalauréat, la production écrite est évaluée sur trois grands critères : la \textbf{richesse du vocabulaire}, la \textbf{correction grammaticale} et la \textbf{cohérence} du texte. Une stratégie payante consiste à \cle{réutiliser les structures du texte support} : les tournures, les verbes et les expressions que vous avez rencontrés dans le texte sont forcément corrects et adaptés au sujet. Recopier intelligemment une structure en changeant les mots n'est pas du plagiat — c'est exactement ce que font les bons candidats. Par exemple, si le texte dit \all{„Die Trommeln rufen die Menschen zusammen“}, vous pouvez écrire dans votre production : \all{„Die Musik ruft die ganze Familie zusammen.“}
\end{savaistu}

\begin{exercice}[À vous de jouer]
\begin{enumerate}
\item Rédigez un commentaire de 10 à 12 lignes sur le texte support de la leçon, avec une introduction, deux idées appuyées par des citations et une conclusion avec votre avis personnel.
\item Rédigez votre propre paragraphe \all{„Ein Fest in meiner Region“} (8 à 10 lignes) en respectant toutes les contraintes, puis évaluez-vous avec la grille.
\item Traduisez en allemand : « À mon avis, les fêtes traditionnelles sont très importantes. On peut dire que la musique unit les peuples. Si j'avais plus de temps, j'apprendrais à jouer de la trommel. »
\end{enumerate}
\end{exercice}

\begin{corrige}[Exercice — question 3]
\all{Meiner Meinung nach sind die traditionellen Feste sehr wichtig. Man kann sagen, dass Musik die Völker verbindet. Wenn ich mehr Zeit hätte, würde ich Trommel spielen lernen.}\\
Points de vigilance : verbe en 2e position après \all{Meiner Meinung nach} ; verbe conjugué à la fin après \all{dass} ; \all{hätte} et \all{würde\dots lernen} pour le subjonctif II ; \all{Trommel spielen} sans article, comme \all{Klavier spielen} ; majuscules à \all{Feste}, \all{Musik}, \all{Völker}, \all{Zeit}, \all{Trommel}.
\end{corrige}

\newpage

\coursec{Activité d'intégration}

\courssub{Objectifs de l'activité}

À l'issue de cette activité d'intégration, l'élève sera capable de :
\begin{itemize}
\item réinvestir le \cle{vocabulaire thématique} de la musique et de la danse (\all{die Musik, der Tanz, das Konzert, die Trommel, der Sänger, das Fest}) dans une situation de communication authentique ;
\item employer correctement la \cle{comparaison} avec \all{wie} et la comparaison irréelle avec \all{als ob} suivi du subjonctif II ;
\item accorder et décliner correctement les \cle{adjectifs} (avec article défini, indéfini ou sans article) ;
\item utiliser le \cle{subjonctif II} pour exprimer des souhaits et des hypothèses (\all{Wenn wir mehr Geld hätten, würden wir ...}) ;
\item produire collectivement un document écrit (affiche + article) et le présenter oralement en deux minutes.
\end{itemize}

\courssub{Situation}

Le lycée de Yaoundé organise, à la fin de l'année scolaire, une semaine culturelle intitulée \all{„Deutsch-Kamerun-Woche“}. Le club d'allemand, en partenariat avec l'Institut culturel allemand, souhaite clôturer la semaine par un grand festival de musique et de danse : le \all{„Deutsch-Kamerun-Musikfest“}. Chaque groupe d'élèves de Terminale A est chargé de concevoir la promotion de ce festival : une affiche et un article pour le journal du lycée.

Pour vous inspirer, lisez ce modèle d'annonce publié dans le journal d'un lycée partenaire de Bonn :

\begin{textebox}[Ankündigung im Schülerjournal — Bonn]
\textbf{Musikfest am Samstag!}\par
Am Samstag, den 15. Juni, findet auf dem Schulhof unser großes Musikfest statt. Es beginnt um 16 Uhr und endet um 22 Uhr. Die bekannte Sängerin Lena Mballa singt deutsche und kamerunische Lieder. Die Schülerband „Die wilden Trommeln“ spielt traditionelle und moderne Musik. Wer tanzen will, kann mit der Tanzgruppe „Makossa Stars“ auf die Bühne kommen. Der Eintritt ist frei, aber wir sammeln Geld für die neue Schulbibliothek. Wenn wir mehr Geld hätten, würden wir auch ein berühmtes Orchester einladen! Kommt alle — es wird ein unvergessliches Fest!
\end{textebox}

\textbf{Glossaire :} \all{der Schulhof} (la cour de l'école) ; \all{die Schülerband} (le groupe des élèves) ; \all{die Bühne, -n} (la scène) ; \all{der Eintritt} (l'entrée) ; \all{sammeln} (collecter) ; \all{einladen} (inviter, verbe à particule séparable) ; \all{unvergesslich} (inoubliable).

\courssub{Consignes}

Travail par groupes de quatre élèves. Chaque groupe produit deux documents et une présentation orale.

\begin{enumerate}
\item \textbf{Brainstorming lexical (5 min).} Sur une feuille, listez au moins dix mots du thème appris dans la leçon (instruments, artistes, lieux, fêtes) avec leur article et leur pluriel : \all{die Trommel, -n ; der Sänger, - ; das Konzert, -e}...
\item \textbf{L'affiche (15 min).} Rédigez une affiche en allemand pour le \all{„Deutsch-Kamerun-Musikfest“} de votre lycée. Elle doit contenir : un titre accrocheur, la date, le lieu (par exemple \all{auf dem Schulhof des Gymnasiums Yaoundé}), les artistes et les instruments. Elle doit employer \textbf{au moins trois adjectifs correctement déclinés} (ex. \all{das große Fest, eine bekannte Sängerin, traditionelle Musik}).
\item \textbf{L'article (20 min).} Rédigez un article d'environ dix lignes pour le journal du lycée. Il doit obligatoirement contenir :
\begin{itemize}
\item une comparaison avec \all{wie} (ex. \all{Die Trommeln klingen wie Donner.}) ;
\item une phrase avec \all{als ob} + subjonctif II (ex. \all{Die Sängerin tanzt, als ob sie zwanzig Jahre alt wäre.}) ;
\item deux phrases au subjonctif II exprimant un souhait (ex. \all{Wenn wir mehr Geld hätten, würden wir ein Orchester einladen.}).
\end{itemize}
\item \textbf{Présentation orale (2 min par groupe).} Présentez votre affiche et lisez votre article devant la classe. Chaque membre du groupe prend la parole.
\item \textbf{Vote et correction collective (10 min).} La classe vote pour le meilleur festival. Le professeur corrige collectivement les erreurs relevées, en ciblant la déclinaison de l'adjectif et le subjonctif II.
\end{enumerate}

\courssub{Ressources à mobiliser}

\begin{aretenir}[La comparaison : wie et als ob]
\begin{itemize}
\item Comparaison réelle : \all{wie} + indicatif. \all{Er singt wie ein Profi.} « Il chante comme un professionnel. »
\item Comparaison irréelle : \all{als ob} + subjonctif II, verbe à la fin. \all{Sie tanzt, als ob sie eine berühmte Tänzerin wäre.} « Elle danse comme si elle était une danseuse célèbre. »
\item Formes utiles du KII : \all{wäre} (serait), \all{hätte} (aurait), \all{würde} + infinitif, \all{könnte} (pourrait), \all{gäbe} (\all{es gäbe} : il y aurait).
\end{itemize}
\end{aretenir}

\begin{aretenir}[La déclinaison de l'adjectif (rappel)]
\begin{center}
\begin{tabular}{|l|l|l|}
\hline
\rowcolor{popblueL} Article défini & Article indéfini & Sans article \\
\hline
\all{das große Fest} & \all{ein großes Fest} & \all{großes Fest} \\
\all{die bekannte Sängerin} & \all{eine bekannte Sängerin} & \all{bekannte Sängerin} \\
\all{der laute Tanz} & \all{ein lauter Tanz} & \all{lauter Tanz} \\
\hline
\end{tabular}
\end{center}
Après article défini : terminaison \all{-e} ou \all{-en}. Après article indéfini : l'adjectif porte la marque du genre (\all{-er, -es, -e}). Sans article : l'adjectif porte seul la marque.
\end{aretenir}

\begin{aretenir}[Lexique utile]
\all{die Musik} (la musique) ; \all{der Tanz, „e} (la danse) ; \all{das Konzert, -e} (le concert) ; \all{die Trommel, -n} (le tambour) ; \all{der Sänger, - / die Sängerin, -nen} (le chanteur / la chanteuse) ; \all{das Fest, -e} (la fête) ; \all{die Bühne, -n} (la scène) ; \all{das Lied, -er} (la chanson) ; \all{auftreten} (se produire sur scène) ; \all{stattfinden} (avoir lieu) ; \all{der Eintritt ist frei} (l'entrée est gratuite).
\end{aretenir}

\begin{attention}[Pièges fréquents]
\begin{itemize}
\item Après \all{als ob}, le verbe conjugué va \textbf{à la fin} : \all{als ob sie krank wäre} (et jamais \all{als ob sie wäre krank}).
\item Ne pas confondre \all{als} (comparaison d'inégalité : \all{größer als}) et \all{wie} (comparaison d'égalité : \all{so laut wie}).
\item Au subjonctif II, on écrit \all{ich hätte}, \all{ich wäre} — pas de confusion avec le passé composé français.
\item Majuscule à tous les noms sur l'affiche : \all{das Fest, die Trommel, der Sänger}.
\end{itemize}
\end{attention}

\courssub{Proposition de production et corrigé commenté}

\begin{exoresolu}[Modèle d'affiche et d'article]
\textbf{Affiche :}\par
\all{\textbf{DEUTSCH-KAMERUN-MUSIKFEST!}}\par
\all{Samstag, 20. Juni, 16 Uhr — auf dem großen Schulhof des Gymnasiums Yaoundé}\par
\all{Die bekannte Sängerin Amina Ngo Bell und die Gruppe „Bendskin Boys“}\par
\all{Traditionelle Trommeln, moderne Musik und wilder Tanz!}\par
\all{Der Eintritt ist frei. Kommt alle!}\par
\tcblower
\textbf{Article modèle :}\par
\all{Am Samstag fand auf dem Schulhof unseres Gymnasiums das erste Deutsch-Kamerun-Musikfest statt. Hunderte von Schülern, Lehrern und Eltern waren gekommen. Die bekannte Sängerin Amina Ngo Bell sang deutsche und kamerunische Lieder, und ihre Stimme klang wie ein Geschenk des Himmels. Die Gruppe „Bendskin Boys“ spielte auf traditionellen Trommeln, und die Tänzer bewegten sich, als ob sie niemals müde würden. Das Publikum klatschte und tanzte bis in die Nacht. Leider hatten wir nur eine kleine Bühne. Wenn wir mehr Geld hätten, würden wir eine größere Bühne bauen. Wenn es nächstes Jahr mehr Sponsoren gäbe, könnten wir auch eine berühmte deutsche Band einladen. Es war ein unvergessliches Fest, und alle hoffen auf eine zweite Auflage!}\par
\medskip
\textbf{Commentaire des choix de langue :}
\begin{itemize}
\item \textbf{Comparaison avec wie :} \all{ihre Stimme klang wie ein Geschenk des Himmels} — comparaison réelle, indicatif (\all{klang}), \all{wie} + nominatif.
\item \textbf{als ob + KII :} \all{als ob sie niemals müde würden} — verbe conjugué à la fin, subjonctif II avec \all{würden} + infinitif.
\item \textbf{Souhaits au KII :} \all{Wenn wir mehr Geld hätten, würden wir eine größere Bühne bauen} (conditionnel : \all{hätten} + \all{würden bauen}) ; \all{Wenn es mehr Sponsoren gäbe, könnten wir ... einladen} (\all{gäbe}, KII de \all{geben}).
\item \textbf{Adjectifs déclinés :} \all{auf dem großen Schulhof} (datif après article défini : \all{-en}) ; \all{die bekannte Sängerin} (nominatif féminin : \all{-e}) ; \all{auf traditionellen Trommeln} (datif pluriel : \all{-en}) ; \all{eine größere Bühne} (accusatif féminin après article indéfini : \all{-e}) ; \all{ein unvergessliches Fest} (neutre après \all{ein} : \all{-es}).
\item \textbf{Lexique du thème :} \all{das Musikfest, die Sängerin, die Trommeln, die Bühne, das Fest, einladen, stattfinden}.
\end{itemize}
\end{exoresolu}

\courssub{Bilan de l'activité}

Cette activité a permis de mobiliser simultanément les quatre savoir-faire de la leçon : la compréhension d'un modèle écrit, l'emploi du vocabulaire de la musique et de la danse, l'application des règles de grammaire (comparaison avec \all{wie} et \all{als ob}, déclinaison de l'adjectif, subjonctif II) et la production écrite et orale. La correction collective a montré que les points les plus délicats restent la place du verbe après \all{als ob} et le choix de la terminaison adjectivale sans article. Pour t'évaluer toi-même, vérifie que ton article contient bien : une comparaison avec \all{wie}, une phrase avec \all{als ob} + KII, deux souhaits au subjonctif II et au moins trois adjectifs correctement déclinés. Si l'un de ces éléments manque, reprends les boîtes « À retenir » de la leçon avant le devoir.

\begin{exercice}[Pour aller plus loin]
Rédige individuellement un court paragraphe (6 lignes) intitulé \all{„Mein Traumkonzert“} (« Mon concert de rêve ») en utilisant au moins deux phrases au subjonctif II (\all{Wenn ich ... hätte / wäre, würde ich ...}) et une comparaison avec \all{als ob}.
\end{exercice}

\newpage

\coursec{Méthodes et exercices résolus}

\coursec{Leçon AL23 : La vie culturelle 2 — La musique et la danse}

\courssub{Texte support : Musik und Tanz in Kamerun und in Deutschland}

\begin{textebox}[Texte original — compréhension écrite/orale]
\all{Musik verbindet Menschen – in Kamerun wie in Deutschland. In Kamerun gehört die Trommel zu jedem wichtigen Ereignis: bei Hochzeiten, Geburtsfeiern und traditionellen Festen bestimmen Rhythmen und Tänze den Ablauf. Die Tänzer bewegen sich so schnell, als ob sie fliegen würden. In Deutschland besuchen junge Leute oft Konzerte und Festivals, wo Bands moderne Musik spielen. Doch auch dort gibt es Volksfeste mit traditioneller Musik. Ob Makossa, Bikutsi oder Techno: Musik ist eine Sprache, die alle verstehen.}
\end{textebox}

\begin{definition}[Vocabulaire clé du texte (avec articles et pluriels)]
\begin{center}
\begin{tabularx}{\linewidth}{|X|X|X|}
\hline
\rowcolor{popblueL} \textbf{Deutsch} & \textbf{Français} & \textbf{Remarque} \\
\hline
\all{die Musik, -en} & la musique & féminin ; \all{Musik hören} (écouter de la musique) \\
\all{der Tanz, die Tänze} & la danse & masculin ; \all{tanzen} (danser) \\
\all{das Konzert, -e} & le concert & neutre ; \all{ein Konzert besuchen} (assister à un concert) \\
\all{die Trommel, -n} & le tam-tam, le tambour & féminin ; \all{Trommel spielen} (jouer du tam-tam) \\
\all{der Sänger, - (Pl. die Sänger)} & le chanteur & masculin ; féminin : \all{die Sängerin, -nen} \\
\all{das Fest, -e} & la fête & neutre ; \all{ein Fest feiern} (fêter) \\
\all{das Lied, -er} & la chanson & neutre \\
\all{der Rhythmus, die Rhythmen} & le rythme & masculin \\
\all{die Band, -s} & le groupe, le band & féminin (emprunt à l'anglais) \\
\all{der Musiker, -} & le musicien & masculin ; féminin : \all{die Musikerin, -nen} \\
\all{auftreten (tritt auf, trat auf, ist aufgetreten)} & se produire (sur scène) & verbe à particule séparable \\
\all{bestimmen (bestimmt, bestimmte, hat bestimmt)} & déterminer, dicter & + accusatif \\
\hline
\end{tabularx}
\end{center}
\end{definition}

\courssub{Compréhension du texte}

\begin{methode}[Lire et comprendre un texte culturel (écrit/oral) — Rappel]
\begin{enumerate}
\item \textbf{Première lecture globale.} Lis le texte une fois sans dictionnaire. Repère le thème (\all{Musik und Tanz}), le type de texte (article d'opinion / reportage) et les lieux cités (\all{Kamerun, Deutschland}).
\item \textbf{Repérage lexical.} Souligne les mots du champ lexical : \all{Trommel, Rhythmus, Tanz, Konzert, Festival, Band, Volksfest, Sprache}. Ces mots t'aident à deviner le sens global.
\item \textbf{Lecture détaillée.} Pour chaque question, localise le passage exact. Ne réponds jamais « de tête » : cite ou reformule la phrase du texte.
\item \textbf{Répondre en phrases complètes.} Une réponse au Bac commence par un sujet et un verbe conjugué : \all{Die Trommel gehört zu jedem wichtigen Ereignis.} Verbe en 2\textsuperscript{e} position, majuscules aux noms, cas corrects.
\item \textbf{Vérification.} Relis : ordre des mots (principal / subordonnée), déclinaison après prépositions (\all{bei} + datif, \all{in} + datif lieu), particule séparable en fin de phrase.
\end{enumerate}
\end{methode}

\begin{exoresolu}[Compréhension détaillée — type Bac]
\textbf{Questions sur le texte :}
1. \all{Wozu gehört die Trommel in Kamerun?}
2. \all{Wie bewegen sich die Tänzer? (Vergleich)}
3. \all{Was machen junge Leute in Deutschland oft?}
4. \all{Was gibt es in Deutschland auch?}
5. \all{Was ist Musik nach dem letzten Satz?}

\tcblower
\textbf{Raisonnement et réponses modèles :}

\textbf{1.} Segment : \all{In Kamerun gehört die Trommel zu jedem wichtigen Ereignis.} \\
Réponse : \all{Die Trommel gehört in Kamerun zu jedem wichtigen Ereignis.} (Préposition \all{zu} + datif : \all{jedem Ereignis}).

\textbf{2.} Segment : \all{Die Tänzer bewegen sich so schnell, als ob sie fliegen würden.} \\
Réponse : \all{Die Tänzer bewegen sich so schnell, als ob sie fliegen würden.} (Comparaison irréelle : \all{als ob} + Konjunktiv II \all{würden fliegen}, verbe à la fin).

\textbf{3.} Segment : \all{In Deutschland besuchen junge Leute oft Konzerte und Festivals...} \\
Réponse : \all{In Deutschland besuchen junge Leute oft Konzerte und Festivals, wo Bands moderne Musik spielen.} (Relative \all{wo} + verbe à la fin).

\textbf{4.} Segment : \all{Doch auch dort gibt es Volksfeste mit traditioneller Musik.} \\
Réponse : \all{In Deutschland gibt es auch Volksfeste mit traditioneller Musik.} (\all{es gibt} + accusatif : \all{Volksfeste}).

\textbf{5.} Segment : \all{Musik ist eine Sprache, die alle verstehen.} \\
Réponse : \all{Musik ist eine Sprache, die alle verstehen.} (Relative \all{die} = \all{die Sprache}, accusatif féminin, verbe \all{verstehen} à la fin).

\textbf{Conclusion :} Les réponses sont courtes, exactes, appuyées sur le texte, avec une syntaxe allemande correcte.
\end{exoresolu}

\begin{exercice}[Compréhension — À vous de jouer]
Répondez en allemand, en phrases complètes, en vous appuyant sur le texte.
1. \all{Bei welchen Anlässen spielt die Trommel eine Rolle?}
2. \all{Wo spielen Bands moderne Musik?}
3. \all{Welche Musikstile nennt der Text am Ende?}
4. \all{Warum versteht alle die Musik?}
\end{exercice}

\begin{corrige}[Exercice Compréhension]
1. \all{Die Trommel spielt bei Hochzeiten, Geburtsfeiern und traditionellen Festen eine Rolle.}
2. \all{Bands spielen moderne Musik auf Konzerten und Festivals.}
3. \all{Der Text nennt Makossa, Bikutsi und Techno.}
4. \all{Weil Musik eine Sprache ist, die alle verstehen.}
\end{corrige}

\courssub{Lexique thématique : die Musik und der Tanz}

\begin{propriete}[Familles de mots et expressions utiles — À mémoriser]
\begin{center}
\begin{tabularx}{\linewidth}{|X|X|}
\hline
\rowcolor{popblueL} \textbf{Famille / Verbe} & \textbf{Noms, adjectifs, expressions} \\
\hline
\all{singen} (chanter) & \all{der Sänger, die Sängerin, das Lied, der Gesang, das Gesangbuch} \\
\all{tanzen} (danser) & \all{der Tanz, der Tänzer, die Tänzerin, die Tanzschule, das Tanzfest} \\
\all{spielen} (jouer) & \all{das Instrument, die Trommel, die Gitarre, das Klavier; ein Instrument spielen} \\
\all{auftreten} (se produire) & \all{der Auftritt, die Bühne, das Konzert, das Festival; auf der Bühne auftreten} \\
\all{besuchen} (assister à) & \all{ein Konzert besuchen, ein Festival besuchen} \\
\all{feiern} (fêter) & \all{das Fest, die Feier, die Party; ein Fest feiern, Geburtstag feiern} \\
\all{hören} (écouter) & \all{Musik hören, ein Lied hören, dem Sänger zuhören (datif)} \\
\hline
\end{tabularx}
\end{center}
\end{propriete}

\begin{aretenir}[Adjectifs fréquents pour décrire la musique]
\begin{center}
\begin{tabular}{|l|l|l|}
\hline
\rowcolor{popblueL} \textbf{Deutsch} & \textbf{Français} & \textbf{Exemple} \\
\hline
\all{traditionell} & traditionnel & \all{traditionelle Musik} \\
\all{modern} & moderne & \all{moderne Bands} \\
\all{berühmt} & célèbre & \all{ein berühmter Sänger} \\
\all{laut / leise} & bruyant / doux & \all{laute Trommeln, leise Melodien} \\
\all{schnell / langsam} & rapide / lent & \all{schneller Rhythmus, langsamer Tanz} \\
\all{lebendig} & vivant, animé & \all{lebendige Musik} \\
\all{beliebt} & populaire & \all{beliebte Lieder} \\
\hline
\end{tabular}
\end{center}
\end{aretenir}

\begin{exoresolu}[Lexique — Complétion et transformation]
\textbf{A.} Complétez avec le mot convenable (forme correcte) : \all{Konzert, Trommel, Sängerin, tanzen, Fest, Band}.

1. \all{Auf dem Dorf ... spielen die Musiker die ... .} \\
2. \all{Die ... singt ein schönes ... .} \\
3. \all{Morgen besuchen wir ein ... in Yaoundé.} \\
4. \all{Die Jugendlichen ... bis Mitternacht.} \\
5. \all{Die ... tritt um 20 Uhr auf.}

\textbf{B.} Formez le nom d'agent (masculin et féminin) : \all{singen, tanzen, musizieren, spielen (Instrument)}.

\tcblower
\textbf{A.}
1. \all{Auf dem Dorffest spielen die Musiker die Trommel.} (\all{das Fest} $\rightarrow$ \all{dem Fest} datif ; \all{die Trommel} accusatif).
2. \all{Die Sängerin singt ein schönes Lied.} (\all{die Sängerin} féminin ; \all{ein schönes Lied} accusatif neutre).
3. \all{Morgen besuchen wir ein Konzert in Yaoundé.} (\all{ein Konzert} accusatif neutre).
4. \all{Die Jugendlichen tanzen bis Mitternacht.} (Verbe \all{tanzen}, 3\textsuperscript{e} pers. plur.).
5. \all{Die Band tritt um 20 Uhr auf.} (Verbe à particule \all{auftreten} : \all{tritt ... auf}).

\textbf{B.}
\all{singen} $\rightarrow$ \all{der Sänger / die Sängerin} \\
\all{tanzen} $\rightarrow$ \all{der Tänzer / die Tänzerin} \\
\all{musizieren} $\rightarrow$ \all{der Musiker / die Musikerin} \\
\all{spielen} $\rightarrow$ \all{der Spieler / die Spielerin} (mais pour musicien : \all{Musiker}).
\end{exoresolu}

\courssub{Grammaire 1 : La comparaison avec « wie » et « als ob »}

\begin{propriete}[Les trois types de comparaison — Forme et emploi]
\begin{enumerate}
\item \textbf{Comparaison d'égalité (aussi ... que) :} \all{so ... wie} / \all{genauso ... wie}. \\
\all{wie} introduit un groupe nominal ou une proposition \textbf{sans déplacement du verbe} (verbe en 2\textsuperscript{e} position si proposition principale). \\
\all{Er singt so schön wie sein Vater.} \all{Sie tanzt genauso gut wie er.}
\item \textbf{Comparaison de supériorité (plus ... que) :} \textbf{comparatif en \all{-er}} + \all{als}. \\
\all{Die Trommel ist lauter als die Gitarre.} \all{Dieses Fest ist wichtiger als jenes.} \\
\all{als} est suivi du même cas que devant (souvent nominatif/accusatif).
\item \textbf{Comparaison irréelle / hypothétique (comme si) :} \all{als ob} + \textbf{Konjunktiv II}, verbe conjugué \textbf{à la fin de la subordonnée}. \\
\all{Sie tanzt, als ob sie schweben würde.} (Présent : \all{würde} + infinitif). \\
\all{Er tut, als ob er krank wäre.} (\all{sein} $\rightarrow$ \all{wäre}). \\
\all{Als ob} peut se réduire à \all{als} + verbe en 1\textsuperscript{re} position : \all{Sie tanzt, als würde sie schweben.}
\item \textbf{Temps du Konjunktiv II :} Présent $\rightarrow$ \all{würde} + infinitif (ou \all{wäre, hätte, könnte}...) ; Passé $\rightarrow$ \all{hätte/wäre} + Participe II : \all{als ob er gesungen hätte}.
\end{enumerate}
\end{propriete}

\begin{attention}[Pièges majeurs pour francophones]
\begin{enumerate}
\item \textbf{Ne JAMAIS traduire « comme » par \all{wie} dans une comparaison irréelle.} \all{wie} = égalité réelle. « Comme si » = \all{als ob} + KII.
\item \textbf{Verbe à la fin :} \all{als ob sie fliegen würden} (pas \all{sie würden fliegen}).
\item \textbf{Comparatif :} \all{wichtig} $\rightarrow$ \all{wichtiger} (pas \all{mehr wichtig}). Adjectifs en \all{-el, -er, -en} perdent souvent le \all{e} : \all{dunkel} $\rightarrow$ \all{dunkler}, \all{teuer} $\rightarrow$ \all{teurer}.
\item \all{als} vs \all{wie} : \all{Er ist größer als ich} (supériorité) ; \all{Er ist so groß wie ich} (égalité).
\end{enumerate}
\end{attention}

\begin{exoresolu}[Comparaison — Transformations et thème]
\textbf{A.} Reliez les deux phrases par une comparaison d'égalité (\all{so ... wie}) : \\
\all{Der Sänger ist laut. Die Trommel ist laut.}

\textbf{B.} Reliez par une comparaison de supériorité : \\
\all{Das Konzert ist interessant. Der Film ist interessant.}

\textbf{C.} Reliez par \all{als ob} + Konjunktiv II (présent) : \\
\all{Die Tänzerin schwebt. Sie ist ein Vogel.}

\textbf{D.} Traduisez : « Le tam-tam est plus important que la guitare à cette fête. » \\
\textbf{E.} Traduisez : « Elle danse comme si elle était une princesse. »

\tcblower
\textbf{A.} \all{Der Sänger ist so laut wie die Trommel.} (Égalité, \all{wie} + nominatif).

\textbf{B.} \all{Das Konzert ist interessanter als der Film.} (Comparatif \all{interessanter} + \all{als}).

\textbf{C.} \all{Die Tänzerin schwebt, als ob sie ein Vogel wäre.} (Irréel : \all{sein} $\rightarrow$ \all{wäre} à la fin). Variante : \all{... als wäre sie ein Vogel.}

\textbf{D.} \all{Die Trommel ist bei diesem Fest wichtiger als die Gitarre.} (\all{wichtiger als} ; \all{bei} + datif : \all{diesem Fest}).

\textbf{E.} \all{Sie tanzt, als ob sie eine Prinzessin wäre.} (\all{als ob} + KII \all{wäre} à la fin ; \all{eine Prinzessin} accusatif après \all{sein}).
\end{exoresolu}

\begin{exercice}[Comparaison — Entraînement]
1. Mettez l'adjectif au comparatif et complétez : \all{Die Musik ist ... (schön) das Gedicht.} \\
2. Transformez avec \all{als ob} + KII : \all{Er singt. Er ist ein Profi.} \\
3. Choisissez \all{wie} ou \all{als} : \all{Sie tanzt ... eine Profi-Tänzerin.} (égalité) \quad \all{Das Lied ist lauter ... das andere.} (supériorité) \\
4. Thème : « Si j'avais du temps, j'apprendrais la danse. » (Indice : \all{wenn} + KII).
\end{exercice}

\begin{corrige}[Exercice Comparaison]
1. \all{schöner als}
2. \all{Er singt, als ob er ein Profi wäre.} / \all{... als wäre er ein Profi.}
3. \all{wie} / \all{als}
4. \all{Wenn ich Zeit hätte, würde ich tanzen lernen.}
\end{corrige}

\courssub{Grammaire 2 : La déclinaison de l'adjectif}

\begin{propriete}[Les trois déclinaisons de l'adjectif épithète — Méthode « Article ? Cas ? Genre ? »]
\begin{enumerate}
\item \textbf{Déclinaison FAIBLE (après article défini \all{der, die, das, die} / \all{dieser, jener, welcher} / pronom relatif) :} l'article montre le cas/genre $\rightarrow$ adjectif en \textbf{-e} (Nominatif/Accusatif sg. masc./neut./fém.) ou \textbf{-en} (tous les autres cas + pluriel). \\
\all{der gute Sänger, die gute Sängerin, das gute Lied, die guten Lieder} ; \all{dem guten Sänger, den guten Sänger}.
\item \textbf{Déclinaison MIXTE (après article indéfini \all{ein, keine, mein, dein, sein, ihr, unser, euer, Ihr} / possessifs) :} l'adjectif \textbf{prend la marque du genre/cas là où l'article ne la montre pas} (Nominatif masc./neut., Accusatif neut. : \all{-er, -es} ; sinon \all{-e} ou \all{-en}). \\
\all{ein guter Sänger, ein gutes Lied, eine gute Sängerin} ; \all{meinem guten Freund} (datif).
\item \textbf{Déclinaison FORTE (SANS article) :} l'adjectif porte \textbf{toutes les marques} du cas/genre (comme l'article défini). \\
\all{guter Wein, gute Musik, laute Trommeln} (Nominatif) ; \all{guten Wein, gute Musik, laute Trommeln} (Accusatif) ; \all{gutem Wein, guter Musik, lauten Trommeln} (Datif) ; \all{guten Weines, guter Musik, lauter Trommeln} (Génitif).
\item \textbf{Règle d'or Datif / Génitif :} presque toujours \textbf{-en} (faible, mixte, forte au pluriel/datif). \all{mit dem berühmten Sänger, wegen des lauten Konzerts, mit lauten Trommeln}.
\end{enumerate}
\end{propriete}

\begin{exemplebox}[Tableaux de déclinaison — Adjectif \all{gut} + \all{Sänger / Lied / Musik / Trommeln}]
\begin{center}
\textbf{Faible (der / die / das / die)}
\begin{tabular}{|l|l|l|l|l|}
\hline
\rowcolor{popblueL} Cas & Masc. & Neutre & Féminin & Pluriel \\
\hline
Nom. & \all{der gute Sänger} & \all{das gute Lied} & \all{die gute Musik} & \all{die guten Trommeln} \\
Akk. & \all{den guten Sänger} & \all{das gute Lied} & \all{die gute Musik} & \all{die guten Trommeln} \\
Dat. & \all{dem guten Sänger} & \all{dem guten Lied} & \all{der guten Musik} & \all{den guten Trommeln} \\
Gen. & \all{des guten Sängers} & \all{des guten Liedes} & \all{der guten Musik} & \all{der guten Trommeln} \\
\hline
\end{tabular}
\end{center}
\vspace{0.3cm}
\begin{center}
\textbf{Mixte (ein / mein / kein ...)}
\begin{tabular}{|l|l|l|l|l|}
\hline
\rowcolor{poporangeL} Cas & Masc. & Neutre & Féminin & Pluriel \\
\hline
Nom. & \all{ein guter Sänger} & \all{ein gutes Lied} & \all{eine gute Musik} & \all{meine guten Trommeln} \\
Akk. & \all{einen guten Sänger} & \all{ein gutes Lied} & \all{eine gute Musik} & \all{meine guten Trommeln} \\
Dat. & \all{einem guten Sänger} & \all{einem guten Lied} & \all{einer guten Musik} & \all{meinen guten Trommeln} \\
Gen. & \all{eines guten Sängers} & \all{eines guten Liedes} & \all{einer guten Musik} & \all{meiner guten Trommeln} \\
\hline
\end{tabular}
\end{center}
\vspace{0.3cm}
\begin{center}
\textbf{Forte (sans article)}
\begin{tabular}{|l|l|l|l|l|}
\hline
\rowcolor{popgreenL} Cas & Masc. & Neutre & Féminin & Pluriel \\
\hline
Nom. & \all{guter Sänger} & \all{gutes Lied} & \all{gute Musik} & \all{gute Trommeln} \\
Akk. & \all{guten Sänger} & \all{gutes Lied} & \all{gute Musik} & \all{gute Trommeln} \\
Dat. & \all{gutem Sänger} & \all{gutem Lied} & \all{guter Musik} & \all{guten Trommeln} \\
Gen. & \all{guten Sängers} & \all{guten Liedes} & \all{guter Musik} & \all{guter Trommeln} \\
\hline
\end{tabular}
\end{center}
\end{exemplebox}

\begin{exoresolu}[Déclinaison — Complétion]
Mettez l'adjectif entre parenthèses à la forme correcte.
1. \all{Der (berühmt) ... Sänger tritt heute auf.} \\
2. \all{Wir besuchen ein (groß) ... Konzert.} \\
3. \all{Sie tanzt mit dem (jung) ... Tänzer.} \\
4. \all{(Traditionell) ... Musik ist in Kamerun sehr lebendig.} \\
5. \all{Er spielt auf (alt) ... Trommeln.} (sans article, accusatif pluriel) \\
6. \all{Wegen (laut) ... Musik können wir nicht schlafen.} (génitif, \all{wegen} + génitif)

\tcblower
\textbf{1.} \all{der} + Nom. masc. $\rightarrow$ faible \textbf{-e} : \all{Der berühmte Sänger tritt heute auf.} \\
\textbf{2.} \all{ein} + Akk. neutre (\all{das Konzert}) : article ne montre pas genre $\rightarrow$ mixte \textbf{-es} : \all{Wir besuchen ein großes Konzert.} \\
\textbf{3.} \all{mit} + Dat. masc. (\all{der Tänzer}) $\rightarrow$ toujours \textbf{-en} : \all{Sie tanzt mit dem jungen Tänzer.} \\
\textbf{4.} Sans article, Nom. fém. (\all{die Musik}) $\rightarrow$ forte \textbf{-e} : \all{Traditionelle Musik ist in Kamerun sehr lebendig.} \\
\textbf{5.} Sans article, Akk. pluriel $\rightarrow$ forte \textbf{-e} : \all{Er spielt auf alte Trommeln.} \\
\textbf{6.} \all{wegen} + Gén. fém. (\all{die Musik}) $\rightarrow$ forte \textbf{-er} : \all{Wegen lauter Musik können wir nicht schlafen.}
\end{exoresolu}

\begin{exercice}[Déclinaison — Entraînement]
Complétez avec la bonne terminaison.
1. \all{Ich höre (gern) ... Musik.} (accusatif, sans article? Non, \all{gern} est adverbe. Phrase : \all{Ich höre gern gute Musik.} $\rightarrow$ \all{gute}).
2. \all{Das ist das (neu) ... Lied der (berühmt) ... Sängerin.} \\
3. \all{Wir sprechen mit (jung) ... Leuten.} (datif pluriel, sans article) \\
4. \all{Er kauft (ein) ... (teuer) ... Instrument.} (accusatif neutre, article indéfini)
\end{exercice}

\begin{corrige}[Exercice Déclinaison]
1. \all{gute} (acc. fém. sans article $\rightarrow$ forte -e) \\
2. \all{neue} (acc. neutre faible -e) ; \all{berühmten} (génitif fém. faible -en) \\
3. \all{jungen} (dat. pluriel forte -en) \\
4. \all{ein} \all{teures} (acc. neutre mixte -es)
\end{corrige}

\courssub{Grammaire 3 : Le subjonctif II (Konjunktiv II)}

\begin{propriete}[Formation et emplois du Konjunktiv II — Niveau Terminale]
\begin{enumerate}
\item \textbf{Forme générale (tous verbes) :} \all{würde} (conjugué) + \textbf{infinitif} à la fin. \\
\all{Ich würde gern tanzen. / Wir würden das Konzert besuchen.}
\item \textbf{Formes irrégulières (à connaître par cœur) — Présent du KII :}
\begin{center}
\begin{tabular}{|l|l|l|l|l|}
\hline
\rowcolor{popblueL} Infinitif & \all{ich} & \all{du} & \all{er/sie/es} & \all{wir / ihr / sie/Sie} \\
\hline
\all{sein} & \all{wäre} & \all{wärest} & \all{wäre} & \all{wären / wäret / wären} \\
\all{haben} & \all{hätte} & \all{hättest} & \all{hätte} & \all{hätten / hättet / hätten} \\
\all{werden} & \all{würde} & \all{würdest} & \all{würde} & \all{würden / würdet / würden} \\
\all{können} & \all{könnte} & \all{könntest} & \all{könnte} & \all{könnten / könntet / könnten} \\
\all{müssen} & \all{müsste} & \all{müsstest} & \all{müsste} & \all{müssten / müsstet / müssten} \\
\all{dürfen} & \all{dürfte} & \all{dürftest} & \all{dürfte} & \all{dürften / dürftet / dürften} \\
\all{sollen} & \all{sollte} & \all{solltest} & \all{sollte} & \all{sollten / solltet / sollten} \\
\all{wissen} & \all{wüsste} & \all{wüsstest} & \all{wüsste} & \all{wüssten / wüsstet / wüssten} \\
\all{geben} & \all{gäbe} & \all{gäbest} & \all{gäbe} & \all{gäben / gäbet / gäben} \\
\hline
\end{tabular}
\end{center}
\item \textbf{Passé du KII (Perfekt du Konjunktiv) :} \all{hätte / wäre} (KII) + \textbf{Participe II}. \\
\all{Ich wäre gern gekommen.} (verbe de mouvement $\rightarrow$ \all{sein}) ; \all{Ich hätte das Lied gesungen.} (autres $\rightarrow$ \all{haben}).
\item \textbf{Emplois principaux :}
\begin{itemize}
\item Souhait / désir : \all{Ich wäre jetzt gern auf dem Fest.} / \all{Ich hätte gern eine Trommel.}
\item Conseil / recommandation : \all{Du solltest mehr üben.} / \all{Du könntest die Sängerin hören.}
\item Hypothèse irréelle (si-clause) : \all{Wenn ich Zeit hätte, würde ich tanzen lernen.} (\all{wenn} + KII à la fin ; principale \all{würde} + inf.).
\item Comparaison irréelle : \all{als ob} + KII (voir Grammaire 1).
\item Politesse : \all{Ich hätte eine Frage.} / \all{Würden Sie mir helfen?}
\end{itemize}
\item \textbf{Place du verbe :} Subordonnée (\all{wenn, als ob, dass}) $\rightarrow$ verbe conjugué \textbf{à la fin}. Principale $\rightarrow$ verbe en \textbf{2\textsuperscript{e} position}.
\end{enumerate}
\end{propriete}

\begin{exoresolu}[Konjunktiv II — Thème et transformation]
\textbf{A.} Traduisez : « Si j'avais de l'argent, j'achèterais un tam-tam. » \\
\textbf{B.} Traduisez : « Tu devrais écouter cette chanteuse ! » \\
\textbf{C.} Traduisez : « J'aimerais danser comme une professionnelle. » \\
\textbf{D.} Mettez au KII passé : \all{Ich komme zum Fest.} / \all{Er singt das Lied.}

\tcblower
\textbf{A.} \all{Wenn ich Geld hätte, würde ich eine Trommel kaufen.} (\all{wenn} + \all{hätte} fin subordonnée ; principale \all{würde ... kaufen} fin). Variante inversée : \all{Hätte ich Geld, so würde ich eine Trommel kaufen.}

\textbf{B.} \all{Du solltest diese Sängerin hören!} (accusatif) ou \all{Du solltest dieser Sängerin zuhören!} (datif + \all{zuhören}). \all{sollen} KII $\rightarrow$ \all{solltest}.

\textbf{C.} \all{Ich würde gern wie eine Profi-Tänzerin tanzen.} (\all{wie} comparaison égalité ; \all{würde gern} + inf.).

\textbf{D.} \all{Ich wäre zum Fest gekommen.} (\all{kommen} mouvement $\rightarrow$ \all{sein}) ; \all{Er hätte das Lied gesungen.} (\all{singen} $\rightarrow$ \all{haben}).
\end{exoresolu}

\begin{exercice}[Konjunktiv II — Entraînement]
1. Conjuguez au KII présent : \all{ich (können)} ... \quad \all{wir (haben)} ... \quad \all{er (sein)} ... \\
2. Complétez la phrase de souhait : \all{Ich ... (würde) gern nach Deutschland ... (reisen).} \\
3. Conseil : \all{Du ... (sollen) weniger ... (telefonieren) und mehr ... (üben).} \\
4. Hypothèse : \all{Wenn er ... (haben) Zeit, ... (er / tanzen) ... (lernen).}
\end{exercice}

\begin{corrige}[Exercice Konjunktiv II]
1. \all{könnte, hätten, wäre}
2. \all{würde ... reisen}
3. \all{solltest ... telefonieren ... üben}
4. \all{hätte, würde er tanzen lernen} (ordre : \all{Wenn er Zeit hätte, würde er tanzen lernen.})
\end{corrige}

\courssub{Méthode du commentaire et production guidée}

\begin{methode}[Rédiger un résumé (Zusammenfassung) ou une courte production — Étape par étape]
\begin{enumerate}
\item \textbf{Résumé (4--6 phrases) :} Repère 3--4 idées forces du texte. Reformule avec TES mots (pas de copier-coller). Début type : \all{Der Text handelt von ...} / \all{In dem Text geht es um ...}. Pas d'opinion personnelle. Temps : présent ou prétérit.
\item \textbf{Production guidée (paragraphe 5--6 phrases) :} Respecte la consigne (sujet, nombre de mots). Utilise le vocabulaire du thème (\all{Konzert, Trommel, tanzen, Fest, Sänger}). Varie les connecteurs : \all{zuerst, dann, außerdem, deshalb, aber, schließlich}. Structure : Introduction $\rightarrow$ Développement (exemples Cameroun/Allemagne) $\rightarrow$ Comparaison \all{als ob} $\rightarrow$ Souhait KII $\rightarrow$ Conclusion (\all{Meiner Meinung nach ...}).
\item \textbf{Check-list relecture (obligatoire) :} Verbe 2\textsuperscript{e} pos. (principal) / fin (subordonnée) ? Majuscules à TOUS les noms ? Adjectifs accordés (déclinaison) ? Particule séparable en fin (\all{tritt auf}) ? Cas après prépositions (\all{mit} + dat, \all{wegen} + gen, \all{zu} + dat) ? Umlaute (\all{ä, ö, ü}) et \all{ß} corrects ?
\end{enumerate}
\end{methode}

\begin{experience}[Activité orale : Jeu de rôle — « Interview nach dem Festival »]
\textbf{Consigne :} Par binômes. A = journaliste (allemand), B = musicien/chanteur camerounais de retour d'un festival en Allemagne.
\textbf{Questions types pour A :} \all{Wie war das Festival in Berlin? Haben Sie deutsche Bands gehört? Was ist der Unterschied zwischen kamerunischer und deutscher Musik? Würden Sie gern wieder kommen?}
\textbf{Réponses types pour B :} Utiliser \all{Perfekt} (récit), \all{als ob} (comparaison), \all{Konjunktiv II} (souhait/conseil). Ex. : \all{Es war toll, als ob ich zu Hause wäre. Ich würde gern nächstes Jahr wiederkommen.}
\textbf{Objectif :} Interaction spontanée, structures du programme, prononciation (Umlaute, \all{ß}, \all{ch}, \all{pf}).
\end{experience}

\begin{savaistu}[Culture : Musique camerounaise et allemande — Ponts et contrastes]
\begin{itemize}
\item \textbf{Cameroun :} \all{Makossa} (Douala, années 70, \all{Manu Dibango}), \all{Bikutsi} (région Centre, \all{Anne-Marie Nzié}, \all{Les Têtes Brûlées}), \all{Assiko}, \all{Bend Skin}. Instruments : \all{die Trommel (Tam-tam), die Balafon, die Marimba, die Gitarre}. Danse : mouvements rapides, épaules, bassin. Fonction sociale : rites, fêtes, contestation.
\item \textbf{Allemagne :} Classique (\all{Bach, Beethoven, Brahms}), \all{Schlager} (populaire), \all{Krautrock} (années 70, \all{Kraftwerk, Can}), \all{Techno} (Berlin, \all{Berghain}), \all{Hip-Hop} (\all{Cro, Sido}). Fêtes : \all{Oktoberfest} (musique traditionnelle bavaroise), \all{Karneval} (Cologne, \all{Bläck Fööss}), Festivals : \all{Rock am Ring, Fusion, Melt!}.
\item \textbf{Échanges :} Tournées d'artistes camerounais en Allemagne (diaspora, festivals \all{Africa Festival Würzburg}). Collaborations (ex. \all{Manu Dibango} avec musiciens allemands). Apprentissage mutuel : rythmes africains dans la pop allemande, technologie studio en Afrique.
\end{itemize}
\end{savaistu}

\begin{exoresolu}[Production écrite guidée — Type Bac : « Musik und Tanz bei einem Fest in deiner Stadt »]
\textbf{Consigne :} Rédigez un paragraphe de 5--6 phrases. Utilisez au moins : 1 comparaison \all{als ob} + KII, 1 phrase au KII (souhait/conseil), 3 mots du lexique thème, 1 adjectif décliné correctement.

\tcblower
\textbf{Plan :} Introduction (fête locale) $\rightarrow$ Description (musiciens, danseurs, instruments) $\rightarrow$ Comparaison irréelle $\rightarrow$ Souhait KII $\rightarrow$ Conclusion (valeur sociale).

\textbf{Modèle de réponse :}

\all{Letztes Wochenende haben wir in Yaoundé ein großes Dorffest gefeiert. Die Musiker haben auf den Trommeln und Balafons gespielt, und die berühmte Sängerin des Viertels hat traditionelle Lieder gesungen. Die Tänzer haben sich so schnell bewegt, als ob sie fliegen würden. Alle Gäste haben bis Mitternacht getanzt und gelacht. Ich würde gern jedes Jahr an so einem Fest teilnehmen. Meiner Meinung nach verbinden Musik und Tanz die Generationen.}

\textbf{Analyse grammaticale (points vérifiés) :}
\begin{itemize}
\item \all{gefeiert, gespielt, gesungen, bewegt, getanzt, gelacht} : Perfekt (auxiliaire \all{haben} / \all{sein} pour \all{bewegen} ? \all{sich bewegen} $\rightarrow$ \all{haben} ici car pas de déplacement de lieu, juste mouvement sur place. \all{haben sich bewegt} correct).
\item \all{als ob sie fliegen würden} : \all{als ob} + KII \all{würden fliegen} (verbe à la fin). \checkmark
\item \all{Ich würde gern ... teilnehmen} : KII souhait, particule \all{teil-} séparée \all{nehmen} $\rightarrow$ \all{teilnehmen} à la fin. \checkmark
\item \all{ein großes Fest} (acc. neutre mixte \all{-es}), \all{die berühmte Sängerin} (nom. fém. faible \all{-e}), \all{traditionelle Lieder} (acc. pluriel faible \all{-e}), \all{so einem Fest} (dat. neutre mixte \all{-em}). \checkmark
\item Connecteurs : \all{und, auch, deshalb (implicite), Meiner Meinung nach}. \checkmark
\end{itemize}
\end{exoresolu}

\begin{exercice}[Production écrite — À vous de jouer]
Rédigez un court paragraphe (5--6 phrases) sur le sujet : \all{Ein Konzertbesuch in Deutschland.} Contraintes : Perfekt (récit), 1 \all{als ob} + KII, 1 KII souhait/conseil, vocabulaire thème (min. 4 mots), 2 adjectifs déclinés.
\end{exercice}

\begin{corrige}[Exercice Production — Proposition de correction]
\all{Vor zwei Wochen habe ich in Berlin ein großes Open-Air-Konzert besucht. Die Band hat moderne und alte Lieder gespielt, und die Sängerin hat eine wunderbare Stimme gehabt. Das Publikum hat so laut gejubelt, als ob es ein Fußballfinale wäre. Ich würde gern nächstes Jahr wieder hingehen. Meiner Meinung nach ist Live-Musik viel besser als Musik vom Handy.}
\end{corrige}

\begin{attention}[Les 5 erreurs qui coûtent des points au Bac — Synthèse finale]
\begin{enumerate}
\item \textbf{Majuscules aux noms :} \all{die trommel, das konzert, der sänger} = \textbf{FAUX}. Toujours : \all{die Trommel, das Konzert, der Sänger, das Fest, die Musik, der Tanz}. \cle{MaJUSCULES}
\item \textbf{\all{als} vs \all{wie} vs \all{als ob} :} « plus ... que » = comparatif + \all{als} (\all{lauter als}) ; « aussi ... que » = \all{so ... wie} ; « comme si » = \all{als ob} + \textbf{Konjunktiv II verbe à la fin}. JAMAIS \all{wie} pour « comme si ».
\item \textbf{Place du verbe :} Subordonnée (\all{weil, dass, wenn, als ob, obwoh}) $\rightarrow$ verbe conjugué \textbf{à la fin}. \all{..., als ob sie fliegen würden}. Principale $\rightarrow$ verbe \textbf{2\textsuperscript{e} position}.
\item \textbf{Déclinaison adjectif :} \all{mit dem jung Tänzer} = \textbf{FAUX}. \all{mit dem jungen Tänzer} (datif $\rightarrow$ \textbf{-en}). Réflexe : \textbf{Datif / Génitif / Pluriel = presque toujours -en}.
\item \textbf{Calques du français (faux amis / structures) :}
\begin{itemize}
\item « Écouter de la musique » $\neq$ \all{hören Musik von} $\rightarrow$ \all{Musik hören} (accusatif).
\item « Assister à un concert » $\neq$ \all{assistieren ein Konzert} $\rightarrow$ \all{ein Konzert besuchen}.
\item « Jouer du tam-tam » $\rightarrow$ \all{Trommel spielen} (accusatif) ou \all{auf der Trommel spielen} (datif).
\item Verbes à particule : \all{Die Band tritt um 20 Uhr auf.} (particule \all{auf} en \textbf{fin de phrase}).
\item « Comme » (comparaison) $\rightarrow$ \all{wie} (égalité) ou \all{als} (supériorité) ou \all{als ob} (irréel). Pas de \all{comme} unique.
\end{itemize}
\end{enumerate}
\end{attention}

\newpage

\coursec{Exercices}

\courssub{Je vérifie mes connaissances}

\begin{exercice}[QCM : vocabulaire et grammaire]
Entourez la bonne réponse.
\begin{enumerate}
\item \all{die Trommel} signifie : a) la trompette \quad b) le tambour / la tam-tam \quad c) la danse
\item Le pluriel de \all{das Konzert} est : a) \all{die Konzerten} \quad b) \all{die Konzerte} \quad c) \all{die Konzerts}
\item \all{Er singt, \_\_\_ er ein Profi wäre.} a) \all{wie} \quad b) \all{als ob} \quad c) \all{als}
\item \all{Sie tanzt wie ein\_\_\_ professionelle Tänzerin.} a) \all{-e} \quad b) \all{-er} \quad c) \all{-en}
\item Le Konjunktiv II de \all{ich kann} est : a) \all{ich kannte} \quad b) \all{ich könnte} \quad c) \all{ich würde können}
\end{enumerate}
\end{exercice}

\begin{exercice}[Vrai ou faux ? Justifiez]
Répondez par \all{Richtig} ou \all{Falsch} et justifiez en une phrase en français.
\begin{enumerate}
\item \all{Der Sänger} est un nom féminin.
\item Après \all{als ob}, le verbe conjugué se place en fin de proposition.
\item \all{Wie} exprime une comparaison réelle, \all{als ob} une comparaison irréelle.
\item Après un article défini (\all{der, die, das}), l'adjectif prend toujours la terminaison \all{-e}.
\end{enumerate}
\end{exercice}

\begin{exercice}[Vocabulaire : articles et pluriels]
Complétez le tableau avec l'article et le pluriel de chaque nom.
\begin{center}
\begin{tabular}{|l|l|l|}
\hline
\rowcolor{popblueL} Nom & Article & Pluriel \\
\hline
Musik & \ldots & \ldots \\
\hline
Tanz & \ldots & \ldots \\
\hline
Fest & \ldots & \ldots \\
\hline
Sänger & \ldots & \ldots \\
\hline
Lied & \ldots & \ldots \\
\hline
Trommel & \ldots & \ldots \\
\hline
\end{tabular}
\end{center}
\end{exercice}

\begin{exercice}[Conjugaison : Konjunktiv II]
Conjuguez les verbes entre parenthèses au Konjunktiv II (Präteritum).
\begin{enumerate}
\item Wenn ich Zeit \ldots\ (haben), würde ich zum Konzert gehen.
\item \ldots\ du gern Sänger? (sein)
\item Wir \ldots\ das Fest besuchen, wenn es nicht regnete. (können)
\item Wenn das Wetter schön \ldots\ (sein), würden wir draußen tanzen.
\item Er \ldots\ gern nach Deutschland reisen. (mögen)
\end{enumerate}
\end{exercice}

\courssub{Je m'entraîne}

\begin{exercice}[Comparaison : \all{wie} ou \all{als ob} ?]
Complétez avec \all{wie} ou \all{als ob} et mettez le verbe à la forme correcte (Konjunktiv II après \all{als ob}).
\begin{enumerate}
\item Der Tänzer bewegt sich, \ldots\ er fliegen \ldots\ (können).
\item Die Sängerin singt \ldots\ ein Vogel.
\item Er tut so, \ldots\ er alles \ldots\ (wissen).
\item Das Mädchen tanzt \ldots\ ihre Mutter.
\item Sie schreit, \ldots\ sie verrückt \ldots\ (sein).
\item In Kamerun liebt man die Musik \ldots\ in Deutschland.
\end{enumerate}
\end{exercice}

\begin{exercice}[Déclinaison de l'adjectif]
Complétez avec la bonne terminaison de l'adjectif.
\begin{enumerate}
\item Das \ldots\ (jung) Mädchen singt ein \ldots\ (traditionell) Lied mit \ldots\ (laut) Stimme.
\item Der \ldots\ (berühmt) Sänger gibt ein \ldots\ (groß) Konzert in der \ldots\ (alt) Stadt.
\item Wir besuchen das \ldots\ (jährlich) Fest mit unseren \ldots\ (gut) Freunden.
\item \ldots\ (Deutsch) Musik ist bei den \ldots\ (jung) Leuten sehr beliebt.
\end{enumerate}
\end{exercice}

\begin{exercice}[Transformation : de l'indicatif au Konjunktiv II]
Transformez les phrases selon le modèle : \all{Ich habe keine Zeit. Ich besuche das Konzert nicht.} $\rightarrow$ \all{Wenn ich Zeit hätte, würde ich das Konzert besuchen.}
\begin{enumerate}
\item Ich habe kein Geld. Ich kaufe die Trommel nicht.
\item Wir haben kein Auto. Wir fahren nicht zum Fest.
\item Er ist krank. Er singt nicht im Chor.
\item Ihr seid müde. Ihr tanzt nicht.
\end{enumerate}
\end{exercice}

\begin{exercice}[Appariement et texte à trous]
\textbf{A.} Reliez chaque mot allemand à sa traduction française.
\begin{center}
\begin{tabular}{|l|l|}
\hline
\rowcolor{popblueL} Deutsch & Français \\
\hline
1. \all{die Bühne} & a) le public \\
\hline
2. \all{das Publikum} & b) la scène \\
\hline
3. \all{der Rhythmus} & c) le rythme \\
\hline
4. \all{auftreten} & d) applaudir \\
\hline
5. \all{klatschen} & e) se produire (sur scène) \\
\hline
\end{tabular}
\end{center}
\textbf{B.} Complétez le texte avec les mots suivants : \all{Konzert, tanzten, Sängerin, Trommeln, Publikum}.
\par
Gestern war ein großes \ldots\ in Yaoundé. Eine bekannte \ldots\ sang traditionelle Lieder. Die Musiker spielten auf den \ldots, und viele Leute \ldots\ bis Mitternacht. Das \ldots\ war begeistert und klatschte laut.
\end{exercice}

\courssub{Je me prépare au Bac}

\begin{exercice}[Compréhension de texte type Bac]
Lisez le texte suivant, puis répondez aux questions.

\begin{textebox}[Ein Konzert in Berlin]
Am Samstagabend besuchte ich mit meinen Freunden ein großes Konzert in Berlin. Die berühmte Sängerin Lena trat auf einer riesigen Bühne im Zentrum der Stadt auf. Tausende Menschen waren gekommen: junge und alte Leute, Deutsche und Ausländer. Als die Musik begann, tanzte das ganze Publikum, als ob es keinen Morgen gäbe. Die alte Frau neben mir bewegte sich wie eine junge Tänzerin. Die Musiker spielten auf modernen Instrumenten, aber auch auf traditionellen Trommeln. Nach zwei Stunden war das Konzert zu Ende, aber niemand wollte nach Hause gehen. „Noch ein Lied!“, rief das begeisterte Publikum. Die Sängerin lächelte und sang ein letztes traditionelles Lied mit leiser Stimme. Es war ein unvergesslicher Abend. Wenn ich mehr Geld hätte, würde ich jedes Wochenende zu solchen Konzerten gehen.
\end{textebox}

\textbf{Glossaire} : \all{auftreten} : se produire ; \all{riesig} : immense ; \all{begeistert} : enthousiaste ; \all{unvergesslich} : inoubliable.

\textbf{A. Questions de compréhension} (répondez en allemand, par des phrases complètes) :
\begin{enumerate}
\item Wann besuchte der Erzähler das Konzert und mit wem ?
\item Wer war auf dem Konzert anwesend ?
\item Wie tanzte das Publikum ?
\item Was rief das Publikum am Ende ? Warum ?
\item Was würde der Erzähler machen, wenn er mehr Geld hätte ?
\end{enumerate}
\textbf{B.} Répondez par \all{Richtig} ou \all{Falsch} et justifiez par une citation du texte :
\begin{enumerate}
\item Das Konzert fand in einem kleinen Dorf statt.
\item Die Musiker benutzten nur moderne Instrumente.
\item Nach dem Konzert wollten alle sofort nach Hause gehen.
\end{enumerate}
\textbf{C. Questions de langue} :
\begin{enumerate}
\item Relevez dans le texte une comparaison avec \all{als ob} et une avec \all{wie}. Expliquez la différence.
\item Relevez deux adjectifs déclinés et précisez le cas et le genre.
\item Mettez au Konjunktiv II : \all{Die Sängerin lächelt und singt ein letztes Lied.}
\end{enumerate}
\end{exercice}

\begin{exercice}[Thème et version]
\textbf{A. Version.} Traduisez en français le paragraphe suivant.

\begin{textebox}[Das Oktoberfest]
Das Oktoberfest in München ist das größte Volksfest der Welt. Jedes Jahr kommen Millionen Besucher aus allen Ländern. Die fröhlichen Menschen tragen traditionelle bayerische Kleidung und tanzen zu lauter Musik. Manche Besucher tanzen auf den Bänken, als ob sie professionelle Tänzer wären. Ein berühmter deutscher Schriftsteller sagte einmal, das Fest sei wie ein großer Traum. Wenn ich dort wäre, würde ich die ganze Nacht tanzen.
\end{textebox}

\textbf{B. Thème.} Traduisez en allemand les phrases suivantes (utilisez le lexique du chapitre, \all{als ob} et le Konjunktiv II) :
\begin{enumerate}
\item Hier à Douala, il y a eu une grande fête avec de la musique traditionnelle.
\item Les jeunes filles dansaient comme des professionnelles.
\item Le vieil homme jouait du tam-tam, comme s'il était jeune.
\item Si j'avais le temps, je visiterais toutes les fêtes de ma région.
\item Le public chantait et dansait jusqu'à minuit.
\end{enumerate}
\end{exercice}

\begin{exercice}[Production écrite type Bac]
\textbf{Sujet :} \all{Beschreiben Sie ein Fest in Ihrer Region: die Musik, die Tänze, die Menschen und Ihre Gefühle.}
\par
Rédigez un texte de \textbf{80 à 100 mots} en allemand. Vous utiliserez obligatoirement :
\begin{itemize}
\item au moins cinq mots du vocabulaire du chapitre (\all{die Musik, der Tanz, das Konzert, die Trommel, der Sänger, das Fest, das Publikum}...) ;
\item une comparaison avec \all{wie} et une comparaison irréelle avec \all{als ob} + Konjunktiv II ;
\item au moins deux adjectifs correctement déclinés ;
\item une phrase au Konjunktiv II avec \all{würden} ou une forme en \all{hätte/wäre}.
\end{itemize}
\textbf{Grille d'évaluation} : vocabulaire thématique (4 pts) ; grammaire : comparaison et Konjunktiv II (6 pts) ; déclinaison de l'adjectif et orthographe (4 pts) ; cohérence et respect du nombre de mots (6 pts). Total : 20 pts.
\end{exercice}

\newpage

\coursec{Corrigés des exercices}

\begin{corrige}[Exercice 1 — QCM : vocabulaire et grammaire]
\begin{enumerate}
\item \textbf{b)} \all{die Trommel} signifie le tambour / le tam-tam. Attention au faux ami : la trompette se dit \all{die Trompete}.
\item \textbf{b)} Le pluriel de \all{das Konzert} est \all{die Konzerte} (pluriel en \all{-e}, sans Umlaut).
\item \textbf{b)} \all{Er singt, als ob er ein Profi wäre.} — \all{als ob} introduit une comparaison irréelle et exige le Konjunktiv II (\all{wäre}) ; le verbe conjugué va en fin de proposition.
\item \textbf{a)} \all{Sie tanzt wie eine professionelle Tänzerin.} — après l'article indéfini \all{eine} au nominatif féminin, l'adjectif prend la terminaison \all{-e} (déclinaison mixte).
\item \textbf{b)} Le Konjunktiv II de \all{ich kann} est \all{ich könnte} (Präteritum \all{konnte} + Umlaut). \all{Ich kannte} est le Präteritum de \all{kennen} : ne pas confondre.
\end{enumerate}
\end{corrige}

\begin{corrige}[Exercice 2 — Vrai ou faux ? Justifiez]
\begin{enumerate}
\item \all{Falsch}. \all{Der Sänger} est un nom masculin (le chanteur) ; le féminin est \all{die Sängerin} (la chanteuse), avec le suffixe féminin \all{-in}.
\item \all{Richtig}. \all{Als ob} est une conjonction de subordination : le verbe conjugué se place en dernière position. Exemple : \all{Er tanzt, als ob er jung wäre.}
\item \all{Richtig}. \all{Wie} exprime une comparaison réelle (\all{Sie singt wie ein Vogel.}), tandis que \all{als ob} exprime une comparaison irréelle, imaginaire, et exige le Konjunktiv II (\all{Er tut so, als ob er alles wüsste.}).
\item \all{Falsch}. Après un article défini, l'adjectif prend \all{-e} seulement au nominatif singulier (et à l'accusatif féminin et neutre) : \all{der junge Mann, die junge Frau, das junge Mädchen}. Ailleurs, il prend \all{-en} : \all{den jungen Mann, dem jungen Mann, die jungen Leute}.
\end{enumerate}
\end{corrige}

\begin{corrige}[Exercice 3 — Vocabulaire : articles et pluriels]
\begin{center}
\begin{tabular}{|l|l|l|}
\hline
\rowcolor{popblueL} Nom & Article & Pluriel \\
\hline
Musik & die & die Musiken (rare ; en général sans pluriel) \\
\hline
Tanz & der & die Tänze \\
\hline
Fest & das & die Feste \\
\hline
Sänger & der & die Sänger \\
\hline
Lied & das & die Lieder \\
\hline
Trommel & die & die Trommeln \\
\hline
\end{tabular}
\end{center}
\textbf{Points de vigilance :} \all{der Tanz} prend l'Umlaut au pluriel (\all{Tänze}) ; \all{das Lied} a un pluriel en \all{-er} (\all{Lieder}), fréquent chez les noms neutres ; \all{der Sänger} est invariable au pluriel. Apprenez toujours chaque nom avec son article et son pluriel.
\end{corrige}

\begin{corrige}[Exercice 4 — Conjugaison : Konjunktiv II]
\begin{enumerate}
\item Wenn ich Zeit \textbf{hätte}, würde ich zum Konzert gehen. (\all{haben} : Präteritum \all{hatte} + Umlaut $\rightarrow$ \all{hätte})
\item \textbf{Wärst} du gern Sänger? (\all{sein} : \all{war} + Umlaut $\rightarrow$ \all{wärst}) — « Aimerais-tu être chanteur ? »
\item Wir \textbf{könnten} das Fest besuchen, wenn es nicht regnete. (\all{können} : \all{konnte} + Umlaut)
\item Wenn das Wetter schön \textbf{wäre}, würden wir draußen tanzen. (\all{sein} $\rightarrow$ \all{wäre})
\item Er \textbf{möchte} gern nach Deutschland reisen. (\all{mögen} : forme de Konjunktiv II \all{möchte}, très fréquente pour exprimer un souhait poli)
\end{enumerate}
\textbf{Rappel :} le Konjunktiv II se forme sur le radical du Präteritum ; les verbes forts prennent l'Umlaut (\all{war}$\rightarrow$\all{wäre}, \all{hatte}$\rightarrow$\all{hätte}, \all{kam}$\rightarrow$\all{käme}) et les terminaisons \all{-e, -est, -e, -en, -et, -en}.
\end{corrige}

\begin{corrige}[Exercice 5 — Comparaison : wie ou als ob ?]
\begin{enumerate}
\item Der Tänzer bewegt sich, \textbf{als ob} er fliegen \textbf{könnte}. — comparaison irréelle : Konjunktiv II de \all{können}, verbe en fin de proposition. « Le danseur se déplace comme s'il pouvait voler. »
\item Die Sängerin singt \textbf{wie} ein Vogel. — comparaison réelle, sans verbe conjugué. « La chanteuse chante comme un oiseau. »
\item Er tut so, \textbf{als ob} er alles \textbf{wüsste}. — irréel : Konjunktiv II de \all{wissen} (\all{wusste} + Umlaut $\rightarrow$ \all{wüsste}). « Il fait comme s'il savait tout. »
\item Das Mädchen tanzt \textbf{wie} ihre Mutter. — comparaison réelle entre deux personnes. « La fillette danse comme sa mère. »
\item Sie schreit, \textbf{als ob} sie verrückt \textbf{wäre}. — irréel : Konjunktiv II de \all{sein}. « Elle crie comme si elle était folle. »
\item In Kamerun liebt man die Musik \textbf{wie} in Deutschland. — comparaison réelle entre deux pays. « Au Cameroun, on aime la musique comme en Allemagne. »
\end{enumerate}
\textbf{Astuce :} si la comparaison est imaginaire ou exagérée, choisissez \all{als ob} + Konjunktiv II ; si elle est réelle et simple, choisissez \all{wie}.
\end{corrige}

\begin{corrige}[Exercice 6 — Déclinaison de l'adjectif]
\begin{enumerate}
\item Das \textbf{junge} Mädchen singt ein \textbf{traditionelles} Lied mit \textbf{lauter} Stimme. — \all{das junge} : nominatif neutre après article défini ($-e$) ; \all{ein traditionelles} : accusatif neutre après article indéfini ($-es$) ; \all{mit lauter Stimme} : datif féminin sans article ($-er$, déclinaison forte).
\item Der \textbf{berühmte} Sänger gibt ein \textbf{großes} Konzert in der \textbf{alten} Stadt. — \all{der berühmte} : nominatif masculin après \all{der} ($-e$) ; \all{ein großes} : accusatif neutre ($-es$) ; \all{in der alten Stadt} : datif féminin après article défini ($-en$).
\item Wir besuchen das \textbf{jährliche} Fest mit unseren \textbf{guten} Freunden. — \all{das jährliche} : accusatif neutre après \all{das} ($-e$) ; \all{mit unseren guten Freunden} : datif pluriel ($-en$ partout).
\item \textbf{Deutsche} Musik ist bei den \textbf{jungen} Leuten sehr beliebt. — \all{deutsche Musik} : nominatif féminin sans article ($-e$, déclinaison forte) ; \all{bei den jungen Leuten} : datif pluriel ($-en$).
\end{enumerate}
\textbf{Rappel :} après article défini, seulement deux terminaisons : \all{-e} ou \all{-en} ; au datif et au génitif, toujours \all{-en}.
\end{corrige}

\begin{corrige}[Exercice 7 — Transformation : de l'indicatif au Konjunktiv II]
\begin{enumerate}
\item \all{Wenn ich Geld hätte, würde ich die Trommel kaufen.} — « Si j'avais de l'argent, j'achèterais le tam-tam. »
\item \all{Wenn wir ein Auto hätten, würden wir zum Fest fahren.} — « Si nous avions une voiture, nous irions à la fête en voiture. »
\item \all{Wenn er nicht krank wäre, würde er im Chor singen.} — « S'il n'était pas malade, il chanterait dans le chœur. » (on peut aussi écrire : \all{Wenn er gesund wäre, ...})
\item \all{Wenn ihr nicht müde wärt, würdet ihr tanzen.} — « Si vous n'étiez pas fatigués, vous danseriez. »
\end{enumerate}
\textbf{Points de vigilance :} dans la subordonnée introduite par \all{wenn}, le verbe conjugué va en fin de phrase ; dans la principale qui suit, le verbe (\all{würde}) est en première position et le sujet juste après (\all{würde ich...}). N'oubliez pas la négation quand la phrase de départ est négative (\all{nicht krank wäre}).
\end{corrige}

\begin{corrige}[Exercice 8 — Appariement et texte à trous]
\textbf{A. Appariement :}
\begin{center}
\begin{tabular}{|l|l|}
\hline
\rowcolor{popblueL} Deutsch & Français \\
\hline
1. \all{die Bühne} & b) la scène \\
\hline
2. \all{das Publikum} & a) le public \\
\hline
3. \all{der Rhythmus} & c) le rythme \\
\hline
4. \all{auftreten} & e) se produire (sur scène) \\
\hline
5. \all{klatschen} & d) applaudir \\
\hline
\end{tabular}
\end{center}
\textbf{B. Texte complété :}
\par
\all{Gestern war ein großes \textbf{Konzert} in Yaoundé. Eine bekannte \textbf{Sängerin} sang traditionelle Lieder. Die Musiker spielten auf den \textbf{Trommeln}, und viele Leute \textbf{tanzten} bis Mitternacht. Das \textbf{Publikum} war begeistert und klatschte laut.}
\par
Traduction : « Hier, il y a eu un grand concert à Yaoundé. Une chanteuse célèbre a chanté des chansons traditionnelles. Les musiciens jouaient des tam-tams, et beaucoup de gens ont dansé jusqu'à minuit. Le public était enthousiaste et applaudissait fort. »
\par
\textbf{Points de vigilance :} \all{auf den Trommeln} : datif pluriel avec \all{-n} ; \all{tanzten} : Präteritum de \all{tanzen} (verbe faible) ; \all{das Publikum} est singulier en allemand, d'où \all{war} et non \all{waren}.
\end{corrige}

\begin{corrige}[Exercice 9 — Compréhension de texte type Bac]
\textbf{A. Questions de compréhension} (réponses modèles, en allemand) :
\begin{enumerate}
\item \all{Der Erzähler besuchte das Konzert am Samstagabend mit seinen Freunden.} — « Le narrateur a assisté au concert samedi soir avec ses amis. »
\item \all{Tausende Menschen waren gekommen: junge und alte Leute, Deutsche und Ausländer.} — « Des milliers de personnes étaient venues : des jeunes et des vieux, des Allemands et des étrangers. »
\item \all{Das ganze Publikum tanzte, als ob es keinen Morgen gäbe.} — « Tout le public dansait comme s'il n'y avait pas de lendemain. »
\item \all{Das Publikum rief „Noch ein Lied!“, weil niemand nach Hause gehen wollte und alle begeistert waren.} — « Le public a crié "Encore une chanson !", parce que personne ne voulait rentrer à la maison. »
\item \all{Wenn er mehr Geld hätte, würde er jedes Wochenende zu solchen Konzerten gehen.} — « S'il avait plus d'argent, il irait à de tels concerts chaque week-end. »
\end{enumerate}
\textbf{B. Richtig oder Falsch :}
\begin{enumerate}
\item \all{Falsch}. — \all{„...ein großes Konzert in Berlin... auf einer riesigen Bühne im Zentrum der Stadt.“} Le concert avait lieu dans une grande ville, pas dans un petit village.
\item \all{Falsch}. — \all{„Die Musiker spielten auf modernen Instrumenten, aber auch auf traditionellen Trommeln.“} Ils utilisaient aussi des instruments traditionnels.
\item \all{Falsch}. — \all{„...aber niemand wollte nach Hause gehen.“} Personne ne voulait rentrer.
\end{enumerate}
\textbf{C. Questions de langue :}
\begin{enumerate}
\item Comparaison avec \all{als ob} : \all{„...tanzte das ganze Publikum, als ob es keinen Morgen gäbe.“} — comparaison irréelle, exagérée, avec Konjunktiv II (\all{gäbe}). Comparaison avec \all{wie} : \all{„Die alte Frau neben mir bewegte sich wie eine junge Tänzerin.“} — comparaison réelle, simple, sans verbe conjugué. Différence : \all{wie} compare deux réalités ; \all{als ob} exprime une hypothèse imaginaire et exige le Konjunktiv II.
\item Exemples : \all{„die berühmte Sängerin“} : nominatif féminin singulier, terminaison \all{-e} après article défini ; \all{„ein unvergesslicher Abend“} : nominatif masculin singulier, terminaison \all{-er} après article indéfini (déclinaison mixte). Autres possibilités : \all{„modernen Instrumenten“} (datif pluriel, \all{-en}), \all{„ein letztes traditionelles Lied“} (accusatif neutre, \all{-es}).
\item \all{Die Sängerin \textbf{würde lächeln} und \textbf{würde} ein letztes Lied \textbf{singen}.} (ou : \all{Die Sängerin lächelte und sänge ein letztes Lied}, forme littéraire). La forme avec \all{würde} + infinitif est la plus sûre.
\end{enumerate}
\textbf{Barème indicatif (20 pts) :} A : 5 $\times$ 2 pts = 10 pts (phrase complète et correcte : 2 pts ; erreur de grammaire : $-0,5$) ; B : 3 $\times$ 2 pts = 6 pts (1 pt pour la réponse, 1 pt pour la citation) ; C : 4 pts (1 + 2 + 1).
\par
\textbf{Points de vigilance :} répondez toujours par des phrases complètes en reprenant les mots de la question ; citez le texte entre guillemets allemands pour justifier ; vérifiez la place du verbe dans chaque phrase.
\end{corrige}

\begin{corrige}[Exercice 10 — Thème et version]
\textbf{A. Version (traduction française) :}
\par
« La fête de la bière (Oktoberfest) à Munich est la plus grande fête populaire du monde. Chaque année, des millions de visiteurs viennent de tous les pays. Les gens joyeux portent des vêtements bavarois traditionnels et dansent au son d'une musique forte. Certains visiteurs dansent sur les bancs, comme s'ils étaient des danseurs professionnels. Un célèbre écrivain allemand a dit un jour que la fête était comme un grand rêve. Si j'y étais, je danserais toute la nuit. »
\par
\textbf{Points de vigilance :} \all{das Volksfest} = la fête populaire ; \all{die bayerische Kleidung} = les vêtements bavarois ; \all{das Fest sei wie ein großer Traum} : discours indirect au Konjunktiv I, à rendre par « que la fête était... » ; \all{wenn ich dort wäre} : « si j'y étais ».
\par
\textbf{B. Thème (traduction allemande) :}
\begin{enumerate}
\item \all{Gestern gab es in Douala ein großes Fest mit traditioneller Musik.} — attention : \all{es gab} (Präteritum de \all{es gibt}) ; \all{traditioneller Musik} : datif féminin sans article après \all{mit}.
\item \all{Die jungen Mädchen tanzten wie professionelle Tänzerinnen.} — comparaison réelle avec \all{wie}, sans Konjunktiv II ; \all{die Tänzerin, -nen} : la danseuse.
\item \all{Der alte Mann spielte auf der Trommel, als ob er jung wäre.} — irréel : \all{als ob} + Konjunktiv II (\all{wäre}), verbe en fin de proposition ; \all{auf der Trommel spielen} : jouer du tam-tam.
\item \all{Wenn ich Zeit hätte, würde ich alle Feste meiner Region besuchen.} — \all{hätte} en fin de subordonnée ; \all{würde} en tête de principale ; \all{meiner Region} : génitif féminin.
\item \all{Das Publikum sang und tanzte bis Mitternacht.} — \all{das Publikum} est singulier : \all{sang}, \all{tanzte} ; \all{bis Mitternacht} : jusqu'à minuit.
\end{enumerate}
\textbf{Barème indicatif (20 pts) :} Version : 8 pts (fidélité du sens 5 pts, qualité du français 3 pts) ; Thème : 5 $\times$ 2,4 pts = 12 pts (vocabulaire 0,8 pt, grammaire : comparaison / Konjunktiv II / place du verbe 1 pt, orthographe et majuscules 0,6 pt).
\end{corrige}

\begin{corrige}[Exercice 11 — Production écrite type Bac]
\textbf{Proposition de rédaction modèle (environ 95 mots) :}

\begin{textebox}[Das Ngondo-Fest]
Letztes Jahr besuchte ich das Ngondo-Fest in Douala. Es war ein großes Fest am Wasser. Viele Menschen kamen: junge und alte Leute, Kinder und Eltern. Die Musiker spielten auf traditionellen Trommeln, und ein berühmter Sänger sang alte Lieder. Die Tänzer bewegten sich wie Vögel im Wind. Ein alter Mann tanzte, als ob er wieder zwanzig Jahre alt wäre. Das begeisterte Publikum klatschte und sang laut mit. Ich war sehr glücklich und stolz auf meine Kultur. Wenn ich mehr Zeit hätte, würde ich jedes Jahr zu diesem wunderbaren Fest gehen. Solche Feste sind wichtig für unsere Tradition.
\end{textebox}

\textbf{Analyse du texte modèle :}
\begin{itemize}
\item Vocabulaire du chapitre : \all{das Fest, die Trommeln, der Sänger, die Tänzer, das Publikum} (5 mots minimum : exigence remplie).
\item Comparaison réelle : \all{wie Vögel im Wind} ; comparaison irréelle : \all{als ob er wieder zwanzig Jahre alt wäre} (Konjunktiv II de \all{sein}).
\item Adjectifs déclinés : \all{traditionellen Trommeln} (datif pluriel), \all{ein berühmter Sänger} (nominatif masculin), \all{das begeisterte Publikum} (nominatif neutre), \all{alte Lieder} (accusatif pluriel sans article).
\item Konjunktiv II : \all{Wenn ich mehr Zeit hätte, würde ich ... gehen.}
\end{itemize}
\textbf{Barème indicatif (20 pts) :}
\begin{center}
\begin{tabular}{|l|l|}
\hline
\rowcolor{popblueL} Critère & Points \\
\hline
Vocabulaire thématique (au moins 5 mots du chapitre) & 4 pts \\
\hline
Grammaire : \all{wie} / \all{als ob} + Konjunktiv II corrects & 6 pts \\
\hline
Déclinaison de l'adjectif, orthographe, majuscules & 4 pts \\
\hline
Cohérence, structure, respect des 80--100 mots & 6 pts \\
\hline
\end{tabular}
\end{center}
\textbf{Points de vigilance :}
\begin{itemize}
\item Majuscule à tous les noms : \all{das Fest, die Musik, die Trommel}.
\item Après \all{als ob}, verbe conjugué en fin de proposition et au Konjunktiv II : \all{als ob er jung wäre}, et non \all{als ob er ist jung}.
\item Dans la principale après une subordonnée en \all{wenn}, le verbe est en première position : \all{Wenn ich Zeit hätte, würde ich...}.
\item Ne traduisez pas mot à mot le français : « il y avait » se dit \all{es gab}, « jusqu'à minuit » \all{bis Mitternacht}.
\item Comptez vos mots et relisez : terminaisons des adjectifs, place du verbe, Umlaute.
\end{itemize}
\end{corrige}

\newpage

\coursec{Fiche bilan}

\begin{popfigure}
\begin{center}\tcbox[colback=popblue,colframe=popblue,coltext=white,arc=9pt,boxsep=4pt,left=6pt,right=6pt]{\bfseries\small Musik und Tanz (AL23)}\end{center}
\vspace{-2pt}
\begin{tcbraster}[raster columns=3,raster equal height=rows,raster column skip=6pt,raster row skip=6pt,enhanced,blankest]
\begin{tcolorbox}[enhanced,colback=white,colframe=poporange,boxrule=0.9pt,arc=3pt,title={Texte support},fonttitle=\bfseries\scriptsize,coltitle=white,colbacktitle=poporange,left=4pt,right=4pt,top=2pt,bottom=2pt,boxsep=1pt,toptitle=1.5pt,bottomtitle=1.5pt,halign=left]
\begin{itemize}[nosep,leftmargin=*,itemsep=1pt,font=\scriptsize]
\item Musik in Kamerun
\item Musik in Deutschland
\end{itemize}
\end{tcolorbox}
\begin{tcolorbox}[enhanced,colback=white,colframe=popgreen,boxrule=0.9pt,arc=3pt,title={Lexique},fonttitle=\bfseries\scriptsize,coltitle=white,colbacktitle=popgreen,left=4pt,right=4pt,top=2pt,bottom=2pt,boxsep=1pt,toptitle=1.5pt,bottomtitle=1.5pt,halign=left]
\begin{itemize}[nosep,leftmargin=*,itemsep=1pt,font=\scriptsize]
\item die Musik: der Sänger, das Konzert
\item der Tanz: die Trommel, das Fest
\end{itemize}
\end{tcolorbox}
\begin{tcolorbox}[enhanced,colback=white,colframe=popblue!70,boxrule=0.9pt,arc=3pt,title={Comparaison},fonttitle=\bfseries\scriptsize,coltitle=white,colbacktitle=popblue!70,left=4pt,right=4pt,top=2pt,bottom=2pt,boxsep=1pt,toptitle=1.5pt,bottomtitle=1.5pt,halign=left]
\begin{itemize}[nosep,leftmargin=*,itemsep=1pt,font=\scriptsize]
\item wie (réelle)
\item als ob + KII (irréelle)
\end{itemize}
\end{tcolorbox}
\begin{tcolorbox}[enhanced,colback=white,colframe=poppurple,boxrule=0.9pt,arc=3pt,title={Déclinaison de l'adjectif},fonttitle=\bfseries\scriptsize,coltitle=white,colbacktitle=poppurple,left=4pt,right=4pt,top=2pt,bottom=2pt,boxsep=1pt,toptitle=1.5pt,bottomtitle=1.5pt,halign=left]
\begin{itemize}[nosep,leftmargin=*,itemsep=1pt,font=\scriptsize]
\item der / ein / Ø article
\end{itemize}
\end{tcolorbox}
\begin{tcolorbox}[enhanced,colback=white,colframe=poppink,boxrule=0.9pt,arc=3pt,title={Konjunktiv II},fonttitle=\bfseries\scriptsize,coltitle=white,colbacktitle=poppink,left=4pt,right=4pt,top=2pt,bottom=2pt,boxsep=1pt,toptitle=1.5pt,bottomtitle=1.5pt,halign=left]
\begin{itemize}[nosep,leftmargin=*,itemsep=1pt,font=\scriptsize]
\item wäre, hätte, würde + Inf.
\end{itemize}
\end{tcolorbox}
\begin{tcolorbox}[enhanced,colback=white,colframe=popgold,boxrule=0.9pt,arc=3pt,title={Expression},fonttitle=\bfseries\scriptsize,coltitle=white,colbacktitle=popgold,left=4pt,right=4pt,top=2pt,bottom=2pt,boxsep=1pt,toptitle=1.5pt,bottomtitle=1.5pt,halign=left]
\begin{itemize}[nosep,leftmargin=*,itemsep=1pt,font=\scriptsize]
\item résumé, commentaire
\end{itemize}
\end{tcolorbox}
\begin{tcolorbox}[enhanced,colback=white,colframe=popdark,boxrule=0.9pt,arc=3pt,title={Méthode Bac},fonttitle=\bfseries\scriptsize,coltitle=white,colbacktitle=popdark,left=4pt,right=4pt,top=2pt,bottom=2pt,boxsep=1pt,toptitle=1.5pt,bottomtitle=1.5pt,halign=left]
\begin{itemize}[nosep,leftmargin=*,itemsep=1pt,font=\scriptsize]
\item lire -- repérer -- rédiger
\end{itemize}
\end{tcolorbox}
\end{tcbraster}

\legende{Carte mentale de la leçon AL23 : la vie culturelle (musique et danse)}
\end{popfigure}

\begin{bilan}
\textbf{1. Lexique clé : die Musik und der Tanz}
\begin{itemize}
  \item \all{die Musik, -en} : la musique ; \all{der Musiker, - / die Musikerin, -nen} : le musicien / la musicienne
  \item \all{das Konzert, -e} : le concert ; \all{das Lied, -er} : la chanson ; \all{die Melodie, -n} : la mélodie
  \item \all{der Sänger, - / die Sängerin, -nen} : le chanteur / la chanteuse ; \all{singen (sang, hat gesungen)} : chanter
  \item \all{der Tanz, \"-e} : la danse ; \all{tanzen} : danser ; \all{der Tänzer, -} : le danseur
  \item \all{die Trommel, -n} : le tambour, la tam-tam ; \all{das Instrument, -e} : l'instrument ; \all{die Gitarre, -n} : la guitare
  \item \all{das Fest, -e} : la fête ; \all{feiern} : fêter, célébrer ; \all{der Rhythmus, die Rhythmen} : le rythme
  \item \all{auftreten (trat auf, ist aufgetreten)} : se produire (sur scène) ; \all{das Publikum} : le public
\end{itemize}

\textbf{2. Grammaire à connaître par c\oe ur}
\begin{center}
\begin{tabularx}{\linewidth}{|l|X|X|}
\hline
\rowcolor{popblueL} \textbf{Point} & \textbf{Règle} & \textbf{Exemple} \\
\hline
Comparaison réelle & \all{wie} + adjectif (comparatif possible : \all{so ... wie}) & \all{Er singt so gut wie sein Vater.} \\
\hline
Comparaison irréelle & \all{als ob} + verbe conjugué au \textbf{Konjunktiv II} en fin de phrase & \all{Sie tanzt, als ob sie schwebte.} \\
\hline
Déclinaison après \all{der/das/die} & adjectif en \textbf{-e} ou \textbf{-en} (faible) & \all{das große Fest, die jungen Sänger} \\
\hline
Déclinaison après \all{ein/kein} & adjectif mixte : \all{ein guter Sänger}, \all{kein gutes Lied} & \all{ein lautes Konzert} \\
\hline
Sans article & adjectif fort (il porte la marque du cas) & \all{laute Musik, junger Wein} \\
\hline
Konjunktiv II & \all{wäre} (sein), \all{hätte} (haben), \all{würde} + infinitif ; irréel, souhait, hypothèse & \all{Ich hätte gern mehr Zeit für Musik.} \\
\hline
\end{tabularx}
\end{center}

\textbf{3. Méthodes express}
\begin{itemize}
  \item \textbf{Compréhension de texte} : lire une première fois pour le sens global, repérer le thème (musique/danse), puis répondre aux questions en citant le texte, sans recopier de longs passages.
  \item \textbf{Phrase avec \all{als ob}} : sujet + verbe + \all{als ob} + ... + verbe au KII \textbf{à la fin}. Ne jamais conjuguer au présent après \all{als ob}.
  \item \textbf{Choisir la déclinaison} : 1) y a-t-il un article ? 2) lequel (défini, indéfini, aucun) ? 3) appliquer la terminaison correspondante.
  \item \textbf{Production guidée} : 5 à 8 phrases, présent ou Perfekt, 3 mots du lexique au minimum, verbe en 2\up{e} position.
\end{itemize}
\end{bilan}

\begin{aretenir}[Checklist avant le Bac]
\begin{itemize}
  \item $\square$ Je connais l'article et le pluriel des noms du thème : \all{das Konzert, -e}, \all{die Trommel, -n}, \all{das Fest, -e}, \all{der Tanz, \"-e}.
  \item $\square$ Je sais conjuguer \all{singen}, \all{tanzen}, \all{feiern}, \all{auftreten} au Präsens et au Perfekt.
  \item $\square$ Je distingue comparaison réelle (\all{wie}) et comparaison irréelle (\all{als ob} + Konjunktiv II).
  \item $\square$ Je place le verbe conjugué en dernière position après \all{als ob}.
  \item $\square$ Je connais par c\oe ur \all{wäre}, \all{hätte}, \all{würde} + infinitif.
  \item $\square$ Je sais décliner l'adjectif après article défini (\all{der junge Sänger}), indéfini (\all{ein junger Sänger}) et sans article (\all{junger Sänger}).
  \item $\square$ Je n'oublie pas la majuscule à tous les noms allemands (\all{die Musik}, \all{das Lied}).
  \item $\square$ Je mets le verbe conjugué en deuxième position dans la phrase déclarative.
  \item $\square$ Je réponds aux questions de compréhension par des phrases complètes, en m'appuyant sur le texte.
  \item $\square$ Je sais rédiger un court paragraphe (5 à 8 phrases) sur une fête ou un concert avec le vocabulaire de la leçon.
  \item $\square$ Je relis ma production : accords, cas, place du verbe, orthographe des Umlaute.
\end{itemize}
\end{aretenir}

\findecours

\end{document}
